% Leçon 14 — Expression écrite : le paludisme — Chinois Tle A — Cours signé OnBuch+
% Programme officiel MINESEC. Compiler avec : tectonic ZH14-ecrit-le-paludisme.tex
\newcommand{\DOCMATIERE}{Chinois}
\newcommand{\DOCNIVEAU}{Tle A}
\newcommand{\DOCMODULE}{Module 2 : Environnement et santé}
\newcommand{\DOCLECON}{ZH14}
\newcommand{\DOCTITRE}{Leçon 14 — Expression écrite : le paludisme}
% OnBuch+ — préambule « cours signés » ENRICHI (compilé avec tectonic / XeLaTeX)
% Système visuel « Pop » d'OnBuch+ : fond crème (jamais blanc), bordures encre
% épaisses, ombres « dures » décalées, titres Archivo Black en capitales.
% Version MATHÉMATIQUES : graphes (pgfplots/tikz), boîtes pédagogiques
% (définition, propriété, méthode, attention, le savais-tu, exercice résolu,
% exercice, TP, bilan), page de garde et signature OnBuch+ ancrée en bas.
%
% Macros attendues dans main.tex AVANT \input{preamble} :
%   \DOCMATIERE  \DOCNIVEAU  \DOCMODULE  \DOCLECON (numéro)  \DOCTITRE
\documentclass[11pt]{article}

\usepackage{fontspec}
\usepackage[french]{babel} % français = langue principale ; le chinois via \zh{..} / textebox (xeCJK)
\usepackage{xeCJK}
\usepackage{pifont} % \ding{51} \ding{55} utilisés par les modèles
\usepackage[a4paper,margin=2cm,top=3.4cm,bottom=2.8cm,headheight=1.4cm,footskip=1.3cm]{geometry}
\usepackage{amsmath,amssymb,amsfonts}
\usepackage{mathtools} % \xmapsto, \coloneqq... souvent utilisés par les modèles
\usepackage{cancel} % simplifications barrées (bilans de forces)
% \abs{z} (module, valeur absolue) et \norm{u} (norme), utilisés par les modèles
\providecommand{\abs}{}\let\abs\relax\DeclarePairedDelimiter{\abs}{\lvert}{\rvert}
\providecommand{\norm}{}\let\norm\relax\DeclarePairedDelimiter{\norm}{\lVert}{\rVert}
\usepackage[table]{xcolor} % [table] : fournit aussi \rowcolors (alternance de lignes)
\usepackage{pagecolor}
\usepackage{tikz}
\usetikzlibrary{arrows.meta,positioning,calc,shapes.geometric,shapes.misc,shapes.arrows,decorations.pathmorphing,decorations.pathreplacing,decorations.markings,patterns,shadows,mindmap,trees,fit,backgrounds,matrix,intersections,angles,quotes}
\usepackage{pgfplots}
\usepackage[version=4]{mhchem} % équations nucléaires \ce{^{235}_{92}U}
\usepackage[european,siunitx]{circuitikz} % schémas électriques
\pgfplotsset{compat=1.18}
\usepgfplotslibrary{fillbetween,groupplots} % aires entre courbes (calcul intégral)
\usepackage{enumitem}
\usepackage{fancyhdr}
% [most] charge toutes les bibliothèques tcolorbox (listings, minted, raster,
% documentation...) et gonfle inutilement la mémoire TeX sur un document long
% (~300 boîtes) : on ne charge que ce qui est réellement utilisé.
\usepackage[skins,breakable]{tcolorbox}
\usepackage{graphicx}
\usepackage{colortbl}
\usepackage{booktabs}
\usepackage{tabularx}
\usepackage{diagbox} % case d'en-tête diagonale des tableaux à double entrée
% Colonnes extensibles centrée / gauche / droite (C, L, R), souvent utilisées par les modèles
\newcolumntype{C}{>{\centering\arraybackslash}X}
\newcolumntype{L}{>{\raggedright\arraybackslash}X}
\newcolumntype{R}{>{\raggedleft\arraybackslash}X}
\usepackage{array}
\usepackage{multicol}
\usepackage{multirow}
\usepackage{siunitx}
% parse-numbers=false : accepte aussi \SI{2,5 \times 10^{-3}}{...}
\sisetup{locale=FR,output-decimal-marker={,},group-minimum-digits=4,parse-numbers=false}
\DeclareSIUnit\ohm{\ensuremath{\symup{\Omega}}} % U+2126 absent de la police
\DeclareSIUnit\division{div} % réglages d'oscilloscope (ms/div, V/div)
\DeclareSIUnit\tr{tr} % tours (vitesse de rotation, ex. \SI{1500}{\tr\per\minute})
\DeclareSIUnit\molar{M} % concentration molaire (ex. \SI{3}{\molar}, \SI{1}{\micro\molar})
\DeclareSIUnit\an{an} % année (le modèle confond parfois avec \year, un compteur TeX) — ex. \SI{5}{\kilo\meter\per\an}
\DeclareSIUnit\FCFA{FCFA} % franc CFA (ex. \SI{500}{\FCFA\per\kilo\gram})
\DeclareSIUnit\degreeBrix{\textdegree Brix} % teneur en sucre d'un sirop/jus (réfractométrie)
\DeclareSIUnit\month{mois} % le modèle utilise parfois \month, un compteur TeX, comme unité de durée
\usepackage{qrcode}
% \degree : commande fréquemment utilisée par les modèles pour un angle ou une
% température en dehors de \SI{}{} (non fournie par les paquets chargés).
\providecommand{\degree}{^{\circ}}
% \qed (fin de démonstration), utilisé par les modèles sans amsthm
\providecommand{\qed}{\hfill$\square$}
% \barre : double barre des valeurs interdites dans un tableau de variations (tkz-tab)
\providecommand{\barre}{\ensuremath{\|}}
% environnement proof (amsthm non chargé) utilisé par les modèles
\ifdefined\proof\else\newenvironment{proof}[1][Démonstration]{\par\noindent\textit{#1.}\ }{\qed\par}\fi
\usepackage{lastpage}
\usepackage{hyperref}
\hypersetup{hidelinks}

% ============================================================
% Palette « Pop » OnBuch+ — mêmes hex que la classe P (plus_ui.dart)
% ============================================================
\definecolor{poporange}{HTML}{F59321}
\definecolor{popdark}{HTML}{E07A0C}
\definecolor{popcream}{HTML}{FFF6E8}
\definecolor{popink}{HTML}{1C1714}
\definecolor{poppurple}{HTML}{7A5AE0}
\definecolor{popgreen}{HTML}{2E9E6A}
\definecolor{poppink}{HTML}{E0457B}
\definecolor{popblue}{HTML}{2F7DE1}
\definecolor{popgold}{HTML}{C98A12}
\definecolor{popmuted}{HTML}{8A7F76}
\definecolor{popline}{HTML}{EADFD2}
\definecolor{popgray}{HTML}{8A7F76}
\colorlet{popgrey}{popgray}
% Alias tolérants (le modèle nomme parfois les couleurs différemment)
\colorlet{popwhite}{white}
\colorlet{popblack}{popink}
\colorlet{popred}{poppink}
\colorlet{popyellow}{popgold}
% Teintes claires (fonds de tableaux/encadrés) — le modèle nomme parfois
% les couleurs avec un suffixe L (ex. popblueL) sans les avoir définies.
\colorlet{poporangeL}{poporange!20}
\colorlet{popdarkL}{popdark!20}
\colorlet{poppurpleL}{poppurple!20}
\colorlet{popgreenL}{popgreen!20}
\colorlet{poppinkL}{poppink!20}
\colorlet{popblueL}{popblue!20}
\colorlet{popgoldL}{popgold!20}
\colorlet{popredL}{popred!20}
\colorlet{popgrayL}{popgray!20}
% Teintes claires pour fonds de tableaux / zones
\colorlet{popblueL}{popblue!10!white}
\colorlet{poporangeL}{poporange!14!white}
\colorlet{popgreenL}{popgreen!12!white}
\colorlet{poppurpleL}{poppurple!12!white}
\colorlet{poppinkL}{poppink!12!white}
\colorlet{popgoldL}{popgold!14!white}

\pagecolor{popcream}

% ============================================================
% Typographie
% ============================================================
\defaultfontfeatures{Ligatures=TeX}
\setmainfont{PlusJakartaSans-Medium.ttf}[Path=../../../fonts/,BoldFont=PlusJakartaSans-Bold.ttf]
% Symboles, API et emojis absents de la police principale : repli sur DejaVu Sans (caractères actifs)
\newfontface\symfont{DejaVuSans.ttf}[Path=../../../fonts/]
\makeatletter
\newcommand\symactive[1]{\catcode#1=13 \begingroup\lccode`\~=#1 \lowercase{\endgroup\def~}{{\symfont\char#1}}}
\newcommand\symblank[1]{\catcode#1=13 \begingroup\lccode`\~=#1 \lowercase{\endgroup\def~}{}}
\newcommand\symhyph[1]{\catcode#1=13 \begingroup\lccode`\~=#1 \lowercase{\endgroup\def~}{\mbox{-}}}
\newcount\symc
\newcommand\symrange[2]{\symc=#1 \@whilenum\symc<#2 \do{\expandafter\symactive\expandafter{\the\symc}\advance\symc by 1 }}
\newcommand\symblankrange[2]{\symc=#1 \@whilenum\symc<#2 \do{\expandafter\symblank\expandafter{\the\symc}\advance\symc by 1 }}
\makeatother
\symrange{"0250}{"02FF}\symactive{"2713}\symactive{"2714}\symactive{"2717}\symactive{"2718}\symactive{"2610}\symactive{"2611}\symactive{"2612}
\symrange{"2600}{"27BF}
\catcode"2705=13 \begingroup\lccode`\~="2705 \lowercase{\endgroup\def~}{{\symfont\char"2713}}\symblank{"26BD}\symblank{"26A1}
\symrange{"2460}{"257F}\symactive{"223C}\symactive{"2248}\symactive{"2260}
\symhyph{"2011}\symhyph{"2E1A}\symblankrange{"1F300}{"1FAFF}\symblank{"FE0F}
% Caractères chinois : WenQuanYi Zen Hei (sous-ensemble GB2312 + pinyin), gras/italique simulés
\setCJKmainfont{WenQuanYiZenHei-subset.ttf}[Path=../../../fonts/,AutoFakeBold=3,AutoFakeSlant=0.2]
\xeCJKsetup{CJKspace=false}
% Pinyin : voyelles à caron (ǎ ǐ ǒ ǔ ǖ ǘ ǚ ǜ) absentes de la police principale -> caractères actifs
\newfontface\pyfont{WenQuanYiZenHei-subset.ttf}[Path=../../../fonts/]
\newcommand\pyactive[1]{\catcode#1=13 \begingroup\lccode`\~=#1 \lowercase{\endgroup\def~}{{\pyfont\char#1}}}
\pyactive{"01CD}\pyactive{"01CE}\pyactive{"01CF}\pyactive{"01D0}\pyactive{"01D1}\pyactive{"01D2}\pyactive{"01D3}\pyactive{"01D4}\pyactive{"01D5}\pyactive{"01D6}\pyactive{"01D7}\pyactive{"01D8}\pyactive{"01D9}\pyactive{"01DA}\pyactive{"01DB}\pyactive{"01DC}
% Maths en sans-serif (Fira Math) avec chiffres et lettres droites de Plus
% Jakarta, pour que les formules se fondent dans le texte.
\usepackage[math-style=ISO,bold-style=ISO]{unicode-math}
\setmathfont{FiraMath-Regular.otf}[Path=../../../fonts/]
\setmathfont{PlusJakartaSans-Medium.ttf}[Path=../../../fonts/,range={up/num,up/{Latin,latin},\mathup}]
\usepackage{icomma} % virgule décimale sans espace en mode math
% Caractères mathématiques Unicode absents des polices de texte (Plus Jakarta,
% Archivo Black) : rendus par leur équivalent mathématique, en texte comme en
% formule — sinon ils s'affichent en carré vide (ex. « Arithmétique dans ℤ »).
\usepackage{newunicodechar}
\newunicodechar{ℝ}{\ensuremath{\mathbb{R}}}
\newunicodechar{ℤ}{\ensuremath{\mathbb{Z}}}
\newunicodechar{ℂ}{\ensuremath{\mathbb{C}}}
\newunicodechar{ℕ}{\ensuremath{\mathbb{N}}}
\newunicodechar{ℚ}{\ensuremath{\mathbb{Q}}}
\newunicodechar{𝔻}{\ensuremath{\mathbb{D}}}
\newunicodechar{α}{\ensuremath{\alpha}}
\newunicodechar{β}{\ensuremath{\beta}}
\newunicodechar{γ}{\ensuremath{\gamma}}
\newunicodechar{δ}{\ensuremath{\delta}}
\newunicodechar{ε}{\ensuremath{\varepsilon}}
\newunicodechar{θ}{\ensuremath{\theta}}
\newunicodechar{λ}{\ensuremath{\lambda}}
\newunicodechar{μ}{\ensuremath{\mu}}
\newunicodechar{σ}{\ensuremath{\sigma}}
\newunicodechar{τ}{\ensuremath{\tau}}
\newunicodechar{φ}{\ensuremath{\varphi}}
\newunicodechar{ψ}{\ensuremath{\psi}}
\newunicodechar{ω}{\ensuremath{\omega}}
\newunicodechar{Σ}{\ensuremath{\Sigma}}
% majuscules produites par \MakeUppercase dans les titres (α-aminés -> Α)
\newunicodechar{Α}{\ensuremath{\alpha}}
\newunicodechar{Β}{\ensuremath{\beta}}
\newunicodechar{Γ}{\ensuremath{\Gamma}}
\newunicodechar{Δ}{\ensuremath{\Delta}}
\newunicodechar{Ε}{\ensuremath{\varepsilon}}
\newunicodechar{Θ}{\ensuremath{\Theta}}
\newunicodechar{Λ}{\ensuremath{\Lambda}}
\newunicodechar{Μ}{\ensuremath{\mu}}
\newunicodechar{Τ}{\ensuremath{\tau}}
\newunicodechar{Φ}{\ensuremath{\Phi}}
\newunicodechar{Ψ}{\ensuremath{\Psi}}
\newunicodechar{Ω}{\ensuremath{\Omega}}
\newunicodechar{↦}{\ensuremath{\mapsto}}
\newunicodechar{∈}{\ensuremath{\in}}
\newunicodechar{∉}{\ensuremath{\notin}}
\newunicodechar{∧}{\ensuremath{\wedge}}
\newunicodechar{⇔}{\ensuremath{\Leftrightarrow}}
\newunicodechar{⟺}{\ensuremath{\Longleftrightarrow}}
\newunicodechar{⇒}{\ensuremath{\Rightarrow}}
\newunicodechar{⟹}{\ensuremath{\Longrightarrow}}
\newunicodechar{∘}{\ensuremath{\circ}}
\newunicodechar{∪}{\ensuremath{\cup}}
\newunicodechar{∩}{\ensuremath{\cap}}
\newunicodechar{≡}{\ensuremath{\equiv}}
\newunicodechar{⊂}{\ensuremath{\subset}}
\newunicodechar{⊥}{\ensuremath{\perp}}
\newunicodechar{∃}{\ensuremath{\exists}}
\newunicodechar{∀}{\ensuremath{\forall}}
\newunicodechar{∛}{\ensuremath{\sqrt[3]{\,}}}
\newunicodechar{∜}{\ensuremath{\sqrt[4]{\,}}}
\newunicodechar{⃗}{\ensuremath{{}^{\to}}}
\newunicodechar{ⁿ}{\ifmmode^{n}\else\textsuperscript{\ensuremath{n}}\fi}
\newunicodechar{ˣ}{\ifmmode^{x}\else\textsuperscript{\ensuremath{x}}\fi}
\newunicodechar{ᵇ}{\ifmmode^{b}\else\textsuperscript{\ensuremath{b}}\fi}
\newunicodechar{ʳ}{\ifmmode^{r}\else\textsuperscript{\ensuremath{r}}\fi}
\newunicodechar{ᵘ}{\ifmmode^{u}\else\textsuperscript{\ensuremath{u}}\fi}
\newunicodechar{ʸ}{\ifmmode^{y}\else\textsuperscript{\ensuremath{y}}\fi}
\newunicodechar{ᵃ}{\ifmmode^{a}\else\textsuperscript{\ensuremath{a}}\fi}
\newunicodechar{ᵉ}{\ifmmode^{e}\else\textsuperscript{\ensuremath{e}}\fi}
\newunicodechar{ᵏ}{\ifmmode^{k}\else\textsuperscript{\ensuremath{k}}\fi}
\newunicodechar{ᵖ}{\ifmmode^{p}\else\textsuperscript{\ensuremath{p}}\fi}
\newunicodechar{ᴺ}{\ifmmode^{N}\else\textsuperscript{\ensuremath{N}}\fi}
\newunicodechar{ᶜ}{\ifmmode^{c}\else\textsuperscript{\ensuremath{c}}\fi}
\newunicodechar{ʰ}{\ifmmode^{h}\else\textsuperscript{\ensuremath{h}}\fi}
\newunicodechar{ᵠ}{\ifmmode^{q}\else\textsuperscript{\ensuremath{q}}\fi}
\newunicodechar{ₙ}{\ifmmode_{n}\else\textsubscript{\ensuremath{n}}\fi}
\newunicodechar{ₐ}{\ifmmode_{a}\else\textsubscript{\ensuremath{a}}\fi}
\newunicodechar{ᵢ}{\ifmmode_{i}\else\textsubscript{\ensuremath{i}}\fi}
\newunicodechar{ᵣ}{\ifmmode_{r}\else\textsubscript{\ensuremath{r}}\fi}
\newunicodechar{ₖ}{\ifmmode_{k}\else\textsubscript{\ensuremath{k}}\fi}
\newunicodechar{ᵦ}{\ifmmode_{\beta}\else\textsubscript{\ensuremath{\beta}}\fi}
\newunicodechar{ₓ}{\ifmmode_{x}\else\textsubscript{\ensuremath{x}}\fi}
\newfontfamily\popdisplay{ArchivoBlack-Regular.ttf}[Path=../../../fonts/]
\newfontfamily\poplogo{SpaceGrotesk-Bold.ttf}[Path=../../../fonts/]
\newfontfamily\popsemi{PlusJakartaSans-SemiBold.ttf}[Path=../../../fonts/]
\newfontfamily\popxbold{PlusJakartaSans-ExtraBold.ttf}[Path=../../../fonts/]
\newfontfamily\popxboldls{PlusJakartaSans-ExtraBold.ttf}[Path=../../../fonts/,LetterSpace=3.0]

\setlist[enumerate]{leftmargin=1.7em,itemsep=5pt,topsep=4pt}
\setlist[itemize]{leftmargin=1.7em,itemsep=4pt,topsep=4pt,label={\color{poporange}\textbullet}}
\setlength{\parindent}{0pt}
\setlength{\parskip}{5pt}
\renewcommand{\arraystretch}{1.25}

% pgfplots : style Pop par défaut
\pgfplotsset{
  every axis/.append style={axis line style={popink,line width=1pt},tick style={popink},
    grid style={popline,line width=0.5pt},label style={font=\small},tick label style={font=\footnotesize}},
  popcurve/.style={poporange,line width=1.8pt,no markers,smooth},
}
% Même style disponible dans un \draw TikZ simple (hors pgfplots)
\tikzset{popcurve/.style={poporange,line width=1.8pt}}
\tikzset{popgrid/.style={popline,very thin}}
\colorlet{popcurve}{poporange} % le nom du style est parfois utilisé comme couleur

% ============================================================
% Pill / squiggle
% ============================================================
\newcommand{\poppill}[3]{%
  \tikz[baseline=-0.55ex]\node[draw=popink,line width=1.2pt,rounded corners=2.6pt,%
    fill=#2,inner xsep=6.5pt,inner ysep=3pt]{%
    \popxboldls\fontsize{7}{7}\selectfont\color{#3}\MakeUppercase{#1}};%
}
\newcommand{\popsquiggle}[1]{%
  \tikz[baseline=-0.4ex]{
    \draw[#1,line width=2.6pt,line cap=round]
      plot [smooth] coordinates {(0,0.09) (0.34,0.01) (0.68,0.21) (1.02,0.01) (1.36,0.10)};
  }%
}
% Mot-clé surligné (marqueur orange sous le texte)
\newcommand{\cle}[1]{{\popxbold\color{popdark}#1}}

% ============================================================
% En-tête / pied
% ============================================================
\pagestyle{fancy}
\fancyhf{}
\renewcommand{\headrulewidth}{0pt}
\renewcommand{\footrulewidth}{0pt}
\lhead{\raisebox{-0.15ex}{\poplogo\fontsize{13}{13}\selectfont\color{popink}OnBuch\color{poporange}+}}
\rhead{\poppill{\DOCMATIERE\ \textbullet\ \DOCNIVEAU\ \textbullet\ Leçon \DOCLECON}{white}{popink}}
\lfoot{\popxbold\fontsize{7.6}{7.6}\selectfont\color{popmuted}\MakeUppercase{Cours signé OnBuch+}\ \textbullet\ {\popsemi\DOCTITRE}}
\rfoot{\tikz[baseline=-0.6ex]\node[rounded corners=8.5pt,fill=popink,minimum height=17pt,minimum width=17pt,inner xsep=5pt,inner sep=0]{\popxbold\fontsize{8}{8}\selectfont\color{popcream}\thepage\,/\,\pageref*{LastPage}};}

% ============================================================
% Titres
% ============================================================
\newcommand{\chaptitre}[2]{%
  \begin{tcolorbox}[enhanced,colback=poporange,colframe=popink,boxrule=2.6pt,arc=11pt,
      drop shadow=popink,left=15pt,right=15pt,top=12pt,bottom=11pt,before skip=3pt,after skip=18pt]
    \poppill{#2}{popink}{popcream}\par\vspace{9pt}
    {\raggedright\hyphenpenalty=10000\exhyphenpenalty=10000\popdisplay\fontsize{22}{24}\selectfont\color{popcream}\MakeUppercase{#1}\par}
  \end{tcolorbox}
}
% Section numérotée : pastille orange avec le numéro + titre Archivo
\newcounter{popsec}
\newcommand{\coursec}[1]{%
  \stepcounter{popsec}\setcounter{popsub}{0}%
  \par\vspace{16pt}\noindent
  \begin{minipage}{\linewidth}
  \tikz[baseline=-0.7ex]\node[draw=popink,line width=1.6pt,fill=poporange,rounded corners=4pt,minimum size=22pt,inner sep=0]{\popdisplay\fontsize{12}{12}\selectfont\color{popcream}\thepopsec};%
  \hspace{8pt}{\popdisplay\fontsize{15}{17}\selectfont\color{popink}\MakeUppercase{#1}}\par
  \vspace{4pt}\noindent{\color{poporange}\rule{36pt}{3.6pt}}
  \end{minipage}\par\nopagebreak\vspace{8pt}
}
\newcounter{popsub}
\newcommand{\courssub}[1]{%
  \stepcounter{popsub}%
  \par\vspace{9pt}\noindent{\popxbold\fontsize{11.5}{13}\selectfont\color{popink}{\color{poporange}\thepopsec.\thepopsub}\hspace{6pt}#1}\par\nopagebreak\vspace{4pt}
}

% ============================================================
% Boîtes pédagogiques — même recette : fond blanc, bordure épaisse colorée,
% ombre dure encre, badge plein. Titre optionnel [#1].
% ============================================================
\tcbset{popbase/.style={breakable,enhanced,colback=white,boxrule=2.3pt,arc=9pt,
  shadow={3pt}{-3pt}{0pt}{popink},left=14pt,right=14pt,top=11pt,bottom=11pt,before skip=11pt,after skip=15pt,
  fontupper=\popsemi\fontsize{10.6}{14.5}\selectfont\color{popink},
  fontlower=\popsemi\fontsize{10.6}{14.5}\selectfont\color{popink}}}
% Coche dessinée (le glyphe ✓ n'existe pas dans les polices du document)
\newcommand{\popcheck}[1]{\tikz[baseline=-0.5ex]\draw[#1,line width=1.6pt,line cap=round,line join=round](-0.13,0.0)--(-0.04,-0.09)--(0.13,0.1);}
\providecommand{\coche}{\popcheck{popgreen}}
\providecommand{\resultat}[1]{\par\smallskip\noindent\fcolorbox{popgreen}{popgreenL}{\parbox{\dimexpr\linewidth-2\fboxsep-2\fboxrule\relax}{\textbf{Résultat.} #1}}\par\smallskip}
\newcommand{\popboxhead}[3]{\poppill{#1}{#2}{white}\ifx\relax#3\relax\else\hspace{6pt}{\popxbold\fontsize{10.6}{12}\selectfont\color{popink}#3}\fi\par\vspace{7pt}}

\newtcolorbox{aretenir}[1][]{popbase,colframe=popgreen,before upper={\popboxhead{À retenir}{popgreen}{#1}}}
% Texte en chinois (document de lecture, dialogue, extrait) : cadre
% d'étude. Un titre optionnel (source, auteur) s'affiche dans l'en-tête.
\newtcolorbox{textebox}[1][]{popbase,colframe=popblue,before upper={\popboxhead{文本 · Texte}{popblue}{#1}},after upper={}}
% Marques ✓ / ✗ (forme correcte / fautive) souvent utilisées par les modèles
\providecommand{\cross}{\textcolor{poppink}{\ensuremath{\times}}}
\providecommand{\xmark}{\cross}\providecommand{\crossmark}{\cross}\providecommand{\cmark}{\ensuremath{\checkmark}}
% Ligne de réponse / justification à compléter (inventée par les modèles)
\providecommand{\justification}[1]{\par\noindent\textit{Justification~:} \dotfill\par}
% \textipa (package tipa non chargé) : la police ne couvre pas l'API -> simple passage
\providecommand{\textipa}[1]{#1}
% Soulignement (ulem non chargé) : \uline, \ul
\providecommand{\uline}[1]{\underline{#1}}\providecommand{\ul}[1]{\underline{#1}}
% \glossaire{\item ...} inventée par les modèles : liste de glossaire
\providecommand{\glossaire}[1]{\begin{itemize}#1\end{itemize}}
% \alert (beamer) inventée par les modèles : texte mis en évidence
\providecommand{\alert}[1]{\textbf{\textcolor{poppink}{#1}}}
% variantes ASCII d'ordinaux écrites en macro
\providecommand{\iere}{\textsuperscript{re}}\providecommand{\ieme}{\textsuperscript{e}}\providecommand{\ere}{\textsuperscript{re}}\providecommand{\eme}{\textsuperscript{e}}\providecommand{\er}{\textsuperscript{er}}\providecommand{\ie}{\textsuperscript{e}}
% Ordinaux français écrits en macro par les modèles : 1\ière, 3\ième
\newcommand{\ière}{\textsuperscript{re}}\newcommand{\ième}{\textsuperscript{e}}
% \exemple (inventée par les modèles) : étiquette « Exemple : »
\providecommand{\exemple}{\textbf{Exemple~:}\ }
% \square utilisable aussi hors mode math (cases à cocher)
\AtBeginDocument{\let\squareMath\square\renewcommand{\square}{\ensuremath{\squareMath}}}
% Mot ou phrase en chinois à l'intérieur d'un paragraphe français
\newcommand{\zh}[1]{\ifmmode\text{#1}\else{#1}\fi}
\let\eng\zh \let\esp\zh \let\all\zh % alias tolérés
% flèches d'intonation inventées par les modèles
\providecommand{\textsearrow}{}\renewcommand{\textsearrow}{\ensuremath{\searrow}}\providecommand{\textnearrow}{}\renewcommand{\textnearrow}{\ensuremath{\nearrow}}
% \fr{..} : mot français (inventé par les modèles)
\providecommand{\fr}[1]{\textit{#1}}
% zhbox : encadré de texte chinois inventé par les modèles
\newenvironment{zhbox}{\begin{quote}}{\end{quote}}
% \vs : « versus » inventé par les modèles
\providecommand{\vs}{\textit{vs}\ }
% \blank : case à compléter inventée par les modèles
\providecommand{\blank}[1][]{\underline{\hspace{1.6em}}}
\newtcolorbox{exemplebox}[1][]{popbase,colframe=poporange,before upper={\popboxhead{Exemple}{poporange}{#1}}}
\newtcolorbox{definition}[1][]{popbase,colframe=popblue,before upper={\popboxhead{Définition}{popblue}{#1}}}
\newtcolorbox{propriete}[1][]{popbase,colframe=poppurple,before upper={\popboxhead{Propriété}{poppurple}{#1}}}
\newtcolorbox{methode}[1][]{popbase,colframe=popgold,before upper={\popboxhead{Méthode}{popgold}{#1}}}
\newtcolorbox{attention}[1][]{popbase,colframe=poppink,before upper={\popboxhead{Attention}{poppink}{#1}}}
\newtcolorbox{experience}[1][]{popbase,colframe=popblue,colback=popblueL,before upper={\popboxhead{Expérience}{popblue}{#1}}}
% « Le savais-tu ? » : carte violette pleine (contexte, culture, Cameroun)
\newtcolorbox{savaistu}[1][]{popbase,colback=poppurple,colframe=popink,
  fontupper=\popsemi\fontsize{10.4}{14.2}\selectfont\color{white},
  before upper={\poppill{Le savais-tu\,?}{white}{poppurple}\ifx\relax#1\relax\else\hspace{6pt}{\popxbold\color{white}#1}\fi\par\vspace{7pt}}}
% Exercice résolu : énoncé en haut, \tcblower puis la solution sur fond crème
\newcounter{popexo}\newcounter{popexr}
\newtcolorbox{exoresolu}[1][]{popbase,colframe=poporange,colbacklower=popcream,
  segmentation style={popink,line width=1.2pt,dashed},
  before upper={\stepcounter{popexr}\popboxhead{Exercice résolu \thepopexr}{poporange}{#1}},
  before lower={\poppill{Solution}{popink}{popcream}\par\vspace{6pt}}}
% Exercice (non résolu en place) — la correction est dans la partie Corrigés
\newtcolorbox{exercice}[1][]{popbase,colframe=popink,
  before upper={\stepcounter{popexo}\popboxhead{Exercice \thepopexo}{popink}{#1}}}
% Corrigé
\newtcolorbox{corrige}[1][]{popbase,colframe=popgreen,colback=popgreenL,
  before upper={\popboxhead{Corrigé}{popgreen}{#1}}}
% Objectifs / prérequis / bilan
\newtcolorbox{objectifs}{popbase,colframe=popink,colback=poporangeL,
  before upper={\popboxhead{Objectifs de la leçon}{popink}{}}}
\newtcolorbox{prerequis}{popbase,colframe=popink,colback=popblueL,
  before upper={\popboxhead{Prérequis}{popblue}{}}}
\newtcolorbox{bilan}{popbase,colframe=popink,colback=white,boxrule=2.8pt,
  before upper={\popboxhead{Fiche bilan}{poporange}{L'essentiel à connaître pour l'examen}}}
% Figure encadrée avec légende
\newtcolorbox{popfigure}[1][]{enhanced,breakable=false,colback=white,colframe=popline,boxrule=1.4pt,arc=8pt,
  left=10pt,right=10pt,top=10pt,bottom=8pt,before skip=10pt,after skip=12pt,halign=center,
  lower separated=false,halign lower=center,
  fontlower=\popsemi\fontsize{9}{11.5}\selectfont\color{popmuted},
  #1}
\newcommand{\legende}[1]{\par\vspace{6pt}{\popsemi\fontsize{9}{11.5}\selectfont\color{popmuted}#1}}

% ============================================================
% Signature OnBuch+
% ============================================================
\newcommand{\poplogoinline}[1]{{\poplogo\fontsize{#1}{#1}\selectfont\color{popink}OnBuch\color{poporange}+}}
\newcommand{\signatureonbuch}{%
  \begin{tcolorbox}[enhanced,colback=popink,colframe=popink,boxrule=2.2pt,arc=10pt,
      drop shadow=poporange,left=14pt,right=14pt,top=10pt,bottom=10pt,before skip=0pt,after skip=0pt]
    \tikz[baseline=-0.7ex]\node[circle,fill=popgreen,draw=popcream,line width=1.3pt,minimum size=22pt,inner sep=0]{\popcheck{white}};%
    \hspace{9pt}\begin{minipage}[c]{0.62\linewidth}
      {\popxbold\fontsize{12}{13}\selectfont\color{popcream}Signé\ \textbullet\ {\poplogo OnBuch\color{poporange}+}}\par\vspace{2pt}
      {\raggedright\popsemi\fontsize{8}{10.5}\selectfont\color{popline}\MakeUppercase{Équipe pédagogique OnBuch+}\\\MakeUppercase{Conforme au programme officiel MINESEC}\par}
    \end{minipage}\hfill
    {\poplogo\fontsize{22}{22}\selectfont\color{popcream}OnBuch\color{poporange}+}
  \end{tcolorbox}%
}

% ============================================================
% Page de garde
% #1 titre · #2 matière · #3 niveau · #4 module · #5 n° leçon · #6 durée ·
% #7 description / objectifs (texte ou itemize)
% ============================================================
\newcommand{\pagegarde}[7]{%
  \thispagestyle{empty}
  \newgeometry{a4paper,margin=1.7cm,top=1.6cm,bottom=1.4cm}%
  \begin{tikzpicture}[remember picture,overlay]
    \fill[poporange!18] ([xshift=-3.2cm,yshift=-3.6cm]current page.north east) circle (4.6cm);
    \fill[poppurple!14] ([xshift=2.4cm,yshift=7.6cm]current page.south west) circle (3.8cm);
  \end{tikzpicture}%
  {\poplogo\fontsize{30}{30}\selectfont\color{popink}OnBuch\color{poporange}+}\hfill
  \raisebox{0.6ex}{\poppill{Cours signé \textbullet\ Premium}{popink}{popcream}}\par
  \vspace{1pt}\popsquiggle{poppurple}\par
  \vspace{8pt}
  \begin{tcolorbox}[enhanced,colback=poporange,colframe=popink,boxrule=2.8pt,arc=13pt,
      drop shadow=popink,left=18pt,right=18pt,top=12pt,bottom=13pt,after skip=10pt]
    \poppill{#2 \textbullet\ #3}{popink}{popcream}\hspace{5pt}\poppill{#4}{white}{popink}\par\vspace{8pt}
    {\popdisplay\fontsize{14}{14}\selectfont\color{popink}LEÇON #5}\par\vspace{4pt}
    {\raggedright\hyphenpenalty=10000\exhyphenpenalty=10000\popdisplay\fontsize{25}{27}\selectfont\color{popcream}\MakeUppercase{#1}\par}
    \vspace{8pt}
    {\popxbold\fontsize{9.5}{9.5}\selectfont\color{popink}\MakeUppercase{Durée conseillée : #6}}
  \end{tcolorbox}
  \begin{tcolorbox}[enhanced,colback=white,colframe=popink,boxrule=2.2pt,arc=10pt,
      drop shadow=popink,left=16pt,right=16pt,top=10pt,bottom=10pt,after skip=8pt]
    \poppill{Dans cette leçon}{poppurple}{white}\par\vspace{5pt}
    \popsemi\fontsize{9.3}{12.6}\selectfont\color{popink}#7
  \end{tcolorbox}
  \noindent\begin{minipage}[t]{0.487\linewidth}
    \begin{tcolorbox}[enhanced,colback=popgreen,colframe=popink,boxrule=2.2pt,arc=10pt,
        drop shadow=popink,left=11pt,right=11pt,top=10pt,bottom=10pt,height=3.1cm,valign=center]
      \tcbox[colback=white,colframe=popink,boxrule=1.3pt,arc=5pt,boxsep=0pt,nobeforeafter,
        left=3pt,right=3pt,top=3pt,bottom=3pt,valign=center]{\qrcode[height=1.9cm]{https://play.google.com/store/apps/details?id=com.luvvix.onbuch&hl=fr}}%
      \hspace{8pt}\begin{minipage}[c]{0.5\linewidth}
        {\popdisplay\fontsize{11}{12.5}\selectfont\color{white}\MakeUppercase{Télécharge OnBuch}}\par\vspace{4pt}
        {\raggedright\popsemi\fontsize{8.4}{11}\selectfont\color{white}Gratuit \textbullet\ Google Play\\Tuteur IA, annales, concours, résultats\par}
      \end{minipage}
    \end{tcolorbox}
  \end{minipage}\hfill
  \begin{minipage}[t]{0.487\linewidth}
    \begin{tcolorbox}[enhanced,colback=poppurple,colframe=popink,boxrule=2.2pt,arc=10pt,
        drop shadow=popink,left=11pt,right=11pt,top=10pt,bottom=10pt,height=3.1cm,valign=center]
      \tikz[baseline=-0.7ex]\node[circle,fill=white,draw=popink,line width=1.3pt,minimum size=1.6cm,inner sep=0]{\popdisplay\fontsize{24}{24}\selectfont\color{poppurple}+};%
      \hspace{8pt}\begin{minipage}[c]{0.6\linewidth}
        {\popdisplay\fontsize{11}{12.5}\selectfont\color{white}\MakeUppercase{Passe à OnBuch+}}\par\vspace{4pt}
        {\raggedright\popsemi\fontsize{8.4}{11}\selectfont\color{white}Tous les cours signés, Tuteur IA illimité, fiches PDF — onglet OnBuch+ de l'app\par}
      \end{minipage}
    \end{tcolorbox}
  \end{minipage}
  \vspace{8pt}
  \popcontenu
  \vfill
  \signatureonbuch
  \restoregeometry
  \newpage
}

% Rangée « Ce cours contient » (page de garde)
\newcommand{\poptile}[3]{% #1 couleur #2 gros texte #3 légende
  \begin{tcolorbox}[enhanced,colback=#1,colframe=popink,boxrule=2pt,arc=9pt,shadow={2.5pt}{-2.5pt}{0pt}{popink},
      left=6pt,right=6pt,top=9pt,bottom=9pt,halign=center,valign=center,height=1.75cm,width=\linewidth,nobeforeafter]
    {\popdisplay\fontsize{12.5}{13}\selectfont\color{white}\MakeUppercase{#2}}\par\vspace{3pt}
    {\centering\popsemi\fontsize{7.8}{9.5}\selectfont\color{white}#3\par}
  \end{tcolorbox}}
\newcommand{\popcontenu}{%
  \par\noindent{\popxboldls\fontsize{7.5}{7.5}\selectfont\color{popmuted}CE COURS SIGNÉ CONTIENT}\par\vspace{7pt}
  \noindent\begin{minipage}[t]{0.235\linewidth}\poptile{popblue}{Cours}{complet, illustré, pas à pas}\end{minipage}\hfill
  \begin{minipage}[t]{0.235\linewidth}\poptile{poporange}{Méthodes}{et exercices résolus}\end{minipage}\hfill
  \begin{minipage}[t]{0.235\linewidth}\poptile{poppink}{Exercices}{type examen + corrigés}\end{minipage}\hfill
  \begin{minipage}[t]{0.235\linewidth}\poptile{popgreen}{Fiche bilan}{l'essentiel à retenir}\end{minipage}\par}

% Sommaire
\newtcolorbox{sommaire}{enhanced,colback=white,colframe=popink,boxrule=2.2pt,arc=10pt,
  drop shadow=popink,left=18pt,right=18pt,top=13pt,bottom=13pt,before skip=6pt,after skip=20pt,
  before upper={\poppill{Sommaire}{poppurple}{white}\par\vspace{9pt}\popsemi\fontsize{10.6}{14.5}\selectfont\color{popink}}}

% Signature de fin de cours — poussée tout en bas de la dernière page
\newcommand{\findecours}{%
  \par\vfill\nopagebreak
  \begin{center}\popsquiggle{poporange}\end{center}\vspace{4pt}
  \signatureonbuch
}
\AtBeginDocument{\def\llbracket{\mathopen{[\mkern-3.5mu[}}\def\rrbracket{\mathclose{]\mkern-3.5mu]}}} % intervalles d'entiers (glyphe ⟦ absent de la police)
% Glyphes absents de Fira Math : redéfinis à partir de glyphes disponibles
\AtBeginDocument{%
  \def\setminus{\mathbin{\backslash}}\let\smallsetminus\setminus
  \def\vdots{\mathord{\vbox{\baselineskip3.2pt\lineskiplimit0pt\kern4pt\hbox{.}\hbox{.}\hbox{.}}}}%
  \def\ddots{\mathinner{\mkern1mu\raise6.5pt\hbox{.}\mkern2mu\raise3.6pt\hbox{.}\mkern2mu\raise0.7pt\hbox{.}\mkern1mu}}%
  \def\triangle{\mathord{\scalebox{0.9}{$\symup{\Delta}$}}}%
  \def\nabla{\mathord{\raisebox{\depth}{\scalebox{1}[-1]{$\symup{\Delta}$}}}}%
  \def\checkmark{\popcheck{popdark}}%
  \def\star{\mathbin{\ast}}\def\bigstar{\mathord{\ast}}%
  \def\diamond{\mathbin{\scalebox{0.6}{$\Diamond$}}}%
  \def\bigcup{\mathop{\mathchoice{\raisebox{-0.3em}{\scalebox{1.7}{$\cup$}}}{\scalebox{1.25}{$\cup$}}{\cup}{\cup}}\displaylimits}%
  \def\bigcap{\mathop{\mathchoice{\raisebox{-0.3em}{\scalebox{1.7}{$\cap$}}}{\scalebox{1.25}{$\cap$}}{\cap}{\cap}}\displaylimits}%
  \def\lesseqqgtr{\mathrel{\lessgtr}}\def\bigoplus{\mathop{\oplus}\displaylimits}%
}
\providecommand{\ssi}{si et seulement si} % abréviation fréquente
\AtBeginDocument{\providecommand{\wideparen}{\overparen}} % arc de cercle

%% FIGFIX-BEGIN v5 — correctifs globaux des figures (docs/onbuch-generation/tools/figfix)
\usepackage{adjustbox}\usepackage{etoolbox}\tcbuselibrary{raster}
% toute figure TikZ/pgfplots plus large que la ligne est réduite à la largeur disponible
\BeforeBeginEnvironment{tikzpicture}{\begin{adjustbox}{max width=\linewidth}}
\AfterEndEnvironment{tikzpicture}{\end{adjustbox}}
% légendes pgfplots lisibles par-dessus les courbes
\pgfplotsset{every axis/.append style={legend style={fill=white,fill opacity=0.92,text opacity=1,draw=black!15,inner sep=3pt}}}
% étiquettes d'arêtes (syntaxe edge["..."]) avec fond blanc
\tikzset{every edge quotes/.append style={fill=white,inner sep=1.2pt}}
%% FIGFIX-END
\begin{document}

\pagegarde{Leçon 14 — Expression écrite : le paludisme}{Chinois}{Tle A}{Module 2 : Environnement et santé}{ZH14}{2 h}{Cette leçon d'expression écrite apprend à rédiger un texte simple en chinois sur le paludisme : symptômes, prévention et traitement. Elle entraîne aussi l'écriture correcte des caractères du thème et la traduction (thème et version), compétences clés de l'épreuve du baccalauréat.
\vspace{4pt}
\begin{multicols}{2}\small
\begin{itemize}
\item À la fin de cette leçon, je sais écrire correctement les caractères clés du thème du paludisme (radicaux, ordre des traits).
\item Je sais distinguer les caractères proches faciles à confondre (疟/疾/病, 蚊/虫, 治/活).
\item Je sais utiliser le vocabulaire essentiel : 疟疾, 发烧, 蚊子, 蚊帐, 预防, 治疗, 医院.
\item Je sais lire et comprendre un texte simple sur le paludisme.
\item Je sais structurer un texte avec les connecteurs 首先、然后、所以、因为、如果.
\item Je sais rédiger un texte de 5 à 8 phrases sur la prévention du paludisme.
\end{itemize}\end{multicols}}

\begin{sommaire}
\begin{multicols}{2}
\begin{enumerate}
\item Texte modèle et découverte du thème
\item Lexique du paludisme
\item Écriture correcte des caractères
\item Structure du texte et connecteurs
\item Thème et version (traduction)
\item Production écrite guidée et évaluation
\item Activité d'intégration
\item Méthodes et exercices résolus
\item Exercices
\item Corrigés des exercices
\item Fiche bilan
\end{enumerate}
\end{multicols}
\end{sommaire}

\begin{objectifs}
\begin{itemize}
\item À la fin de cette leçon, je sais écrire correctement les caractères clés du thème du paludisme (radicaux, ordre des traits).
\item Je sais distinguer les caractères proches faciles à confondre (疟/疾/病, 蚊/虫, 治/活).
\item Je sais utiliser le vocabulaire essentiel : 疟疾, 发烧, 蚊子, 蚊帐, 预防, 治疗, 医院.
\item Je sais lire et comprendre un texte simple sur le paludisme.
\item Je sais structurer un texte avec les connecteurs 首先、然后、所以、因为、如果.
\item Je sais rédiger un texte de 5 à 8 phrases sur la prévention du paludisme.
\item Je sais traduire des phrases simples du français vers le chinois (thème).
\item Je sais traduire un court texte du chinois vers le français (version).
\item Je sais m'auto-évaluer à l'aide d'une grille d'évaluation type examen.
\end{itemize}
\end{objectifs}

\begin{prerequis}
\begin{itemize}
\item Vocabulaire du corps et de la santé vu en Première (头、疼、医生、医院、药)
\item Structure de base Sujet–Verbe–Objet et phrase avec 是
\item Expression de la cause avec 因为…所以… (révision)
\item Négation avec 不 et 没有
\item Radicaux de base (疒 maladie, 虫 insecte) rencontrés dans les leçons précédentes
\end{itemize}
\end{prerequis}

\newpage

\coursec{Texte modèle et découverte du thème}

\courssub{Situation d'accroche et problématique}

À Yaoundé comme à Douala, la saison des pluies ramène les moustiques… et le paludisme. Dans les quartiers, on entend le soir le bourdonnement familier des \zh{蚊子} (\emph{wénzi}, moustiques), et au centre de santé de ton quartier, une affiche en chinois explique comment se protéger : dormir sous moustiquaire, consulter vite en cas de fièvre. Un ami chinois en visite au Cameroun te demande : \zh{你们怎么预防疟疾？} (\emph{Nǐmen zěnme yùfáng nüèjí ?}) — « Comment prévenez-vous le paludisme ? »

Sauras-tu lui répondre \textbf{par écrit}, en chinois correct, avec des caractères bien tracés et des connecteurs clairs ? C'est exactement ce que cette leçon t'apprendra. Mais avant d'écrire, il faut d'abord \textbf{lire et comprendre} un texte modèle : c'est la première étape de toute production écrite réussie.

\begin{aretenir}[Problématique de la leçon]
Comment lire, comprendre puis imiter un texte chinois simple sur un sujet de santé (le paludisme) afin de produire ensuite son propre texte ? Nous allons : lire le texte modèle, repérer le vocabulaire nouveau, comprendre sa structure (introduction, développement, conclusion), puis nous préparer à écrire.
\end{aretenir}

\courssub{Lecture du texte modèle}

Voici le texte modèle de cette leçon. Lis-le d'abord en silence, sans le pinyin, pour tester ta lecture des caractères. Puis lis-le une deuxième fois avec le pinyin, et enfin avec la traduction.

\begin{textebox}[疟疾是一种常见的病 — Le paludisme est une maladie courante]
疟疾是一种常见的病。\\
\emph{Nüèjí shì yì zhǒng chángjiàn de bìng.}\\
« Le paludisme est une maladie courante. »\\[0.4em]
蚊子把病毒传给人。\\
\emph{Wénzi bǎ bìngdú chuángěi rén.}\\
« Le moustique transmet le virus aux hommes. »\\[0.4em]
得了疟疾的人会发烧、头疼。\\
\emph{Déle nüèjí de rén huì fāshāo, tóuténg.}\\
« Les personnes qui ont attrapé le paludisme ont de la fièvre et mal à la tête. »\\[0.4em]
我们应该睡在蚊帐里，并且去医院治疗。\\
\emph{Wǒmen yīnggāi shuì zài wénzhàng lǐ, bìngqiě qù yīyuàn zhìliáo.}\\
« Nous devons dormir sous une moustiquaire, et de plus aller à l'hôpital pour nous faire soigner. »\\[0.4em]
预防疟疾很重要。\\
\emph{Yùfáng nüèjí hěn zhòngyào.}\\
« Prévenir le paludisme est très important. »
\end{textebox}

\begin{methode}[Comment lire un texte chinois : la méthode en 4 étapes]
Face à un texte chinois inconnu, ne panique jamais. Suis toujours le même ordre :
\begin{enumerate}
\item \textbf{Lire le titre.} Le titre annonce le thème. Ici : \zh{疟疾是一种常见的病} — on sait déjà qu'on parlera d'une maladie (\zh{病}).
\item \textbf{Repérer les mots connus.} Souligne mentalement les mots déjà appris : \zh{是} (être), \zh{人} (personne), \zh{我们} (nous), \zh{应该} (devoir), \zh{去} (aller), \zh{医院} (hôpital), \zh{很} (très), \zh{重要} (important). Ces mots forment le « squelette » du texte.
\item \textbf{Deviner les mots nouveaux par le contexte.} Exemple : \zh{蚊子把病毒传给人} — tu connais \zh{人} ; le mot \zh{传} précédé de \zh{把} évoque une action faite sur un objet ; avec \zh{蚊子} (que le contexte « maladie + saison des pluies » aide à deviner), tu peux supposer « transmettre ».
\item \textbf{Vérifier au glossaire.} Confirme tes hypothèses dans le tableau de vocabulaire ci-dessous. Ne devine jamais sans vérifier : c'est ainsi qu'on évite les contresens.
\end{enumerate}
\end{methode}

\courssub{Repérage des mots nouveaux et glossaire}

Dans le texte modèle, les \cle{mots nouveaux} à surligner dans ton cahier sont : \zh{疟疾}, \zh{常见}, \zh{病}, \zh{蚊子}, \zh{病毒}, \zh{传给}, \zh{发烧}, \zh{头疼}, \zh{蚊帐}, \zh{并且}, \zh{治疗}, \zh{预防}. Voici le glossaire complet :

\begin{center}
\begin{tabularx}{\linewidth}{|X|X|X|X|}
\hline
\rowcolor{popblueL} Caractères & Pinyin & Nature & Traduction \\
\hline
疟疾 & nüèjí & nom & paludisme, malaria \\
\hline
种 & zhǒng & classificateur / nom & espèce, sorte \\
\hline
常见 & chángjiàn & adjectif & courant, fréquent \\
\hline
病 & bìng & nom & maladie \\
\hline
蚊子 & wénzi & nom & moustique \\
\hline
把 & bǎ & préposition & (marque de l'objet déplacé) \\
\hline
病毒 & bìngdú & nom & virus, agent pathogène \\
\hline
传给 & chuángěi & verbe & transmettre à \\
\hline
得了 & déle & verbe + particule & attraper (une maladie) \\
\hline
会 & huì & auxiliaire & pouvoir, risquer de \\
\hline
发烧 & fāshāo & verbe & avoir de la fièvre \\
\hline
头疼 & tóuténg & verbe / adjectif & avoir mal à la tête \\
\hline
应该 & yīnggāi & auxiliaire & devoir (il faut) \\
\hline
蚊帐 & wénzhàng & nom & moustiquaire \\
\hline
并且 & bìngqiě & conjonction & et de plus, en outre \\
\hline
医院 & yīyuàn & nom & hôpital \\
\hline
治疗 & zhìliáo & verbe & soigner, traiter \\
\hline
预防 & yùfáng & verbe & prévenir \\
\hline
重要 & zhòngyào & adjectif & important \\
\hline
\end{tabularx}
\end{center}

Observe bien la famille de mots construite sur \zh{病} (\emph{bìng}, maladie) : \zh{病毒} (virus), \zh{病人} (malade), \zh{看病} (consulter un médecin), \zh{生病} (tomber malade). Le chinois construit son vocabulaire médical par \textbf{composition} : chaque caractère garde son sens et s'assemble à un autre, comme des briques. De même, \zh{蚊} (moustique) donne \zh{蚊子} (le moustique) et \zh{蚊帐} (littéralement « tente à moustiques », donc moustiquaire).

\begin{savaistu}[Le radical des maladies : 疒]
Le caractère \zh{疟} (\emph{nüè}) contient le radical \zh{疒} (\emph{nè}), appelé « radical de la maladie ». Ce radical apparaît dans presque tous les caractères liés aux maladies : \zh{病} (bìng, maladie), \zh{疼} (téng, douloureux), \zh{疾} (jí, maladie, dans \zh{疾病}), \zh{痛} (tòng, douleur). Astuce de mémorisation : quand tu vois \zh{疒} au-dessus ou à gauche d'un caractère, pense « santé, maladie ». Le reste du caractère donne souvent le son : dans \zh{疟}, la partie \zh{虐} indique la prononciation \emph{nüè}. C'est le principe des caractères « sémantico-phonétiques », qui représentent la grande majorité des sinogrammes.
\end{savaistu}

\courssub{Questions de compréhension du texte}

Après la lecture, vérifie ta compréhension globale. Réponds d'abord en français, puis essaie en chinois avec les mots du texte.

\begin{exoresolu}[Comprendre le texte modèle]
\textbf{Énoncé.} Réponds aux quatre questions suivantes d'après le texte \zh{疟疾是一种常见的病} : 1) De quelle maladie parle le texte ? 2) Qui transmet cette maladie ? 3) Quels sont les symptômes ? 4) Comment se protéger et se soigner ?
\tcblower
\textbf{Solution détaillée.}
\begin{enumerate}
\item Le texte parle du \cle{paludisme} : \zh{疟疾} (\emph{nüèjí}). Première phrase : \zh{疟疾是一种常见的病。} « Le paludisme est une maladie courante. »
\item C'est le \cle{moustique} qui transmet la maladie : \zh{蚊子把病毒传给人。} « Le moustique transmet le virus aux hommes. »
\item Les symptômes sont la \cle{fièvre} et le \cle{mal de tête} : \zh{得了疟疾的人会发烧、头疼。} « Les malades ont de la fièvre et mal à la tête. »
\item Pour se protéger : \zh{睡在蚊帐里} « dormir sous une moustiquaire » ; pour se soigner : \zh{去医院治疗} « aller à l'hôpital pour se faire soigner ». La phrase de conclusion résume : \zh{预防疟疾很重要。} « Prévenir le paludisme est très important. »
\end{enumerate}
\end{exoresolu}

\courssub{Commentaire du modèle : la structure du texte}

Un bon texte chinois simple, comme un bon texte français, suit un plan en \textbf{trois parties}. Observons comment le texte modèle est construit :

\begin{center}
\begin{tabularx}{\linewidth}{|X|X|X|}
\hline
\rowcolor{popblueL} Partie & Phrase du texte & Fonction \\
\hline
Titre & 疟疾是一种常见的病 & Annoncer le sujet \\
\hline
Introduction & 疟疾是一种常见的病。 & Présenter le problème (définition) \\
\hline
Développement 1 & 蚊子把病毒传给人。 & La cause : le moustique \\
\hline
Développement 2 & 得了疟疾的人会发烧、头疼。 & Les symptômes \\
\hline
Conclusion & 我们应该睡在蚊帐里，并且去医院治疗。预防疟疾很重要。 & Les conseils (protection + traitement) \\
\hline
\end{tabularx}
\end{center}

\begin{propriete}[Le plan type d'un texte informatif simple en chinois]
Pour écrire un court texte informatif (santé, environnement, vie quotidienne), retiens ce schéma en trois temps :
\begin{enumerate}
\item \textbf{Introduction} : présenter le sujet avec la structure \zh{X 是 yì zhǒng … de …} « X est une sorte de… ».
\item \textbf{Développement} : donner la cause, puis les effets ou symptômes, avec des phrases courtes Sujet + Verbe + Complément.
\item \textbf{Conclusion} : donner un conseil avec \zh{我们应该…} « nous devons… », éventuellement renforcé par \zh{并且} « et de plus », puis une phrase d'importance : \zh{…很重要。}
\end{enumerate}
Ce plan sera ta « grille » pour la production écrite de la fin de la leçon.
\end{propriete}

Deux structures grammaticales du texte méritent une observation attentive dès maintenant (elles seront réutilisées dans ta production) :

\begin{exemplebox}[Deux structures-clés du texte]
\textbf{1. La construction avec \zh{把}} : \zh{蚊子把病毒传给人。} (\emph{Wénzi bǎ bìngdú chuángěi rén.}) — Schéma : Sujet + \zh{把} + Objet + Verbe + complément. Le français dirait « le moustique transmet le virus à l'homme » ; le chinois place l'objet (\zh{病毒}) \emph{avant} le verbe grâce à \zh{把}.\\[0.3em]
\textbf{2. La subordonnée avec \zh{的}} : \zh{得了疟疾的人} (\emph{déle nüèjí de rén}) — « les personnes \emph{qui ont attrapé le paludisme} ». En chinois, toute la proposition relative se place \textbf{avant} le nom, suivie de \zh{的} : littéralement « attrapé-paludisme-DE-personnes ». C'est l'inverse du français !
\end{exemplebox}

\begin{attention}[Pièges fréquents des francophones]
\begin{itemize}
\item \textbf{Ne traduis pas mot à mot} « avoir de la fièvre » par \zh{有发烧} : on dit \zh{发烧} (\emph{fāshāo}), verbe à part entière : \zh{他发烧了。} « Il a de la fièvre. »
\item \textbf{Attention à \zh{会}} : dans \zh{会发烧}, il n'exprime pas le futur mais la possibilité probable (« risque d'avoir »).
\item \textbf{Ordre de la relative} : ne dis jamais \zh{人得了疟疾的} ; la bonne forme est \zh{得了疟疾的人} (relative + \zh{的} + nom).
\item \textbf{Énumération} : entre deux symptômes, on utilise la virgule chinoise d'énumération \zh{、} : \zh{发烧、头疼}, et non \zh{和} qui relie plutôt des noms en sujet ou en objet.
\end{itemize}
\end{attention}

\courssub{Carte mentale : organiser les idées du texte}

Avant d'écrire, les bons élèves visualisent le contenu. Voici la carte mentale du texte modèle : elle montre les quatre « branches » du thème \zh{疟疾} — la cause, les symptômes, la prévention, le traitement. Tu la réutiliseras comme plan pour ta propre rédaction.

\begin{popfigure}
\begin{tikzpicture}[
  mindmap,
  grow cyclic,
  every node/.style={concept, align=center, font=\small},
  root concept/.append style={concept color=popblue, font=\large\bfseries},
  level 1/.append style={level distance=3.2cm, sibling angle=90},
  level 2/.append style={level distance=2.2cm, sibling angle=50},
]
\node[root concept, align=center]{疟疾\\ \emph{nüèjí}\\ paludisme}
  child[concept color=poporange]{ node[align=center]{原因\\ \emph{yuányīn}\\ cause}
    child[concept color=poporangeL]{ node[align=center]{蚊子\\ \emph{wénzi}\\ moustique} }
  }
  child[concept color=poppink]{ node[align=center]{症状\\ \emph{zhèngzhuàng}\\ symptômes}
    child[concept color=poppinkL]{ node[align=center]{发烧\\ \emph{fāshāo}\\ fièvre} }
    child[concept color=poppinkL]{ node[align=center]{头疼\\ \emph{tóuténg}\\ mal de tête} }
  }
  child[concept color=popgreen]{ node[align=center]{预防\\ \emph{yùfáng}\\ prévention}
    child[concept color=popgreenL]{ node[align=center]{蚊帐\\ \emph{wénzhàng}\\ moustiquaire} }
  }
  child[concept color=poppurple]{ node[align=center]{治疗\\ \emph{zhìliáo}\\ traitement}
    child[concept color=poppurpleL]{ node[align=center]{医院\\ \emph{yīyuàn}\\ hôpital} }
    child[concept color=poppurpleL]{ node[align=center]{药\\ \emph{yào}\\ médicaments} }
  };
\end{tikzpicture}
\legende{Carte mentale du texte modèle : les quatre branches du thème 疟疾 (cause, symptômes, prévention, traitement).}
\end{popfigure}

\begin{experience}[Activité de classe : lecture à voix haute et repérage]
Travail en binôme (10 minutes) :
\begin{enumerate}
\item Lis le texte modèle à voix haute, en faisant attention aux tons (surtout \zh{nüèjí}, 4\up{e} + 2\up{e} ton, et \zh{wénzi}, 2\up{e} ton + ton neutre). Ton binôme corrige ta prononciation.
\item Surligne dans le texte les douze mots nouveaux du glossaire.
\item Sans regarder la traduction, raconte le texte en français à ton binôme en suivant la carte mentale : cause, symptômes, prévention, traitement.
\item Question flash du professeur : \zh{我们怎么预防疟疾？} — réponds à l'oral avec une phrase complète : \zh{我们应该睡在蚊帐里。}
\end{enumerate}
\end{experience}

\begin{aretenir}[Ce qu'il faut retenir de cette première étape]
\begin{itemize}
\item Le texte modèle \zh{疟疾是一种常见的病} suit un plan en trois parties : \cle{introduction} (présenter le problème), \cle{développement} (cause + symptômes), \cle{conclusion} (conseils avec \zh{应该} et \zh{并且}).
\item La méthode de lecture en 4 étapes : titre, mots connus, devinette par le contexte, vérification au glossaire.
\item Le vocabulaire médical chinois se construit par composition autour de \zh{病} ; le radical \zh{疒} signale les caractères de maladie.
\item Deux structures à réutiliser : Sujet + \zh{把} + Objet + Verbe, et la relative placée avant le nom avec \zh{的}.
\end{itemize}
\end{aretenir}

\coursec{Lexique du paludisme}

Dans la section précédente, nous avons lu et commenté un texte modèle sur le paludisme au Cameroun. Vous avez remarqué que certains mots revenaient souvent : \zh{疟疾}, \zh{蚊子}, \zh{发烧}, \zh{预防}… Pour pouvoir, à la fin de cette leçon, rédiger votre propre texte, il faut d'abord \cle{maîtriser ce vocabulaire} : le reconnaître, le prononcer, l'écrire et savoir l'employer dans une phrase correcte. C'est l'objectif de cette deuxième partie : construire, pas à pas, votre \emph{boîte à outils lexicale} sur le thème de la santé et du paludisme.

\courssub{Observation : les mots dans leur contexte}

Avant tout tableau, observons. Lisez ce court paragraphe et soulignez mentalement les mots qui parlent de la maladie, de sa transmission, de ses symptômes et de sa prise en charge.

\begin{textebox}[Mini-texte d'observation]
\zh{在喀麦隆，疟疾是一种很常见的病。蚊子传播疟疾。得了疟疾的人会发烧、头疼。为了预防疟疾，我们应该睡在蚊帐里。如果生病了，要去医院看病，医生会给你药。}\\
\emph{Zài Kāmàilóng, nüèjí shì yì zhǒng hěn chángjiàn de bìng. Wénzi chuánbō nüèjí. Dé le nüèjí de rén huì fāshāo, tóuténg. Wèile yùfáng nüèjí, wǒmen yīnggāi shuì zài wénzhàng lǐ. Rúguǒ shēngbìng le, yào qù yīyuàn kànbìng, yīshēng huì gěi nǐ yào.}\\
« Au Cameroun, le paludisme est une maladie très courante. Les moustiques transmettent le paludisme. Les personnes qui ont attrapé le paludisme ont de la fièvre et mal à la tête. Pour prévenir le paludisme, nous devons dormir sous une moustiquaire. Si l'on tombe malade, il faut aller à l'hôpital consulter, le médecin vous donnera des médicaments. »
\end{textebox}

Vous l'avez constaté : en dix lignes, presque tout notre lexique apparaît. Chaque mot a une \cle{fonction précise} : nommer la maladie, nommer le vecteur, décrire les symptômes, parler de la prévention ou du soin. Organisons maintenant ce vocabulaire.

\courssub{Glossaire du thème}

\begin{definition}[Les dix mots-clés de la leçon]
Voici le vocabulaire officiel du thème « le paludisme ». Apprenez chaque mot avec ses \cle{trois identités} : ses caractères, son pinyin (avec les tons !) et son sens.
\end{definition}

\begin{center}
\begin{tabularx}{\linewidth}{|X|X|X|X|}
\hline
\rowcolor{popblueL} Caractères & Pinyin & Nature & Traduction \\
\hline
疟疾 & nüèjí & n. & paludisme \\
\hline
蚊子 & wénzi & n. & moustique \\
\hline
发烧 & fāshāo & v. & avoir de la fièvre \\
\hline
头疼 & tóuténg & v. & avoir mal à la tête \\
\hline
蚊帐 & wénzhàng & n. & moustiquaire \\
\hline
预防 & yùfáng & v. & prévenir \\
\hline
治疗 & zhìliáo & v. & traiter, soigner \\
\hline
病毒 & bìngdú & n. & virus \\
\hline
医院 & yīyuàn & n. & hôpital \\
\hline
药 & yào & n. & médicament \\
\hline
\end{tabularx}
\end{center}

Quelques remarques de prononciation pour les francophones :

\begin{itemize}
\item \zh{疟疾} \emph{nüèjí} : le son \emph{üe} n'existe pas en français ; arrondissez les lèvres comme pour dire « u » tout en prononçant « i ». Premier caractère au 4\up{ème} ton, deuxième au 2\up{ème} ton : \emph{nüè-jí}.
\item \zh{蚊子} \emph{wénzi} : le deuxième caractère \zh{子} est au \cle{ton neutre}, très léger, comme souvent dans les noms d'animaux et d'objets (\zh{孩子} háizi, « enfant »).
\item \zh{头疼} \emph{tóuténg} : attention à la finale \emph{-eng} du deuxième caractère, nasale, différente de \emph{-en}.
\item \zh{治疗} \emph{zhìliáo} : les initiales \emph{zh-} et \emph{l-} se succèdent ; articulez bien les deux syllabes séparément.
\item \zh{病毒} \emph{bìngdú} : le \emph{u} de \emph{dú} se prononce comme le « ou » français « roue », pas comme le « u » de « lu ».
\end{itemize}

\begin{definition}[预防 yùfáng : prévenir]
\zh{预防} signifie « prévenir », c'est-à-dire \cle{prendre des mesures avant} que la maladie n'arrive. Structure : \textbf{预防 + nom de maladie}. Exemple : \zh{预防疟疾} \emph{yùfáng nüèjí} « prévenir le paludisme ». Ne le confondez pas avec \zh{治疗} \emph{zhìliáo} « traiter », qui intervient \emph{après} l'apparition de la maladie. La logique temporelle est : d'abord \zh{预防} (avant), ensuite \zh{治疗} (après).
\end{definition}

\begin{savaistu}[Paludisme : virus ou parasite ?]
Scientifiquement, le paludisme n'est pas causé par un \zh{病毒} (virus) mais par un \emph{parasite} du genre \emph{Plasmodium} (en chinois : \zh{疟原虫} \emph{nüèyuánchóng}). Dans le langage courant ou scolaire (niveau HSK), le mot \zh{病毒} est parfois utilisé de façon générique pour « agent pathogène / germe ». Dans un texte rigoureux, on dira : \zh{蚊子传播疟疾} (les moustiques transmettent le paludisme) ou \zh{蚊子传播疟原虫}. Retenez \zh{病毒} pour le vocabulaire général (grippe, COVID, etc.), mais sachez la distinction pour le paludisme.
\end{savaistu}

\begin{savaistu}[Une moustiquaire, littéralement une « tente à moustiques »]
Le mot \zh{蚊帐} \emph{wénzhàng} se compose de \zh{蚊} « moustique » et de \zh{帐} « tente, rideau, moustiquaire » : littéralement « tente à moustiques ». Le chinois construit très souvent ses mots ainsi, par combinaison logique : \zh{医院} « hôpital » = \zh{医} « médecine » + \zh{院} « bâtiment, cour » ; \zh{病毒} « virus » = \zh{病} « maladie » + \zh{毒} « poison ». Quand vous rencontrez un mot nouveau, essayez de deviner son sens à partir de ses composants : c'est une excellente stratégie de lecture.
\end{savaistu}

\courssub{Les collocations : des mots qui vont ensemble}

Un mot appris seul ne sert à rien si l'on ne sait pas avec quel verbe il s'emploie. En chinois comme en français, certains couples verbe + nom sont \cle{fixes} : on les appelle des collocations. Les voici, classées selon la logique du texte (cause $\rightarrow$ symptômes $\rightarrow$ prévention $\rightarrow$ soin) :

\begin{center}
\begin{tabularx}{\linewidth}{|X|X|X|}
\hline
\rowcolor{popblueL} Collocation & Pinyin & Traduction \\
\hline
得疟疾 & dé nüèjí & attraper le paludisme \\
\hline
传播疟疾 & chuánbō nüèjí & transmettre le paludisme \\
\hline
发烧、头疼 & fāshāo, tóuténg & avoir de la fièvre, mal à la tête \\
\hline
预防疟疾 & yùfáng nüèjí & prévenir le paludisme \\
\hline
睡在蚊帐里 & shuì zài wénzhàng lǐ & dormir sous une moustiquaire \\
\hline
去医院看病 & qù yīyuàn kànbìng & aller à l'hôpital consulter \\
\hline
吃药 & chī yào & prendre des médicaments \\
\hline
治疗疟疾 & zhìliáo nüèjí & traiter le paludisme \\
\hline
\end{tabularx}
\end{center}

\begin{exemplebox}[Des exemples à imiter]
\zh{他得了疟疾，发烧了。}\\
\emph{Tā dé le nüèjí, fāshāo le.}\\
« Il a attrapé le paludisme, il a de la fièvre. »\\[4pt]
\zh{蚊子传播疟疾。}\\
\emph{Wénzi chuánbō nüèjí.}\\
« Les moustiques transmettent le paludisme. »\\[4pt]
\zh{我们应该预防疟疾。}\\
\emph{Wǒmen yīnggāi yùfáng nüèjí.}\\
« Nous devons prévenir le paludisme. »\\[4pt]
\zh{妈妈生病了，她去医院看病。}\\
\emph{Māma shēngbìng le, tā qù yīyuàn kànbìng.}\\
« Maman est tombée malade, elle est allée à l'hôpital consulter. »\\[4pt]
\zh{每天晚上，我都睡在蚊帐里。}\\
\emph{Měi tiān wǎnshang, wǒ dōu shuì zài wénzhàng lǐ.}\\
« Chaque soir, je dors sous une moustiquaire. »\\[4pt]
\zh{医生给他开药，他按时吃药。}\\
\emph{Yīshēng gěi tā kāi yào, tā ànshí chī yào.}\\
« Le médecin lui a prescrit des médicaments, il prend ses médicaments à l'heure. »
\end{exemplebox}

\begin{attention}[得疟疾 ou 有疟疾 ?]
Pour dire « attraper une maladie » ou « être malade de », le chinois utilise le verbe \zh{得} \emph{dé} : \zh{得疟疾} « attraper le paludisme », \zh{得病} « tomber malade ». \textbf{N'utilisez pas} \zh{有} pour traduire « avoir une maladie » à propos d'une personne ! La phrase \zh{他有疟疾} ne veut pas dire « il a le paludisme » (état de santé) mais « il y a du paludisme (chez lui / dans cette région) » (constat d'existence). Comparez :\\
\zh{他得了疟疾。} \emph{Tā dé le nüèjí.} « Il a attrapé le paludisme / Il est atteint de paludisme. » (la personne est malade)\\
\zh{这个地区有疟疾。} \emph{Zhège dìqū yǒu nüèjí.} « Il y a du paludisme dans cette région. » (constat d'existence)\\
Autre piège fréquent : \zh{吃药} se construit avec \zh{吃} « manger » — en chinois, on « mange » les médicaments, on ne les « prend » pas (\zh{拿} \emph{ná}) ni ne les « boit » (\zh{喝} \emph{hē}). Ne dites jamais \zh{拿药} pour « prendre un traitement » (cela signifie « aller chercher/prendre le médicament physiquement »).
\end{attention}

\courssub{Le champ lexical associé (révision et extension HSK 2-3)}

Autour de nos dix mots-clés gravitent des mots déjà connus qui complètent le champ lexical de la santé et permettent de construire un raisonnement complet :

\begin{center}
\begin{tabularx}{\linewidth}{|X|X|X|X|}
\hline
\rowcolor{popblueL} Caractères & Pinyin & Nature & Traduction \\
\hline
生病 & shēngbìng & v. & tomber malade \\
\hline
健康 & jiànkāng & n./adj. & santé ; en bonne santé \\
\hline
医生 & yīshēng & n. & médecin \\
\hline
病人 & bìngrén & n. & malade, patient \\
\hline
休息 & xiūxi & v. & se reposer \\
\hline
症状 & zhèngzhuàng & n. & symptôme \\
\hline
按时 & ànshí & adv. & à l'heure, ponctuellement \\
\hline
开药 & kāi yào & v.o. & prescrire des médicaments \\
\hline
\end{tabularx}
\end{center}

Ces mots permettent de structurer le \cle{raisonnement complet} d'un texte sur le paludisme : la cause (\zh{蚊子}, \zh{传播}), la maladie (\zh{得疟疾}, \zh{生病}), les symptômes (\zh{发烧}, \zh{头疼}, \zh{症状}), la prévention (\zh{预防}, \zh{蚊帐}) et le soin (\zh{医院}, \zh{医生}, \zh{病人}, \zh{治疗}, \zh{开药}, \zh{吃药}, \zh{休息}). La figure ci-dessous classe les mots en deux grandes familles.

\begin{popfigure}
\begin{center}
\begin{tikzpicture}[
  box/.style={draw=popline, rounded corners=3pt, fill=popcream, align=center, minimum width=3.6cm, minimum height=0.85cm, font=\small},
  head/.style={draw=popblue, rounded corners=3pt, fill=popblueL, align=center, minimum width=3.8cm, minimum height=1.0cm, font=\small\bfseries},
  head2/.style={draw=popgreen, rounded corners=3pt, fill=popgreenL, align=center, minimum width=3.8cm, minimum height=1.0cm, font=\small\bfseries}
]
\node[head, align=center] (s) at (0,0){病因与症状\\[-1pt] \footnotesize Cause et symptômes};
\node[box] (s1) at (0,-1.3) {蚊子 wénzi};
\node[box] (s2) at (0,-2.3) {传播疟疾 chuánbō nüèjí};
\node[box] (s3) at (0,-3.3) {得疟疾 dé nüèjí};
\node[box] (s4) at (0,-4.3) {发烧 fāshāo / 头疼 tóuténg};
\node[box] (s5) at (0,-5.3) {症状 zhèngzhuàng};
\node[head2, align=center] (p) at (7,0){预防与治疗\\[-1pt] \footnotesize Prévention et soin};
\node[box] (p1) at (7,-1.3) {预防疟疾 yùfáng nüèjí};
\node[box] (p2) at (7,-2.3) {睡在蚊帐里 shuì zài wénzhàng lǐ};
\node[box] (p3) at (7,-3.3) {去医院看病 qù yīyuàn kànbìng};
\node[box] (p4) at (7,-4.3) {医生开药 yīshēng kāi yào};
\node[box] (p5) at (7,-5.3) {吃药 chī yào / 休息 xiūxi};
\draw[-{Stealth}, popblue, thick] (s) -- (s1);
\draw[-{Stealth}, popgreen, thick] (p) -- (p1);
\end{tikzpicture}
\end{center}
\legende{Carte mentale lexicale : deux familles de mots autour du paludisme (cause/symptômes vs prévention/soin).}
\end{popfigure}

\courssub{Exercice résolu et entraînement}

\begin{exoresolu}[Compléter avec le bon mot]
Complétez les phrases avec : \zh{蚊子} — \zh{发烧} — \zh{蚊帐} — \zh{预防} — \zh{药} — \zh{吃药}.\\
1. 晚上我要睡在\blank{}\blank{}里。\\
2. 他\blank{}\blank{}了，头也很疼。\\
3. 我们应该\blank{}\blank{}疟疾。\\
4. 医生给了他一些\blank{}，让他按时\blank{}\blank{}。\\
5. \blank{}\blank{}传播疟疾。
\tcblower
\textbf{Solutions :}\\
1. \zh{蚊帐} (wénzhàng) : on dort « dans » une moustiquaire, d'où la structure \zh{睡在蚊帐里}.\\
2. \zh{发烧} (fāshāo) : la présence de \zh{了} (changement d'état) et de \zh{头也很疼} indique l'apparition d'un symptôme.\\
3. \zh{预防} (yùfáng) : structure \zh{应该 + verbe} ; seul « prévenir » a un sens logique avec \zh{疟疾} dans une optique de santé publique.\\
4. \zh{药} (yào) / \zh{吃药} (chī yào) : le médecin donne des médicaments (\zh{给...药}), le patient les prend (\zh{吃药}).\\
5. \zh{蚊子} (wénzi) : sujet de la phrase ; c'est le vecteur qui transmet la maladie.\\
\textbf{Traductions :} 1. Le soir, je dors sous une moustiquaire. 2. Il a de la fièvre et il a aussi mal à la tête. 3. Nous devons prévenir le paludisme. 4. Le médecin lui a donné des médicaments pour qu'il les prenne à l'heure. 5. Les moustiques transmettent le paludisme.
\end{exoresolu}

\begin{exercice}[Association, dictée et production]
\textbf{A. Association mot-sens.} Reliez chaque mot chinois à sa traduction :\\
\zh{医院}, \zh{病毒}, \zh{治疗}, \zh{吃药}, \zh{得疟疾}, \zh{传播}, \zh{症状}, \zh{按时} \\
/ a. traiter, b. attraper le paludisme, c. hôpital, d. prendre des médicaments, e. virus, f. transmettre, g. symptôme, h. à l'heure.\\

\textbf{B. Dictée de 10 mots.} Votre professeur dicte en chinois (ou en français) : \zh{疟疾}, \zh{蚊子}, \zh{发烧}, \zh{头疼}, \zh{蚊帐}, \zh{预防}, \zh{治疗}, \zh{医院}, \zh{药}, \zh{症状}. Écrivez les caractères, puis ajoutez le pinyin avec les tons.\\

\textbf{C. Phrases libres.} Écrivez une phrase avec la collocation \zh{睡在蚊帐里} et une phrase avec \zh{去医院看病}. Utilisez si possible un connecteur logique (\zh{因为}...\zh{所以}..., \zh{如果}...\zh{就}...).
\end{exercice}

\begin{corrige}[Exercice : association, dictée et production]
\textbf{A.} \zh{医院} → c (hôpital) ; \zh{病毒} → e (virus) ; \zh{治疗} → a (traiter) ; \zh{吃药} → d (prendre des médicaments) ; \zh{得疟疾} → b (attraper le paludisme) ; \zh{传播} → f (transmettre) ; \zh{症状} → g (symptôme) ; \zh{按时} → h (à l'heure).\\

\textbf{B.} Vérifiez chaque caractère et chaque ton soigneusement :\\
\zh{疟疾} nüèjí (4\up{ème} + 2\up{ème} ton) — \zh{蚊子} wénzi (2\up{ème} ton + neutre) — \zh{发烧} fāshāo (1\up{er} + 1\up{er}) — \zh{头疼} tóuténg (2\up{ème} + 2\up{ème}) — \zh{蚊帐} wénzhàng (2\up{ème} + 4\up{ème}) — \zh{预防} yùfáng (4\up{ème} + 2\up{ème}) — \zh{治疗} zhìliáo (4\up{ème} + 2\up{ème}) — \zh{医院} yīyuàn (1\up{er} + 4\up{ème}) — \zh{药} yào (4\up{ème} ton) — \zh{症状} zhèngzhuàng (4\up{ème} + 4\up{ème} ton).\\

\textbf{C.} Modèles possibles :\\
\zh{因为疟疾很危险，所以我每天晚上都睡在蚊帐里。} \emph{Yīnwèi nüèjí hěn wéixiǎn, suǒyǐ wǒ měi tiān wǎnshang dōu shuì zài wénzhàng lǐ.} « Comme le paludisme est dangereux, je dors sous une moustiquaire chaque soir. »\\
\zh{如果发烧头疼，就要去医院看病。} \emph{Rúguǒ fāshāo tóuténg, jiù yào qù yīyuàn kànbìng.} « Si on a de la fièvre et mal à la tête, il faut aller à l'hôpital consulter. »
\end{corrige}

\begin{aretenir}[L'essentiel du lexique]
\begin{itemize}
\item Dix mots-clés : \zh{疟疾}, \zh{蚊子}, \zh{发烧}, \zh{头疼}, \zh{蚊帐}, \zh{预防}, \zh{治疗}, \zh{病毒}, \zh{医院}, \zh{药}.
\item Huit collocations à connaître par cœur : \zh{得疟疾}, \zh{传播疟疾}, \zh{预防疟疾}, \zh{去医院看病}, \zh{睡在蚊帐里}, \zh{吃药}, \zh{治疗疟疾}, \zh{按时吃药}.
\item « Attraper / être atteint de » une maladie = \zh{得} + maladie ; « prendre » un médicament = \zh{吃药}.
\item \zh{预防} (avant la maladie) $\neq$ \zh{治疗} (après la maladie).
\item Distinction \zh{得} (personne malade) vs \zh{有} (existence de la maladie dans un lieu).
\item Précision scientifique : le paludisme est une maladie parasitaire (\zh{疟原虫}), pas virale (\zh{病毒}), mais \zh{传播} « transmettre » s'utilise pour les deux.
\end{itemize}
\end{aretenir}

Ce lexique est désormais votre matière première. Dans la section suivante, nous apprendrons à \emph{écrire correctement ces caractères} : ordre des traits, radicaux (clés) et pièges des caractères proches.

\coursec{Écriture correcte des caractères}

Dans la section précédente, nous avons découvert le lexique du paludisme : \zh{疟疾} (nüèjí, « paludisme »), \zh{蚊子} (wénzi, « moustique »), \zh{蚊帐} (wénzhàng, « moustiquaire »), \zh{预防} (yùfáng, « prévenir »), \zh{治疗} (zhìliáo, « soigner »). Mais reconnaître un caractère à la lecture ne suffit pas : en Terminale, vous devez être capable de l'\textbf{écrire} correctement, trait par trait, sans faute. Un trait manquant ou mal placé peut changer le sens, voire rendre le caractère illisible. Cette section vous donne une méthode rigoureuse pour maîtriser les caractères du thème.

\courssub{Observer avant d'écrire : la logique des caractères chinois}

Observez d'abord ces cinq caractères : \zh{病}、\zh{疾}、\zh{疟}、\zh{疼}、\zh{痛}. Que remarquez-vous ? Tous contiennent, en haut à gauche, le même élément graphique : \zh{疒}. Ce n'est pas un hasard.

\begin{propriete}[Le radical 疒, clé des maladies]
Les caractères liés aux maladies partagent presque tous le radical \cle{疒} (nè), appelé « clé de la maladie ». Il est lui-même issu d'un ancien pictogramme représentant un homme alité. Ainsi : \zh{病} bìng (maladie), \zh{疾} jí (maladie rapide), \zh{疟} nüè (paludisme), \zh{疼} téng (douloureux), \zh{痛} tòng (douleur). Quand vous voyez \zh{疒} dans un caractère inconnu, pensez immédiatement : « cela concerne probablement la santé ou la maladie ».
\end{propriete}

C'est la grande logique de l'écriture chinoise : un caractère composé réunit généralement un élément de \textbf{sens} (le radical ou clé) et un élément de \textbf{son} (la phonétique) — ou parfois deux éléments de sens. Par exemple, dans \zh{蚊} wén (moustique), \zh{虫} chóng indique le sens (« insecte ») et \zh{文} wén indique la prononciation. Retenir cette double information, c'est retenir le caractère deux fois plus facilement qu'en apprenant par cœur une suite de traits.

\begin{methode}[Mémoriser un caractère en trois étapes]
\begin{enumerate}
\item \textbf{Identifier le radical (sens)} : quel élément donne la catégorie de sens ? (疒 = maladie, 虫 = insecte, 巾 = tissu, 氵 = eau, 页 = tête/page, 阝 = mound/terre…)
\item \textbf{Identifier l'autre composant (son ou sens)} : quel élément donne le son (phonétique) ou précise le sens (composant sémantique) ? (文 → wén, 长 → zhǎng, 矢 → sens « flèche », 台 → tái, 方 → fāng, 予 → yǔ/yù…)
\item \textbf{Écrire 5 fois} le caractère dans des cases tianzige (田字格), en \textbf{comptant les traits à voix haute} et en respectant l'ordre : de haut en bas, de gauche à droite, l'horizontal avant le vertical, l'enveloppe avant le contenu (sauf pour les enveloppes fermées en bas).
\end{enumerate}
\end{methode}

\courssub{Le radical 疒 et ses caractères : 疟 et 疾}

\textbf{Ordre des traits du radical 疒 (5 traits).} On écrit d'abord le \textbf{point} en haut (\zh{丶}), puis le \textbf{trait horizontal} (\zh{一}), puis le grand \textbf{trait oblique} vers la gauche (\zh{丿}) qui forme l'enveloppe, puis un \textbf{point} à l'intérieur (en haut à gauche), et enfin un \textbf{trait horizontal} qui ferme l'enveloppe en bas. Le contenu du caractère vient \emph{ensuite}, sous cette enveloppe. Erreur classique du francophone : écrire le contenu d'abord puis « coiffer » le tout — le caractère devient déséquilibré.

\textbf{Le caractère 疟 nüè (8 traits).} Il se compose du radical \zh{疒} (5 traits) + \zh{丆} (3 traits : un trait horizontal, puis un trait vertical-oblique avec crochet \zh{㇇}, puis un trait horizontal). Ordre complet : 1. point, 2. horizontal, 3. oblique de l'enveloppe, 4. point intérieur, 5. horizontal de fermeture — puis sous l'enveloppe : 6. horizontal, 7. vertical-oblique avec crochet, 8. horizontal. \zh{疟} ne s'emploie pratiquement que dans le mot \zh{疟疾} nüèjí, « paludisme ».

\textbf{Le caractère 疾 jí (10 traits).} Radical \zh{疒} (5 traits) + \zh{矢} shǐ (5 traits), qui signifie « flèche ». L'image est parlante : une maladie qui frappe \emph{comme une flèche}, vite et fort — d'où le sens ancien de « maladie rapide, soudaine ». \zh{矢} est ici un \textbf{composant sémantique} (et non phonétique, car shǐ ≠ jí). On le retrouve dans \zh{疾病} jíbìng (maladie, terme général et soutenu) et bien sûr \zh{疟疾}. Ordre : les 5 traits de \zh{疒}, puis \zh{矢} : 6. petit oblique, 7. horizontal, 8. horizontal, 9. grand oblique, 10. point final (na) en bas à droite.

\begin{attention}[疒 : le point qui change tout]
Deux erreurs reviennent sans cesse chez les élèves francophones :
\begin{itemize}
\item \textbf{Oublier le point initial} de \zh{疒} : sans ce point, on obtient un élément qui ressemble au radical de la « falaise » \zh{厂} — le caractère est fautif.
\item \textbf{Écrire 治 avec un 氵 incomplet} : \zh{治} zhì (soigner, gouverner) porte le radical de l'\textbf{eau} \zh{氵}, qui compte \textbf{trois traits} (deux points puis un trait relevé \zh{㇀}), pas deux. \zh{治} n'a \emph{rien} à voir avec \zh{疒} : ne le confondez pas avec un caractère de maladie, c'est un verbe d'action.
\end{itemize}
\end{attention}

\courssub{Le moustique et la moustiquaire : 蚊 et 帐}

\textbf{Le caractère 蚊 wén (10 traits).} Structure gauche-droite : à gauche, le radical \zh{虫} chóng (6 traits), clé des insectes et petites bêtes ; à droite, la phonétique \zh{文} wén (4 traits : point, horizontal, oblique croisée, point final). La phonétique donne ici le son \emph{exact} : 文 wén → 蚊 wén. Cadeau pour la mémoire !

\textbf{Le caractère 帐 zhàng (7 traits).} Structure gauche-droite : à gauche, le radical \zh{巾} jīn (3 traits : oblique, relevé, oblique), clé du « tissu, morceau d'étoffe » ; à droite, la phonétique \zh{长} cháng/zhǎng (4 traits : oblique, horizontal, oblique, point final na). Le sens est limpide : une moustiquaire, c'est un \textbf{morceau de tissu} tendu. Le mot complet : \zh{蚊帐} wénzhàng, littéralement « tissu à moustiques ».

\begin{popfigure}
\begin{center}
\begin{tikzpicture}[node distance=0.35cm and 0.6cm]
\node[draw=popblue, fill=popblueL, rounded corners, minimum width=1.1cm, minimum height=1.1cm, font=\Large] (chong) {虫};
\node[right=of chong, font=\Large, text=popdark] (plus1) {$+$};
\node[right=of plus1, draw=poporange, fill=poporangeL, rounded corners, minimum width=1.1cm, minimum height=1.1cm, font=\Large] (wen) {文};
\node[right=of wen, font=\Large, text=popdark] (eq) {$=$};
\node[right=of eq, draw=popgreen, fill=popgreenL, rounded corners, minimum width=1.1cm, minimum height=1.1cm, font=\Large] (wenzi) {蚊};
\node[below=0.15cm of chong, align=center, font=\small, text=popblue] {sens\\ insecte};
\node[below=0.15cm of wen, align=center, font=\small, text=poporange] {son\\ wén};
\node[below=0.15cm of wenzi, align=center, font=\small, text=popgreen] {wén\\ moustique};
\node[below=1.6cm of chong, draw=popblue, fill=popblueL, rounded corners, minimum width=1.1cm, minimum height=1.1cm, font=\Large] (jin) {巾};
\node[right=of jin, font=\Large, text=popdark] (plus2) {$+$};
\node[right=of plus2, draw=poporange, fill=poporangeL, rounded corners, minimum width=1.1cm, minimum height=1.1cm, font=\Large] (chang) {长};
\node[right=of chang, font=\Large, text=popdark] (eq2) {$=$};
\node[right=of eq2, draw=popgreen, fill=popgreenL, rounded corners, minimum width=1.1cm, minimum height=1.1cm, font=\Large] (zhang) {帐};
\node[below=0.15cm of jin, align=center, font=\small, text=popblue] {sens\\ tissu};
\node[below=0.15cm of chang, align=center, font=\small, text=poporange] {son\\ zhǎng};
\node[below=0.15cm of zhang, align=center, font=\small, text=popgreen] {zhàng\\ moustiquaire};
\node[below=1.6cm of jin, draw=popblue, fill=popblueL, rounded corners, minimum width=1.1cm, minimum height=1.1cm, font=\Large] (ne) {疒};
\node[right=of ne, font=\Large, text=popdark] (plus3) {$+$};
\node[right=of plus3, draw=poporange, fill=poporangeL, rounded corners, minimum width=1.1cm, minimum height=1.1cm, font=\Large] (shi) {矢};
\node[right=of shi, font=\Large, text=popdark] (eq3) {$=$};
\node[right=of eq3, draw=popgreen, fill=popgreenL, rounded corners, minimum width=1.1cm, minimum height=1.1cm, font=\Large] (ji) {疾};
\node[below=0.15cm of ne, align=center, font=\small, text=popblue] {sens\\ maladie};
\node[below=0.15cm of shi, align=center, font=\small, text=poporange] {sens\\ flèche};
\node[below=0.15cm of ji, align=center, font=\small, text=popgreen] {jí\\ maladie rapide};
\end{tikzpicture}
\legende{Décomposition des caractères 蚊, 帐 et 疾 : élément de sens + élément de son (ou de sens).}
\end{center}
\end{popfigure}

\courssub{Tableau récapitulatif des caractères du thème}

\begin{center}
\begin{tabularx}{\linewidth}{|X|X|X|X|}
\hline
\rowcolor{popblueL} Caractère & Pinyin & Radical + traits & Traduction \\
\hline
疟 & nüè & 疒 (nè), 8 traits & paludisme (dans 疟疾) \\
\hline
疾 & jí & 疒 (nè), 10 traits & maladie rapide \\
\hline
病 & bìng & 疒 (nè), 10 traits & maladie, tomber malade \\
\hline
蚊 & wén & 虫 (chóng), 10 traits & moustique \\
\hline
帐 & zhàng & 巾 (jīn), 7 traits & tente, moustiquaire \\
\hline
预 & yù & 页 (yè), 10 traits & d'avance, pré- \\
\hline
防 & fáng & 阝 (fù, oreille gauche), 6 traits & prévenir, se protéger \\
\hline
治 & zhì & 氵 (shuǐ, eau), 8 traits & soigner, gouverner \\
\hline
\end{tabularx}
\end{center}

\courssub{Caractères proches à ne pas confondre}

Le chinois regorge de paires de caractères visuellement ou sémantiquement proches. Voici les quatre pièges de notre thème.

\textbf{1. 病 / 疾 : deux « maladies ».} Tous deux portent \zh{疒} et signifient « maladie », mais \zh{病} est le mot courant (\zh{我病了。} Wǒ bìng le. « Je suis tombé malade. »), tandis que \zh{疾} est plus soutenu et entre dans des composés : \zh{疾病} jíbìng (les maladies en général), \zh{疟疾} nüèjí. Ne dites pas \zh{我疾了} : c'est une faute.

\textbf{2. 蚊 / 虫 : moustique / insecte.} \zh{虫} chóng est le terme générique (« insecte, ver ») et sert de radical ; \zh{蚊} wén désigne précisément le moustique. Un moustique est un 虫, mais un 虫 n'est pas forcément un 蚊.

\textbf{3. 治 / 活 : soigner / vivre.} Tous deux commencent par \zh{氵} (l'eau), mais la partie droite diffère : \zh{治} = 氵 + 台, \zh{活} = 氵 + 舌. \zh{治疗} zhìliáo = soigner ; \zh{生活} shēnghuó = vivre, la vie. Un trait de travers et votre phrase sur le paludisme devient une phrase sur la vie quotidienne !

\textbf{4. 帐 / 账 : moustiquaire / compte.} Même phonétique \zh{长}, même son zhàng, mais radicaux différents : \zh{帐} avec \zh{巾} (tissu) → moustiquaire ; \zh{账} avec \zh{贝} (coquillage, ancienne monnaie) → compte, comptabilité. Le radical vous donne la clé : dormir sous une \zh{蚊帐}, faire ses \zh{账} à la banque.

\begin{exemplebox}[治疗 : deux synonymes réunis]
Le mot \zh{治疗} zhìliáo (« soigner, traitement ») illustre un procédé très fréquent du chinois moderne : la réunion de \textbf{deux caractères synonymes}. \zh{治} zhì signifie « gouverner, soigner » et \zh{疗} liáo (radical 疒 !) signifie « soigner ». Ensemble, ils forment un verbe dissyllabique clair et stable. Exemple : \zh{他在医院治疗疟疾。} \emph{Tā zài yīyuàn zhìliáo nüèjí.} « Il soigne son paludisme à l'hôpital. » Même logique dans \zh{疾病} jíbìng (疾 + 病 = maladie) et \zh{预防} yùfáng (prévenir + se protéger = prévention).
\end{exemplebox}

\begin{exoresolu}[Identifier radicaux et composants]
Énoncé. Pour chaque caractère, indiquez le radical (sens), l'autre composant (son ou sens) et le nombre total de traits : 蚊, 帐, 疾, 治.
\tcblower
Solution détaillée.
\begin{itemize}
\item \zh{蚊} wén : radical \zh{虫} (insecte, sens) + \zh{文} (phonétique wén) ; 10 traits.
\item \zh{帐} zhàng : radical \zh{巾} (tissu, sens) + \zh{长} (phonétique zhǎng) ; 7 traits.
\item \zh{疾} jí : radical \zh{疒} (maladie, sens) + \zh{矢} (flèche, composant sémantique) ; 10 traits.
\item \zh{治} zhì : radical \zh{氵} (eau, sens) + \zh{台} (phonétique tái) ; 8 traits. Attention : trois traits pour \zh{氵} (deux points + relevé), pas deux.
\end{itemize}
\end{exoresolu}

\begin{exercice}[Tracé en cases tianzige]
Sur votre cahier, tracez des cases tianzige (田字格) et écrivez \textbf{5 fois chacun} les caractères suivants, en comptant les traits à voix haute et en respectant l'ordre étudié : \zh{疟}、\zh{疾}、\zh{蚊}、\zh{帐}、\zh{预}、\zh{防}. Puis recopiez une fois les mots complets : \zh{疟疾}、\zh{蚊帐}、\zh{预防}、\zh{治疗}.
\end{exercice}

\begin{aretenir}[L'essentiel de la section]
\begin{itemize}
\item Les caractères de maladie portent le radical \zh{疒} : on écrit le point, l'horizontal, l'enveloppe oblique, le point intérieur, l'horizontal de fermeture ; le contenu vient ensuite.
\item Chaque caractère composé = radical (sens) + phonétique (son) \textbf{ou} composant sémantique (sens) : exploitez cette double piste pour mémoriser.
\item Paires à surveiller : 病/疾, 蚊/虫, 治/活, 帐/账.
\item On écrit 5 fois en comptant les traits : la main et la voix aident la mémoire.
\end{itemize}
\end{aretenir}

\coursec{Structure du texte et connecteurs}

Dans la section précédente, vous avez appris à écrire correctement les caractères du thème du paludisme : \zh{疟疾} (nüèjí, « paludisme »), \zh{蚊子} (wénzi, « moustique »), \zh{发烧} (fāshāo, « avoir de la fièvre »), \zh{预防} (yùfáng, « prévenir »). Savoir écrire les mots ne suffit pas : il faut maintenant les \cle{organiser} dans un texte cohérent. Un bon texte chinois, comme un bon texte français, suit un plan clair et relie ses phrases par des \cle{connecteurs} (连接词, liánjiēcí). C'est exactement ce que nous allons construire ici, étape par étape.

\courssub{Observer : un texte modèle complet}

Lisons d'abord le texte modèle que nous avons commencé à étudier dans cette leçon. Observez comment les phrases s'enchaînent et repérez les mots qui font le lien entre elles.

\begin{textebox}[Texte modèle : 预防疟疾 (Yùfáng nüèjí — Prévenir le paludisme)]
疟疾是一种很常见的病。\\
Nüèjí shì yì zhǒng hěn chángjiàn de bìng.\\
Le paludisme est une maladie très courante.\\[4pt]
在喀麦隆，因为蚊子很多，所以人们容易得疟疾。\\
Zài Kāmàilóng, yīnwèi wénzi hěn duō, suǒyǐ rénmen róngyì dé nüèjí.\\
Au Cameroun, comme il y a beaucoup de moustiques, les gens attrapent facilement le paludisme.\\[4pt]
蚊子把疟疾传给人。得了疟疾，人会发烧、头疼，还会觉得没力气。\\
Wénzi bǎ nüèjí chuán gěi rén. Dé le nüèjí, rén huì fāshāo, tóuténg, hái huì juéde méi lìqì.\\
Le moustique transmet le paludisme à l'homme. Quand on a le paludisme, on a de la fièvre, mal à la tête, et on se sent sans force.\\[4pt]
我们应该怎么预防疟疾呢？首先，我们应该睡在蚊帐里；然后，要保持环境干净，并且按时吃药。\\
Wǒmen yīnggāi zěnme yùfáng nüèjí ne? Shǒuxiān, wǒmen yīnggāi shuì zài wénzhàng lǐ; ránhòu, yào bǎochí huánjìng gānjìng, bìngqiě ànshí chī yào.\\
Comment devons-nous prévenir le paludisme ? D'abord, nous devons dormir sous une moustiquaire ; ensuite, il faut garder l'environnement propre et prendre les médicaments à temps.\\[4pt]
如果发烧了，就要马上去医院看医生。\\
Rúguǒ fāshāo le, jiù yào mǎshàng qù yīyuàn kàn yīshēng.\\
Si l'on a de la fièvre, il faut aller immédiatement à l'hôpital voir le médecin.\\[4pt]
我觉得预防疟疾很重要，因为健康比什么都重要。\\
Wǒ juéde yùfáng nüèjí hěn zhòngyào, yīnwèi jiànkāng bǐ shénme dōu zhòngyào.\\
Je pense que prévenir le paludisme est très important, car la santé est plus importante que tout.
\end{textebox}

\textbf{Questions de compréhension :}
\begin{enumerate}
  \item Pourquoi les gens attrapent-ils facilement le paludisme au Cameroun ?
  \item Quels sont les trois symptômes mentionnés dans le texte ?
  \item Quels sont les trois conseils de prévention donnés ?
  \item Que faut-il faire si l'on a de la fièvre ?
\end{enumerate}

\textbf{Glossaire du texte :}

\begin{center}
\begin{tabularx}{\linewidth}{|X|X|X|X|}
\hline
\rowcolor{popblueL} Caractères & Pinyin & Nature & Traduction \\
\hline
常见 & chángjiàn & adj. & courant, fréquent \\
\hline
容易 & róngyì & adj. & facile ; facilement \\
\hline
传 & chuán & v. & transmettre \\
\hline
力气 & lìqì & n. & force physique \\
\hline
蚊帐 & wénzhàng & n. & moustiquaire \\
\hline
保持 & bǎochí & v. & maintenir, garder \\
\hline
环境 & huánjìng & n. & environnement \\
\hline
干净 & gānjìng & adj. & propre \\
\hline
按时 & ànshí & adv. & à temps, ponctuellement \\
\hline
马上 & mǎshàng & adv. & immédiatement \\
\hline
\end{tabularx}
\end{center}

\courssub{Le plan type d'un texte sur une maladie}

En relisant le texte, vous remarquez qu'il suit toujours le même schéma, en cinq étapes. C'est le \cle{plan type} à réutiliser dans toutes vos productions écrites sur une maladie (paludisme, typhoïde, choléra…).

\begin{enumerate}
  \item \textbf{Présentation} : on définit la maladie. Modèle : \zh{疟疾是一种很常见的病。} (Nüèjí shì yì zhǒng hěn chángjiàn de bìng.) « Le paludisme est une maladie très courante. » Structure : X + 是 + 一种 + adjectif + 的 + 病.
  \item \textbf{Causes} : on explique l'origine. Modèle : \zh{蚊子把疟疾传给人。} (Wénzi bǎ nüèjí chuán gěi rén.) « Le moustique transmet le paludisme à l'homme. » Structure \cle{把} (bǎ) : Sujet + 把 + Objet + Verbe + 给 + Bénéficiaire.
  \item \textbf{Symptômes} : on décrit les effets. Modèle : \zh{人会发烧、头疼。} (Rén huì fāshāo, tóuténg.) « On a de la fièvre, mal à la tête. » Énumération avec la virgule chinoise \zh{、}.
  \item \textbf{Prévention et conseils} : on donne des solutions. Modèle : \zh{我们应该睡在蚊帐里。} (Wǒmen yīnggāi shuì zài wénzhàng lǐ.) « Nous devons dormir sous une moustiquaire. » Structure : Sujet + 应该/要 + Verbe + Complément de lieu.
  \item \textbf{Conclusion} : on donne son avis. Modèle : \zh{我觉得预防疟疾很重要。} (Wǒ juéde yùfáng nüèjí hěn zhòngyào.) « Je pense que prévenir le paludisme est très important. »
\end{enumerate}

\begin{popfigure}
\begin{tikzpicture}[
  node distance=0.55cm,
  etape/.style={rectangle, rounded corners=3pt, draw=popblue, fill=popblueL, text width=4.6cm, align=center, font=\small, inner sep=5pt},
  conn/.style={rectangle, rounded corners=3pt, draw=poporange, fill=poporangeL, text width=3.4cm, align=center, font=\small, inner sep=4pt},
  fleche/.style={-{Stealth[length=2.5mm]}, thick, popdark}
]
\node[etape, align=center] (intro){\textbf{1. Présentation}\\ 疟疾是一种…的病。};
\node[etape, below=of intro, align=center] (cause){\textbf{2. Causes}\\ 蚊子把疟疾传给人。};
\node[etape, below=of cause, align=center] (symp){\textbf{3. Symptômes}\\ 人会发烧、头疼。};
\node[etape, below=of symp, align=center] (prev){\textbf{4. Prévention}\\ 我们应该…};
\node[etape, below=of prev, align=center] (concl){\textbf{5. Conclusion}\\ 预防疟疾很重要。};
\node[conn, right=1.1cm of intro] (c1) {définition avec \zh{是}};
\node[conn, right=1.1cm of cause] (c2) {structure \zh{把} ; \zh{因为…所以…}};
\node[conn, right=1.1cm of symp] (c3) {énumération avec \zh{、}, \zh{还}};
\node[conn, right=1.1cm of prev, align=center] (c4){\zh{首先…然后…并且}\\ \zh{如果…就…}};
\node[conn, right=1.1cm of concl] (c5) {\zh{我觉得…很重要}};
\draw[fleche] (intro) -- (cause);
\draw[fleche] (cause) -- (symp);
\draw[fleche] (symp) -- (prev);
\draw[fleche] (prev) -- (concl);
\draw[fleche, poporange] (c1) -- (intro);
\draw[fleche, poporange] (c2) -- (cause);
\draw[fleche, poporange] (c3) -- (symp);
\draw[fleche, poporange] (c4) -- (prev);
\draw[fleche, poporange] (c5) -- (concl);
\end{tikzpicture}
\legende{Organigramme du plan type : à chaque étape correspond un connecteur ou une structure associée.}
\end{popfigure}

\courssub{Le connecteur de cause : 因为…所以…}

\begin{propriete}[因为…所以… : « parce que… donc »]
Structure : \textbf{因为 + cause，所以 + conséquence。}\\
En chinois, les deux moitiés \zh{因为} (yīnwèi, « parce que ») et \zh{所以} (suǒyǐ, « donc ») peuvent — et c'est même le cas le plus fréquent à l'écrit — apparaître \textbf{ensemble dans la même phrase}. C'est une différence majeure avec le français, où « parce que… donc » est une faute de syntaxe : on n'utilise qu'un seul des deux connecteurs. On peut aussi n'employer que \zh{因为} (en fin de phrase : « C'est parce que… ») ou que \zh{所以} (en début de phrase : « Donc… »), mais on ne les sépare jamais par un point final à l'intérieur d'une même énonciation.
\end{propriete}

\begin{exemplebox}[Exemples avec 因为…所以…]
\zh{因为蚊子很多，所以人们容易得疟疾。}\\
\emph{Yīnwèi wénzi hěn duō, suǒyǐ rénmen róngyì dé nüèjí.}\\
« Parce qu'il y a beaucoup de moustiques, les gens attrapent facilement le paludisme. »\\[4pt]
\zh{因为他没有睡在蚊帐里，所以生病了。}\\
\emph{Yīnwèi tā méiyǒu shuì zài wénzhàng lǐ, suǒyǐ shēngbìng le.}\\
« Comme il n'a pas dormi sous la moustiquaire, il est tombé malade. »\\[4pt]
\zh{因为昨天下雨了，所以蚊子更多了。}\\
\emph{Yīnwèi zuótiān xiàyǔ le, suǒyǐ wénzi gèng duō le.}\\
« Parce qu'il a plu hier, il y a encore plus de moustiques. »
\end{exemplebox}

\courssub{Le connecteur de condition : 如果…就…}

\begin{propriete}[如果…就… : « si… alors »]
Structure : \textbf{如果 + condition，(sujet) + 就 + conséquence。}\\
\zh{如果} (rúguǒ, « si ») introduit l'hypothèse ; \zh{就} (jiù, « alors, aussitôt ») introduit la conséquence et se place \textbf{devant le verbe} (ou l'auxiliaire modal), après le sujet de la deuxième proposition. Attention : \zh{就} n'est pas un mot de liaison autonome comme « alors » en français qui peut commencer une phrase ; il reste à l'intérieur de la proposition, collé au prédicat.
\end{propriete}

\begin{exemplebox}[Exemples avec 如果…就…]
\zh{如果发烧了，就要去医院。}\\
\emph{Rúguǒ fāshāo le, jiù yào qù yīyuàn.}\\
« Si on a de la fièvre, il faut aller à l'hôpital. » (Sujet omis, générique)\\[4pt]
\zh{如果你头疼，就应该告诉老师。}\\
\emph{Rúguǒ nǐ tóuténg, jiù yīnggāi gàosu lǎoshī.}\\
« Si tu as mal à la tête, tu dois le dire au professeur. »\\[4pt]
\zh{如果环境不干净，蚊子就会很多。}\\
\emph{Rúguǒ huánjìng bù gānjìng, wénzi jiù huì hěn duō.}\\
« Si l'environnement n'est pas propre, il y aura beaucoup de moustiques. »
\end{exemplebox}

\courssub{Les connecteurs de séquence : 首先…然后… (…最后…)}

Pour énumérer des conseils ou des étapes dans l'ordre chronologique ou logique, le chinois utilise \zh{首先} (shǒuxiān, « d'abord, en premier lieu ») puis \zh{然后} (ránhòu, « ensuite »). On peut prolonger la série avec \zh{最后} (zuìhòu, « enfin »).

\begin{exemplebox}[Exemples avec 首先…然后…]
\zh{首先我们应该睡在蚊帐里，然后保持环境干净。}\\
\emph{Shǒuxiān wǒmen yīnggāi shuì zài wénzhàng lǐ, ránhòu bǎochí huánjìng gānjìng.}\\
« D'abord, nous devons dormir sous la moustiquaire ; ensuite, garder l'environnement propre. »\\[4pt]
\zh{首先洗手，然后吃饭，最后刷牙。}\\
\emph{Shǒuxiān xǐ shǒu, ránhòu chīfàn, zuìhòu shuā yá.}\\
« D'abord se laver les mains, ensuite manger, enfin se brosser les dents. »
\end{exemplebox}

\courssub{Les connecteurs d'addition : 并且 et 还}

\begin{itemize}
  \item \zh{并且} (bìngqiě, « et de plus, et également ») relie deux propositions complètes (sujet + verbe) dans un style soutenu, très utile à l'écrit formel : \zh{我们应该睡在蚊帐里，并且按时吃药。} (Wǒmen yīnggāi shuì zài wénzhàng lǐ, bìngqiě ànshí chī yào.) « Nous devons dormir sous la moustiquaire et prendre les médicaments à temps. »
  \item \zh{还} (hái, « aussi, en plus ») est plus courant à l'oral et dans l'écrit courant. Il se place \textbf{devant le verbe} (ou l'auxiliaire) et suppose souvent un sujet commun : \zh{他会发烧，还会头疼。} (Tā huì fāshāo, hái huì tóuténg.) « Il a de la fièvre et aussi mal à la tête. »
\end{itemize}

\courssub{La phrase d'opinion : 我觉得…}

Pour conclure un texte argumentatif ou explicatif, le modèle d'opinion personnel est : \textbf{我觉得 + proposition。} (Wǒ juéde…, « je pense que… / je trouve que… »). On y ajoute souvent un adjectif évaluatif précédé de \zh{很} (hěn) ou un comparatif.

\zh{我觉得预防疟疾很重要。} \emph{Wǒ juéde yùfáng nüèjí hěn zhòngyào.} « Je pense que prévenir le paludisme est très important. »

\zh{我觉得健康比钱更重要。} \emph{Wǒ juéde jiànkāng bǐ qián gèng zhòngyào.} « Je pense que la santé est plus importante que l'argent. »

\courssub{Tableau récapitulatif des connecteurs}

\begin{center}
\begin{tabularx}{\linewidth}{|X|X|X|}
\hline
\rowcolor{popblueL} Structure & Exemple (caractères + pinyin) & Traduction \\
\hline
因为…所以… & 因为蚊子很多，所以人们容易得疟疾。Yīnwèi wénzi hěn duō, suǒyǐ rénmen róngyì dé nüèjí. & Parce qu'il y a beaucoup de moustiques, les gens attrapent facilement le paludisme. \\
\hline
如果…就… & 如果发烧了，就要去医院。Rúguǒ fāshāo le, jiù yào qù yīyuàn. & Si on a de la fièvre, il faut aller à l'hôpital. \\
\hline
首先…然后… & 首先睡在蚊帐里，然后保持环境干净。Shǒuxiān shuì zài wénzhàng lǐ, ránhòu bǎochí huánjìng gānjìng. & D'abord dormir sous la moustiquaire, ensuite garder l'environnement propre. \\
\hline
并且 & 我们应该睡在蚊帐里，并且按时吃药。Wǒmen yīnggāi shuì zài wénzhàng lǐ, bìngqiě ànshí chī yào. & Nous devons dormir sous la moustiquaire et prendre les médicaments à temps. \\
\hline
还 & 人会发烧，还会头疼。Rén huì fāshāo, hái huì tóuténg. & On a de la fièvre et aussi mal à la tête. \\
\hline
我觉得… & 我觉得预防疟疾很重要。Wǒ juéde yùfáng nüèjí hěn zhòngyào. & Je pense que prévenir le paludisme est très important. \\
\hline
\end{tabularx}
\end{center}

\courssub{Comparaison chinois / français : les pièges de la correspondance}

\begin{center}
\begin{tabularx}{\linewidth}{|X|X|X|}
\hline
\rowcolor{popblueL} Point de grammaire & En chinois & En français \\
\hline
Cause + conséquence & \zh{因为} et \zh{所以} coexistent dans la même phrase & « parce que… donc » est incorrect : un seul connecteur \\
\hline
Condition & \zh{如果…就…} : les deux éléments coexistent ; \zh{就} devant le verbe & « si… alors » : « alors » est facultatif, souvent en tête de proposition \\
\hline
« Il y a beaucoup de moustiques » & \zh{蚊子很多} (sujet + 很 + adjectif) & construction impersonnelle « il y a » \\
\hline
Énumération de noms/verbes & virgule chinoise \zh{、} entre items : \zh{发烧、头疼} & virgule simple « , » \\
\hline
Transmettre une maladie & structure \zh{把} : \zh{蚊子把疟疾传给人} & « Le moustique transmet le paludisme aux gens » \\
\hline
\end{tabularx}
\end{center}

\begin{attention}[Ne traduisez pas « il y a » par \zh{有} partout]
Les francophones traduisent mécaniquement « il y a beaucoup de moustiques » par \zh{有很多蚊子}. Cette phrase n'est pas fausse (existence), mais le chinois préfère de loin la structure \textbf{sujet + 很 + adjectif} pour décrire une quantité importante : \zh{蚊子很多。} (Wénzi hěn duō.), littéralement « les moustiques sont nombreux ». De même, « il y a beaucoup de malades » se dira plus naturellement \zh{病人很多} (bìngrén hěn duō). Réservez \zh{有} (yǒu) à la possession (\zh{我有两本书}, « j'ai deux livres ») et à l'existence localisée précise (\zh{桌子上有一杯水}, « il y a un verre d'eau sur la table »).
\end{attention}

\begin{attention}[Le paludisme n'est pas un virus]
Dans le texte modèle, on écrit \zh{蚊子把疟疾传给人} (le moustique transmet \textbf{le paludisme} à l'homme) et non \zh{传病毒} (transmet un virus). Le paludisme est une maladie parasitaire (parasite \zh{疟原虫}, nüèyuánchóng), pas virale. Évitez le terme \zh{病毒} (bìngdú) pour cette pathologie.
\end{attention}

\begin{methode}[Rédiger un texte sur une maladie en quatre étapes]
\begin{enumerate}
  \item \textbf{Écrire le plan en français} : présentation, causes, symptômes, prévention, conclusion. Une idée principale par étape, une phrase par idée.
  \item \textbf{Placer les connecteurs} : décider où insérer \zh{因为…所以…} (cause), \zh{如果…就…} (condition/conseil), \zh{首先…然后…} (séquence), \zh{并且}/\zh{还} (addition), et terminer par \zh{我觉得…}.
  \item \textbf{Traduire phrase par phrase} avec le lexique du thème (\zh{疟疾}, \zh{蚊子}, \zh{发烧}, \zh{蚊帐}, \zh{预防}…), en vérifiant l'ordre des mots chinois canonique : \cle{Sujet + Temps + Lieu + Verbe (+ Objet)}.
  \item \textbf{Relire et corriger} : vérifier les tons du pinyin, l'écriture des caractères (\zh{疟} vs \zh{疾}, \zh{帐} vs \zh{账}), la ponctuation chinoise （，。；？！）、 et la cohérence logique des connecteurs.
\end{enumerate}
\end{methode}

\begin{exoresolu}[Assembler les phrases avec le connecteur correct]
Énoncé : reliez les deux phrases avec le connecteur qui convient (\zh{因为…所以…} ou \zh{如果…就…}), écrivez en caractères, pinyin et traduisez.\\
1. 蚊子很多。/ 人们容易得疟疾。\\
2. 你发烧了。/ 你要去医院。
\tcblower
Solution détaillée :\\
1. Les deux phrases expriment une cause réelle (fait constaté : les moustiques sont nombreux) et sa conséquence directe : on utilise \zh{因为…所以…}.\\
\zh{因为蚊子很多，所以人们容易得疟疾。}\\
\emph{Yīnwèi wénzi hěn duō, suǒyǐ rénmen róngyì dé nüèjí.}\\
« Parce qu'il y a beaucoup de moustiques, les gens attrapent facilement le paludisme. »\\[4pt]
2. La première phrase pose une hypothèse (condition) : on utilise \zh{如果…就…}, avec \zh{就} placé devant le verbe \zh{要}. Le sujet \zh{你} peut être répété ou omis dans la seconde proposition.\\
\zh{如果你发烧了，(你)就要去医院。}\\
\emph{Rúguǒ nǐ fāshāo le, (nǐ) jiù yào qù yīyuàn.}\\
« Si tu as de la fièvre, tu dois aller à l'hôpital. »
\end{exoresolu}

\begin{exercice}[À vous de jouer]
1. Traduisez en chinois : « Parce que l'environnement n'est pas propre, il y a beaucoup de moustiques. »\\
2. Traduisez en chinois : « D'abord nous devons dormir sous la moustiquaire, ensuite prendre les médicaments à temps. »\\
3. Complétez avec \zh{并且} ou \zh{还} : 我们应该保持环境干净，\_\_\_ 睡在蚊帐里。\\
4. Rédigez une phrase de conclusion avec \zh{我觉得} sur la prévention du paludisme.
\end{exercice}

\begin{corrige}[Exercice « À vous de jouer »]
1. \zh{因为环境不干净，所以蚊子很多。} (Yīnwèi huánjìng bù gānjìng, suǒyǐ wénzi hěn duō.) — On préfère la structure sujet + 很 + adjectif (\zh{蚊子很多}) à \zh{有很多蚊子}.\\
2. \zh{首先我们应该睡在蚊帐里，然后按时吃药。} (Shǒuxiān wǒmen yīnggāi shuì zài wénzhàng lǐ, ránhòu ànshí chī yào.)\\
3. \zh{并且} : \zh{我们应该保持环境干净，并且睡在蚊帐里。} (Wǒmen yīnggāi bǎochí huánjìng gānjìng, bìngqiě shuì zài wénzhàng lǐ.) \zh{并且} relie deux propositions complètes. \zh{还} aurait exigé une structure à sujet unique : \zh{我们还要睡在蚊帐里。}\\
4. Exemple : \zh{我觉得预防疟疾很重要。} (Wǒ juéde yùfáng nüèjí hěn zhòngyào.) Ou : \zh{我觉得睡蚊帐最重要。} (Wǒ juéde shuì wénzhàng zuì zhòngyào.)
\end{corrige}

\begin{aretenir}[L'essentiel de la section]
\begin{itemize}
  \item Un texte sur une maladie suit cinq étapes : présentation (\zh{是一种…的病}) $\rightarrow$ causes (\zh{把} structure) $\rightarrow$ symptômes (\zh{、}, \zh{还}) $\rightarrow$ prévention (\zh{首先…然后…}, \zh{并且}) $\rightarrow$ conclusion (\zh{我觉得…很重要}).
  \item \zh{因为…所以…} : les deux connecteurs coexistent en chinois (cause réelle), contrairement au français.
  \item \zh{如果…就…} : \zh{就} se place obligatoirement devant le verbe/auxiliaire de la conséquence.
  \item \zh{首先…然后…(最后…)} pour ordonner les conseils ; \zh{并且} (propositions) / \zh{还} (verbe) pour ajouter une idée.
  \item Conclusion d'opinion : \zh{我觉得} + proposition + \zh{很重要} / \zh{比…更重要}.
  \item « Il y a beaucoup de X » $\rightarrow$ \zh{X 很多} (description), pas \zh{有很多 X} (existence).
  \item Précision scientifique : le paludisme se transmet (\zh{传}), ce n'est pas un \zh{病毒}.
\end{itemize}
\end{aretenir}

\coursec{Thème et version (traduction)}

Dans la section précédente, nous avons appris à organiser un texte sur le paludisme grâce aux connecteurs \zh{因为……所以……}, \zh{并且} ou \zh{应该}. Nous savons donc maintenant \emph{quoi dire} et \emph{dans quel ordre enchaîner les idées}. Il reste une dernière compétence, décisive au baccalauréat : la \cle{traduction}, c'est-à-dire le passage correct du français au chinois (le \cle{thème}) et du chinois au français (la \cle{version}). C'est ici que tout se joue : un bon texte écrit est d'abord une suite de phrases bien traduites.

\courssub{Le thème : du français vers le chinois}

\begin{definition}[Qu'est-ce que le thème ?]
Le \cle{thème} consiste à traduire une phrase ou un texte du \textbf{français vers le chinois}. C'est l'exercice le plus exigeant : il faut à la fois trouver le bon vocabulaire, choisir la bonne structure de phrase chinoise et écrire des caractères corrects. Au baccalauréat, chaque erreur de caractère, de ton ou d'ordre des mots est sanctionnée.
\end{definition}

\textbf{Observer avant de traduire.} Regardons d'abord comment une phrase française devient une phrase chinoise :

\begin{exemplebox}[Observer la transformation]
Français : « Le moustique transmet le paludisme. »\\
Chinois : \zh{蚊子传播疟疾。}\\
\emph{Wénzi chuánbō nüèjí.}\\
On remarque : sujet (\zh{蚊子}) + verbe (\zh{传播}) + objet (\zh{疟疾}) : le même ordre Sujet--Verbe--Objet qu'en français. Pas d'article, pas de conjugaison.
\end{exemplebox}

\begin{methode}[Méthode du thème en quatre étapes]
\begin{enumerate}
\item \textbf{Souligner les mots du glossaire} : repérer dans la phrase française les mots qui figurent dans le vocabulaire de la leçon (moustique = \zh{蚊子}, paludisme = \zh{疟疾}, moustiquaire = \zh{蚊帐}, fièvre = \zh{发烧}, prévenir = \zh{预防}, environnement = \zh{环境}).
\item \textbf{Choisir la structure} : identifier sujet, verbe, objet, compléments ; décider de la charpente chinoise (S + V + O, \zh{因为……所以……}, \zh{应该} + verbe, etc.).
\item \textbf{Écrire} : rédiger la phrase en caractères, sans se presser, trait par trait.
\item \textbf{Relire en vérifiant} : caractères (traits, radicaux), grammaire (place des compléments, particules \zh{了}, \zh{把}), et cohérence du sens.
\end{enumerate}
\end{methode}

\textbf{Thème guidé 1.} « Le moustique transmet le paludisme. »

Étape 1 : moustique = \zh{蚊子}, transmettre = \zh{传播}, paludisme = \zh{疟疾}. Étape 2 : structure simple S + V + O. Étape 3 : \zh{蚊子传播疟疾。} Étape 4 : vérification — \zh{蚊} a bien la clé de l'insecte \zh{虫} ; \zh{疟疾} a la clé de la maladie \zh{疒}. Traduction correcte.

\textbf{Thème guidé 2.} « Parce qu'il a attrapé le paludisme, il a de la fièvre. »

Étape 1 : attraper (une maladie) = \zh{得}, paludisme = \zh{疟疾}, avoir de la fièvre = \zh{发烧}. Étape 2 : la cause s'exprime avec \zh{因为……所以……} ; l'action accomplie demande la particule \zh{了}. Étape 3 : \zh{因为他得了疟疾，所以发烧了。} \emph{Yīnwèi tā dé le nüèjí, suǒyǐ fāshāo le.} Étape 4 : \zh{了} apparaît deux fois, une fois par action accomplie (attraper, puis avoir de la fièvre).

\begin{attention}[因为……所以…… : les deux moitiés ensemble]
En français on dit « parce qu'il est malade, il reste au lit » \emph{ou} « il est malade, donc il reste au lit », mais jamais « parce que... donc... ». En chinois, c'est l'inverse : on utilise volontiers \textbf{les deux} connecteurs ensemble : \zh{因为……所以……}. Ne traduisez donc pas « parce que » seul par \zh{因为} en oubliant le \zh{所以} dans la seconde partie.
\end{attention}

\textbf{Thème guidé 3.} « Nous devons dormir sous une moustiquaire. »

Étape 1 : devoir = \zh{应该}, dormir = \zh{睡}, moustiquaire = \zh{蚊帐}. Étape 2 : attention à la place du complément de lieu ! Avec le verbe \zh{睡}, le lieu introduit par \zh{在} se place \textbf{après} le verbe : \zh{睡在蚊帐里}. Étape 3 : \zh{我们应该睡在蚊帐里。} \emph{Wǒmen yīnggāi shuì zài wénzhàng lǐ.} Étape 4 : le verbe modal \zh{应该} précède bien le verbe principal.

\begin{attention}[在 + lieu : avant ou après le verbe ?]
Règle générale : \zh{在} + lieu se place \textbf{avant} le verbe principal : \zh{我在雅温得学习。} \emph{Wǒ zài Yǎwēndé xuéxí.} « J'étudie à Yaoundé. »\\
\textbf{Exception capitale} : avec les verbes de position ou de déplacement d'objet comme \zh{睡} (dormir), \zh{住} (habiter), \zh{放} (poser), \zh{坐} (s'asseoir), le complément de lieu se place \textbf{après} le verbe : \zh{睡在蚊帐里} (dormir sous une moustiquaire), \zh{住在杜阿拉} (habiter à Douala), \zh{把书放在桌子上} (poser le livre sur la table). Ne dites jamais \zh{在蚊帐里睡} dans ce contexte.
\end{attention}

\begin{popfigure}
\begin{tikzpicture}[font=\small]
\node[draw=popblue, fill=popblueL, rounded corners, align=center, minimum width=4.2cm] (fr) at (0,2.6) {\textbf{Français}\\dormir / sous une moustiquaire\\(verbe puis lieu)};
\node[draw=popgreen, fill=popgreenL, rounded corners, align=center, minimum width=4.2cm] (zh) at (8,2.6) {\textbf{Chinois}\\睡 / 在蚊帐里\\(verbe puis lieu avec 在)};
\draw[-{Stealth}, very thick, poporange] (fr) -- (zh) node[midway, above, text=popdark]{même ordre !};
\node[draw=popblue, fill=white, rounded corners, align=center] (fr2) at (0,0) {étudier / à Yaoundé};
\node[draw=popgreen, fill=white, rounded corners, align=center] (zh2) at (8,0) {在雅温得 / 学习\\(lieu AVANT le verbe)};
\draw[-{Stealth}, very thick, poppink] (fr2) -- (zh2) node[midway, above, text=popdark]{ordre inversé !};
\end{tikzpicture}
\legende{Ordre des mots : avec \zh{睡}, le lieu suit le verbe comme en français ; avec la plupart des autres verbes, \zh{在} + lieu précède le verbe.}
\end{popfigure}

\courssub{La version : du chinois vers le français}

\begin{definition}[Qu'est-ce que la version ?]
La \cle{version} consiste à traduire du \textbf{chinois vers le français}. Le danger n'est pas seulement de comprendre, mais de \emph{reformuler} en français naturel et correct : une traduction mot à mot est toujours fautive.
\end{definition}

\begin{methode}[Méthode de la version en trois étapes]
\begin{enumerate}
\item \textbf{Lire globalement} : lire toute la phrase (ou le texte) une première fois pour saisir le sens général, sans s'arrêter sur un caractère inconnu.
\item \textbf{Traduire par groupes de sens} : découper la phrase en blocs (sujet / verbe / compléments) et traduire bloc par bloc.
\item \textbf{Reformuler en français correct} : rétablir les articles, les conjugaisons, les accords ; vérifier que la phrase finale est du vrai français.
\end{enumerate}
\end{methode}

\textbf{Version guidée.} \zh{预防疟疾很重要。我们应该保持环境干净，并且睡在蚊帐里。}

Lecture globale : il s'agit de la prévention du paludisme. Groupes de sens : \zh{预防疟疾} (prévenir le paludisme) + \zh{很重要} (est très important) ; \zh{我们应该} (nous devons) + \zh{保持环境干净} (garder l'environnement propre) + \zh{并且} (et) + \zh{睡在蚊帐里} (dormir sous une moustiquaire). Reformulation : « Prévenir le paludisme est très important. Nous devons garder l'environnement propre et dormir sous une moustiquaire. »

\begin{exemplebox}[Exemples de version]
\zh{我们应该保持环境干净。} \emph{Wǒmen yīnggāi bǎochí huánjìng gānjìng.} « Nous devons garder l'environnement propre. »\\
\zh{蚊子把病菌传给人。} \emph{Wénzi bǎ bìngjūn chuán gěi rén.} « Le moustique transmet les agents pathogènes à l'homme. »\\
\zh{他昨天发烧了，现在在医院。} \emph{Tā zuótiān fāshāo le, xiànzài zài yīyuàn.} « Il a eu de la fièvre hier, il est maintenant à l'hôpital. »
\end{exemplebox}

\courssub{Les trois grands pièges de traduction}

\begin{propriete}[Piège 1 : l'ordre des mots]
Le complément de lieu avec \zh{在} se place avant le verbe, \textbf{sauf} avec \zh{睡}, \zh{住}, \zh{放}, \zh{坐} où il suit le verbe : \zh{睡在蚊帐里}. De même, le complément de temps se place avant le verbe : \zh{他昨天发烧了。} « Il a eu de la fièvre hier » (et non « il a eu hier de la fièvre » en chinois).
\end{propriete}

\begin{propriete}[Piège 2 : la particule 了]
\zh{了} marque l'action \textbf{accomplie} ; en français, il correspond souvent au passé composé : \zh{他得了疟疾。} « Il a attrapé le paludisme. » Attention : \zh{了} ne s'utilise pas avec les verbes d'état ni pour une habitude ; on ne dit pas \zh{我每天睡觉了} pour « je dors tous les jours ».
\end{propriete}

\begin{propriete}[Piège 3 : la construction 把]
La structure \zh{Sujet + 把 + objet + verbe + complément} met l'accent sur ce que le sujet \emph{fait subir} à l'objet : \zh{蚊子把病菌传给人。} « Le moustique transmet les agents pathogènes à l'homme. » En version, on la traduit par une phrase française ordinaire ; en thème, on l'utilise quand l'objet subit une transformation ou un déplacement.
\end{propriete}

\begin{center}
\begin{tabularx}{\linewidth}{|X|X|X|}
\hline
\rowcolor{popblueL} Structure & Exemple (caractères + pinyin) & Traduction \\
\hline
S + V + O & \zh{蚊子传播疟疾。} \emph{Wénzi chuánbō nüèjí.} & Le moustique transmet le paludisme. \\
\hline
因为……所以…… & \zh{因为他得了疟疾，所以发烧了。} \emph{Yīnwèi tā dé le nüèjí, suǒyǐ fāshāo le.} & Parce qu'il a attrapé le paludisme, il a de la fièvre. \\
\hline
应该 + V & \zh{我们应该睡在蚊帐里。} \emph{Wǒmen yīnggāi shuì zài wénzhàng lǐ.} & Nous devons dormir sous une moustiquaire. \\
\hline
把 + O + V & \zh{蚊子把病菌传给人。} \emph{Wénzi bǎ bìngjūn chuán gěi rén.} & Le moustique transmet les agents pathogènes à l'homme. \\
\hline
保持 + O + adj. & \zh{保持环境干净} \emph{bǎochí huánjìng gānjìng} & garder l'environnement propre \\
\hline
\end{tabularx}
\end{center}

\begin{exoresolu}[Thème : traduire en chinois]
Énoncé : Traduisez en chinois : « Parce que les moustiques transmettent le paludisme, nous devons dormir sous une moustiquaire. »
\tcblower
Solution détaillée : 1) Glossaire : moustique = \zh{蚊子}, transmettre = \zh{传播}, paludisme = \zh{疟疾}, devoir = \zh{应该}, dormir sous une moustiquaire = \zh{睡在蚊帐里}. 2) Structure : \zh{因为……所以……} avec deux propositions. 3) Rédaction : \zh{因为蚊子传播疟疾，所以我们应该睡在蚊帐里。} \emph{Yīnwèi wénzi chuánbō nüèjí, suǒyǐ wǒmen yīnggāi shuì zài wénzhàng lǐ.} 4) Vérification : \zh{应该} précède le verbe \zh{睡} ; le lieu \zh{在蚊帐里} suit \zh{睡} (exception) ; pas de \zh{了} car il s'agit d'un fait général, pas d'une action accomplie.
\end{exoresolu}

\begin{exercice}[À vous de traduire]
Thème : 1) « Il a attrapé le paludisme, il est à l'hôpital. » 2) « Nous devons garder l'environnement propre. »\\
Version : \zh{因为预防疟疾很重要，所以我们应该保持环境干净。}
\end{exercice}

\begin{corrige}[Exercice]
Thème 1 : \zh{他得了疟疾，他在医院。} \emph{Tā dé le nüèjí, tā zài yīyuàn.} — \zh{了} marque l'accompli ; \zh{在医院} : lieu avec le verbe \zh{在} (être à), sujet \zh{他} repris pour la clarté.\\
Thème 2 : \zh{我们应该保持环境干净。} \emph{Wǒmen yīnggāi bǎochí huánjìng gānjìng.}\\
Version : « Parce que prévenir le paludisme est très important, nous devons garder l'environnement propre. »
\end{corrige}

\begin{savaistu}[Le thème et la version au baccalauréat]
Au baccalauréat camerounais, les exercices de thème et de version portent souvent sur des sujets de \textbf{santé publique} comme le paludisme, l'hygiène ou la vaccination : ce sont des thèmes d'actualité au Cameroun, où le paludisme reste l'une des premières causes de consultation médicale. En Chine aussi, la lutte contre les maladies transmises par les moustiques a été une grande campagne nationale de santé publique (élimination du paludisme certifiée par l'OMS en 2021). Maîtriser le lexique de cette leçon, c'est donc se préparer à un sujet très probable à l'examen !
\end{savaistu}

\begin{aretenir}[L'essentiel de la section]
\begin{itemize}
\item Thème : glossaire $\rightarrow$ structure $\rightarrow$ écriture $\rightarrow$ relecture (caractères, tons, grammaire).
\item Version : lecture globale $\rightarrow$ groupes de sens $\rightarrow$ reformulation en vrai français.
\item Trois pièges : place de \zh{在} + lieu (après \zh{睡}, \zh{住}, \zh{放}), particule \zh{了} pour l'accompli, construction \zh{把}.
\item Connecteur de cause : \zh{因为……所以……} s'utilise en deux parties.
\end{itemize}
\end{aretenir}

\coursec{Production écrite guidée et évaluation}

Après avoir travaillé la traduction dans les deux sens (thème et version), vous disposez maintenant de tout le matériel nécessaire : le vocabulaire du paludisme, les caractères du thème, les connecteurs et les structures grammaticales. Il est temps de passer à l'étape finale et la plus importante : \cle{produire votre propre texte} en chinois. C'est exactement ce qui vous sera demandé à l'examen. Cette section vous guide pas à pas, du brainstorming à la relecture, puis vous montre comment votre copie sera évaluée.

\courssub{La tâche à accomplir}

\begin{definition}[Sujet d'expression écrite]
Rédiger un texte de \cle{5 à 8 phrases} en chinois simplifié, intitulé \zh{预防疟疾} (Yùfáng nüèji, « Prévenir le paludisme »). Le texte doit présenter le problème, ses causes, ses symptômes et des conseils de prévention, avec une conclusion personnelle.
\end{definition}

Ce type de sujet est classique : on vous donne un thème de santé publique et on attend un petit texte argumentatif simple. Votre objectif n'est pas d'impressionner le correcteur avec des phrases compliquées, mais de démontrer que vous maîtrisez le vocabulaire du thème, les caractères et les structures de base.

\begin{methode}[Réussir l'expression écrite au Bac]
\begin{enumerate}
\item \textbf{Un plan simple avant tout.} Notez en français ou en chinois les 4 ou 5 idées que vous allez développer : introduction, causes, symptômes, conseils, conclusion. Ne rédigez jamais sans plan.
\item \textbf{Des phrases courtes et correctes valent mieux que des phrases longues fautives.} Une phrase de 8 caractères sans faute rapporte des points ; une phrase de 20 caractères avec trois erreurs en perd. Au Bac, la \cle{correction} prime sur l'ambition.
\item \textbf{Réutilisez le texte modèle.} Les structures vues dans la leçon (\zh{因为…所以…}, \zh{应该}, \zh{为了…}) sont vos meilleures alliées : reprenez-les telles quelles.
\item \textbf{Comptez vos phrases.} 5 à 8 phrases : si vous en avez 4, ajoutez un conseil ; si vous en avez 10, coupez.
\item \textbf{Relisez 3 minutes} avec la grille de relecture (voir plus bas).
\end{enumerate}
\end{methode}

\courssub{Étape 1 : brainstormer le vocabulaire utile}

Avant d'écrire, faites remonter de votre mémoire tout le lexique de la leçon. Voici le glossaire de référence à mobiliser :

\begin{center}
\begin{tabularx}{\linewidth}{|X|X|X|X|}
\hline
\rowcolor{popblueL} Caractères & Pinyin & Nature & Traduction \\
\hline
疟疾 & nüèji & n. & le paludisme \\
\hline
预防 & yùfáng & v. & prévenir \\
\hline
蚊子 & wénzi & n. & le moustique \\
\hline
发烧 & fāshāo & v./n. & avoir de la fièvre ; la fièvre \\
\hline
头疼 & tóuténg & v. & avoir mal à la tête \\
\hline
病人 & bìngrén & n. & le malade, le patient \\
\hline
医院 & yīyuàn & n. & l'hôpital \\
\hline
药 & yào & n. & le médicament \\
\hline
蚊帐 & wénzhàng & n. & la moustiquaire \\
\hline
睡觉 & shuìjiào & v. & dormir \\
\hline
脏水 & zāngshuǐ & n. & l'eau sale, les eaux usées \\
\hline
干净 & gānjìng & adj. & propre \\
\hline
健康 & jiànkāng & n./adj. & la santé ; en bonne santé \\
\hline
战胜 & zhànshèng & v. & vaincre, triompher de \\
\hline
努力 & nǔlì & v./adj. & faire des efforts ; assidu \\
\hline
\end{tabularx}
\end{center}

\begin{popfigure}
\begin{tikzpicture}[mindmap, grow cyclic, every node/.style={concept, font=\small},
root concept/.append style={concept color=popblue, font=\bfseries},
level 1/.append style={level distance=3.2cm, sibling angle=72},
level 1 concept/.append style={concept color=poporange!70}]
\node[root concept, align=center]{预防疟疾\\ Yùfáng nüèji}
child { node[align=center]{原因\\ 蚊子、脏水} }
child { node[align=center]{症状\\ 发烧、头疼} }
child { node[align=center]{预防\\ 蚊帐、干净的水} }
child { node[align=center]{治疗\\ 医院、药} }
child { node[align=center]{结论\\ 努力、战胜} };
\end{tikzpicture}
\legende{Carte mentale du lexique : organiser le vocabulaire par familles avant de rédiger.}
\end{popfigure}

\courssub{Étape 2 : rédiger le plan}

Un bon texte de 5 à 8 phrases suit un plan en cinq parties, chaque partie pouvant tenir en une ou deux phrases :

\begin{center}
\begin{tabularx}{\linewidth}{|X|X|X|}
\hline
\rowcolor{popblueL} Partie & Contenu attendu & Structure utile \\
\hline
Introduction & Le paludisme est une maladie dangereuse (au Cameroun). & 疟疾是一种危险的病。 \\
\hline
Causes & Les moustiques, les eaux stagnantes. & 因为…所以… \\
\hline
Symptômes & Fièvre, maux de tête, fatigue. & 会发烧、头疼。 \\
\hline
Conseils & Moustiquaire, eau propre, hôpital. & 我们应该…；为了… \\
\hline
Conclusion & Avis personnel + appel à l'effort collectif. & 我觉得…；就能… \\
\hline
\end{tabularx}
\end{center}

\courssub{Étape 3 : rédiger avec le texte modèle et les connecteurs}

Appuyez-vous sur les connecteurs étudiés dans la leçon :

\begin{itemize}
\item \zh{因为…所以…} (yīnwèi… suǒyǐ…) : « parce que… donc… » — pour les causes ;
\item \zh{应该} (yīnggāi) : « devoir (conseil) » — placé \textbf{avant} le verbe, pour les recommandations ;
\item \zh{为了} (wèile) : « afin de » — pour introduire un but ;
\item \zh{如果…就…} (rúguǒ… jiù…) : « si… alors… » — pour l'hypothèse ;
\item \zh{我觉得} (wǒ juéde) : « je pense que » — pour la conclusion personnelle.
\end{itemize}

\begin{textebox}[Modèle de conclusion à réutiliser]
我觉得预防疟疾很重要。大家一起努力，就能战胜疟疾！\\
Wǒ juéde yùfáng nüèji hěn zhòngyào. Dàjiā yìqǐ nǔlì, jiù néng zhànshèng nüèji!\\
« Je pense que prévenir le paludisme est très important. En travaillant tous ensemble, nous pouvons vaincre le paludisme ! »
\end{textebox}

Observez la structure de cette conclusion : \cle{avis personnel} (\zh{我觉得} + phrase) puis \cle{appel collectif} (\zh{大家一起努力}) puis \cle{résultat} (\zh{就能} + verbe). C'est un schéma en trois temps que vous pouvez reproduire pour n'importe quel sujet de santé ou d'environnement.

\begin{savaistu}[Une formule valorisée]
\zh{战胜疟疾} (zhànshèng nüèji, « vaincre le paludisme ») est une formule très appréciée dans les conclusions de texte argumentatif. Le verbe \zh{战胜} (vaincre, triompher de) donne une fin dynamique et optimiste, très valorisée dans la tradition rédactionnelle chinoise, qui aime les conclusions tournées vers l'action collective. Vous pouvez l'adapter : \zh{战胜疾病} (vaincre la maladie), \zh{战胜困难} (surmonter les difficultés).
\end{savaistu}

\courssub{Étape 4 : la relecture avec la grille}

\begin{attention}[Chaque caractère compte]
Au Bac, chaque caractère mal écrit — trait manquant, radical confondu — coûte des points. Les confusions les plus fréquentes dans ce thème : \zh{疟} (paludisme) qui perd un trait, \zh{蚊} (moustique) dont le radical 虫 (insecte) est oublié ou remplacé, \zh{帐} (moustiquaire, radical 巾 « tissu ») confondu avec \zh{账} (compte, radical 贝 « argent »). \textbf{Réservez toujours 3 minutes à la relecture finale}, uniquement pour vérifier les caractères, un par un.
\end{attention}

La grille de relecture, dans l'ordre :

\begin{enumerate}
\item \textbf{Caractères} : chaque caractère est-il complet et correct ? (radicaux 疒, 虫, 巾 bien écrits)
\item \textbf{Connecteurs} : ai-je utilisé au moins deux connecteurs de la leçon ? Sont-ils bien placés ?
\item \textbf{Grammaire} : ordre des mots (sujet + verbe + complément), \zh{应该} placé avant le verbe, pas de \zh{是} devant un adjectif seul.
\item \textbf{Plan} : introduction, causes, symptômes, conseils, conclusion — les cinq parties sont-elles présentes ?
\end{enumerate}

\begin{exoresolu}[Corriger un texte d'élève]
Voici le texte d'un élève anonyme. Relevez les erreurs, puis comparez avec la correction collective.

\zh{疟疾是一种很危险的病，很多人得这个病。因为有很多蚊子，所以人生病。得了疟疾的人发烧和头疼。我们应该睡觉用蚊帐，还要喝干净的水。我觉得预防疟疾重要，大家一起努力，就能战胜疟疾！}

\textit{Nüèji shì yī zhǒng hěn wēixiǎn de bìng, hěn duō rén dé zhège bìng. Yīnwèi yǒu hěn duō wénzi, suǒyǐ rén shēngbìng. Déle nüèji de rén fāshāo hé tóuténg. Wǒmen yīnggāi shuìjiào yòng wénzhàng, hái yào hē gānjìng de shuǐ. Wǒ juéde yùfáng nüèji zhòngyào, dàjiā yìqǐ nǔlì, jiù néng zhànshèng nüèji!}
\tcblower
\textbf{Correction collective — erreurs relevées ensemble :}
\begin{enumerate}
\item \zh{很多人得这个病} : correct, mais \zh{得这个病} est un peu vague ; mieux : \zh{很多人得疟疾} (beaucoup de gens attrapent le paludisme).
\item \zh{所以人生病} : trop général et maladroit ; préciser : \zh{所以人们容易生病} (donc les gens tombent facilement malades).
\item \zh{发烧和头疼} : erreur classique du francophone ! \zh{和} (hé, « et ») ne relie pas deux verbes ou deux propositions, seulement des noms. Correction : \zh{会发烧、头疼} avec la virgule chinoise 、, ou \zh{会发烧，也会头疼}.
\item \zh{睡觉用蚊帐} : ordre des mots fautif. Le complément de moyen se place avant le verbe : \zh{应该用蚊帐睡觉} (dormir sous une moustiquaire).
\item \zh{我觉得预防疟疾重要} : il manque l'intensifieur ; un adjectif seul après \zh{觉得} sonne incomplet : \zh{我觉得预防疟疾很重要}.
\end{enumerate}
\textbf{Bilan} : le plan est respecté (5/5 parties) et la conclusion est bonne, mais 3 erreurs de grammaire et de structure coûteraient environ 4 points. Texte corrigé :

\zh{疟疾是一种很危险的病，很多人得疟疾。因为有很多蚊子，所以人们容易生病。得了疟疾的人会发烧、头疼。我们应该用蚊帐睡觉，还要喝干净的水。我觉得预防疟疾很重要，大家一起努力，就能战胜疟疾！}

\textit{Nüèji shì yī zhǒng hěn wēixiǎn de bìng, hěn duō rén dé nüèji. Yīnwèi yǒu hěn duō wénzi, suǒyǐ rénmen róngyì shēngbìng. Déle nüèji de rén huì fāshāo、tóuténg. Wǒmen yīnggāi yòng wénzhàng shuìjiào, hái yào hē gānjìng de shuǐ. Wǒ juéde yùfáng nüèji hěn zhòngyào, dàjiā yìqǐ nǔlì, jiù néng zhànshèng nüèji!}
\end{exoresolu}

\courssub{Grille d'évaluation type examen}

Votre production sera notée sur 20 selon cinq critères, exactement comme à l'examen :

\begin{popfigure}
\begin{tikzpicture}
\fill[popblueL] (0,0) rectangle (10.5,0.8);
\node[anchor=west, font=\bfseries] at (0.2,0.4) {Critère};
\node[anchor=east, font=\bfseries] at (10.3,0.4) {Barème};
\fill[popcream] (0,-0.8) rectangle (10.5,0);
\node[anchor=west] at (0.2,-0.4) {1. Exactitude des caractères (traits, radicaux)};
\node[anchor=east, font=\bfseries, text=poporange] at (10.3,-0.4) {/4};
\fill[white] (0,-1.6) rectangle (10.5,-0.8);
\node[anchor=west] at (0.2,-1.2) {2. Grammaire et structures (ordre des mots, 应该, 因为…所以…)};
\node[anchor=east, font=\bfseries, text=poporange] at (10.3,-1.2) {/4};
\fill[popcream] (0,-2.4) rectangle (10.5,-1.6);
\node[anchor=west] at (0.2,-2.0) {3. Vocabulaire du thème (lexique du paludisme)};
\node[anchor=east, font=\bfseries, text=poporange] at (10.3,-2.0) {/4};
\fill[white] (0,-3.2) rectangle (10.5,-2.4);
\node[anchor=west] at (0.2,-2.8) {4. Organisation et connecteurs (plan en 5 parties)};
\node[anchor=east, font=\bfseries, text=poporange] at (10.3,-2.8) {/4};
\fill[popcream] (0,-4.0) rectangle (10.5,-3.2);
\node[anchor=west] at (0.2,-3.6) {5. Présentation générale (propreté, ponctuation chinoise)};
\node[anchor=east, font=\bfseries, text=poporange] at (10.3,-3.6) {/4};
\fill[popblue!60] (0,-4.8) rectangle (10.5,-4.0);
\node[anchor=west, font=\bfseries] at (0.2,-4.4) {Total};
\node[anchor=east, font=\bfseries] at (10.3,-4.4) {/20};
\draw[popline, thick] (0,0.8) rectangle (10.5,-4.8);
\draw[popline] (0,0) -- (10.5,0); \draw[popline] (0,-0.8) -- (10.5,-0.8);
\draw[popline] (0,-1.6) -- (10.5,-1.6); \draw[popline] (0,-2.4) -- (10.5,-2.4);
\draw[popline] (0,-3.2) -- (10.5,-3.2); \draw[popline] (0,-4.0) -- (10.5,-4.0);
\end{tikzpicture}
\legende{Grille d'évaluation de l'expression écrite : cinq critères à 4 points, total sur 20.}
\end{popfigure}

\begin{aretenir}[Les réflexes du jour de l'examen]
\begin{itemize}
\item Je lis le sujet deux fois et je note mon plan en 5 parties.
\item Je pioche dans le glossaire du thème : au moins 6 mots du lexique du paludisme.
\item J'écris des phrases courtes avec les structures apprises.
\item J'utilise au moins deux connecteurs (\zh{因为…所以…}, \zh{为了}, \zh{如果…就…}).
\item Je termine par une conclusion en trois temps : \zh{我觉得…} + effort collectif + \zh{就能战胜…！}
\item Je garde 3 minutes pour relire chaque caractère.
\end{itemize}
\end{aretenir}

\begin{exercice}[À vous d'écrire]
Rédigez maintenant votre texte \zh{预防疟疾} en 5 à 8 phrases. Respectez le plan en cinq parties, utilisez au moins six mots du glossaire et deux connecteurs. Terminez par la conclusion modèle adaptée à votre style. Échangez ensuite votre copie avec un voisin et corrigez-vous mutuellement avec la grille sur 20.
\end{exercice}

\begin{corrige}[Exercice — éléments de correction]
Il n'existe pas une seule bonne copie, mais un bon texte doit contenir : (1) une introduction du type \zh{疟疾是一种很危险的病。} ; (2) une cause avec \zh{因为…所以…} ; (3) des symptômes avec \zh{会发烧、头疼} (sans \zh{和} entre les verbes) ; (4) deux conseils avec \zh{应该} et \zh{为了} ; (5) une conclusion avec \zh{我觉得…很重要} et \zh{就能战胜疟疾！}. Vérifiez chaque caractère du glossaire : 疟疾、蚊子、发烧、蚊帐、预防、战胜. Attribuez-vous une note sur 20 avec la grille, puis réécrivez les phrases fautives.
\end{corrige}

\newpage

\coursec{Activité d'intégration}

\courssub{Objectifs de l'activité}

À l'issue de cette activité d'intégration, l'élève sera capable de :

\begin{itemize}
\item produire un \cle{texte fonctionnel} en chinois (une affiche de sensibilisation) destiné à un public réel, en respectant les contraintes du genre (titre accrocheur, phrases courtes, conseils impératifs) ;
\item mobiliser le \cle{lexique du paludisme} étudié dans la leçon : 疟疾 (nüèji, paludisme), 蚊子 (wénzi, moustique), 发烧 (fāshāo, avoir de la fièvre), 蚊帐 (wénzhàng, moustiquaire), 预防 (yùfáng, prévenir), 医院 (yīyuàn, hôpital) ;
\item employer correctement au moins trois \cle{connecteurs} : 因为…所以… (parce que… donc…), 如果…就… (si… alors…), 首先…然后… (d'abord… ensuite…) ;
\item écrire sans faute les caractères clés du thème, en respectant l'ordre des traits et en distinguant les caractères proches (蚊 / 纹, 帐 / 账, 疾 / 病) ;
\item pratiquer la \cle{version inversée} : traduire sa propre production chinoise en français correct, ce qui consolide la compétence de traduction dans les deux sens ;
\item évaluer la production d'autrui à l'aide d'une grille, puis réviser la sienne.
\end{itemize}

\courssub{Situation}

Le centre de santé de votre quartier à Yaoundé (quartier Mvog-Ada) organise, pendant la saison des pluies, une campagne de sensibilisation contre le paludisme. Le médecin-chef, qui sait que votre lycée enseigne le chinois, demande aux élèves de Terminale A de concevoir des affiches \emph{bilingues chinois-français} : elles seront exposées dans la salle d'attente du centre et montrées à une délégation de médecins chinois en visite d'échange au Cameroun. Votre affiche doit informer, convaincre et donner des conseils pratiques.

Voici le message que le médecin-chef vous a remis, et qui servira de cahier des charges :

\begin{textebox}[Message du médecin-chef (cahier des charges)]
同学们，雨季到了，得疟疾的病人越来越多。请你们做一张海报：\\
Tóngxuémen, yǔjì dào le, dé nüèji de bìngrén yuèláiyuè duō. Qǐng nǐmen zuò yì zhāng hǎibào :\\
« Élèves, la saison des pluies arrive, les malades atteints du paludisme sont de plus en plus nombreux. Je vous demande de réaliser une affiche. »\\[4pt]
海报上要有：一个题目、五到八个句子、三条建议，还要有法文翻译。\\
Hǎibào shang yào yǒu : yí ge tímù, wǔ dào bā ge jùzi, sān tiáo jiànyì, hái yào yǒu Fǎwén fānyì.\\
« L'affiche doit comporter : un titre, cinq à huit phrases, trois conseils, et aussi une traduction française. »
\end{textebox}

\begin{attention}[Lecture fine du cahier des charges]
Notez l'emploi de \zh{得} (dé) dans \zh{得疟疾} : 得 signifie ici « contracter, attraper » une maladie (得 + nom de maladie). Notez aussi le classificateur \zh{条} (tiáo) employé avec 建议 « conseil » : 三条建议 « trois conseils ». Enfin, 张 (zhāng) est le classificateur des objets plats : 一张海报 « une affiche ».
\end{attention}

\courssub{Consignes}

Travail en groupes de 3 ou 4 élèves. Durée : 2 heures (préparation, rédaction, présentation).

\begin{enumerate}
\item \textbf{Compréhension et repérage (10 min).} Lisez le message du médecin-chef. Soulignez les mots du lexique du paludisme et identifiez les quatre éléments obligatoires de l'affiche (titre, nombre de phrases, conseils, traduction).
\item \textbf{Brainstorming lexical (10 min).} En groupe, listez sur le brouillon au moins dix mots de la leçon que vous pourrez utiliser : 疟疾, 蚊子, 发烧, 头疼, 蚊帐, 药, 医院, 医生, 预防, 健康, 睡觉, 水… Vérifiez chaque caractère dans le tableau de vocabulaire.
\item \textbf{Choix du titre (5 min).} Reprenez le titre modèle 预防疟疾，人人有责 (Yùfáng nüèji, rénrén yǒu zé — « Prévenir le paludisme, l'affaire de tous ») ou composez un titre équivalent de 4 à 8 caractères.
\item \textbf{Rédaction du corps (30 min).} Rédigez 5 à 8 phrases courtes comportant \emph{au minimum} : une phrase avec 因为…所以…, une phrase avec 如果…就…, une séquence avec 首先…然后…, et \textbf{trois conseils} formulés à l'impératif (用蚊帐睡觉。/ 发烧了，要马上去医院。…).
\item \textbf{Version inversée (20 min).} Traduisez votre affiche en français, en bas de page, sans dictionnaire : c'est l'exercice de version appliqué à votre propre texte.
\item \textbf{Mise en forme (15 min).} Disposez le titre en grand, les phrases au centre, les trois conseils sous forme de liste visible (vous pouvez numéroter 一、二、三). Relisez : tons du pinyin sur le brouillon, caractères sur l'affiche.
\item \textbf{Présentation et évaluation par les pairs (30 min).} Chaque groupe présente son affiche en lisant le texte chinois à voix haute. Les autres groupes remplissent la grille d'évaluation ci-dessous. Les affiches validées seront exposées dans la salle de classe.
\end{enumerate}

\begin{attention}[Grille d'évaluation par les pairs (sur 20)]
\begin{itemize}
\item Lexique du thème correctement employé (au moins 8 mots) : /4
\item Caractères écrits sans faute (traits, caractères proches distingués) : /4
\item Au moins 3 connecteurs corrects (因为…所以…, 如果…就…, 首先…然后…) : /4
\item 3 conseils clairs à l'impératif : /3
\item Traduction française fidèle et correcte : /3
\item Présentation orale et lisibilité de l'affiche : /2
\end{itemize}
\end{attention}

\courssub{Ressources à mobiliser}

\begin{center}
\begin{tabularx}{\linewidth}{|X|X|X|X|}
\hline
\rowcolor{popblueL} Caractères & Pinyin & Nature & Traduction \\
\hline
疟疾 & nüèji & n. & paludisme \\
\hline
预防 & yùfáng & v. & prévenir \\
\hline
蚊子 & wénzi & n. & moustique \\
\hline
蚊帐 & wénzhàng & n. & moustiquaire \\
\hline
发烧 & fāshāo & v. & avoir de la fièvre \\
\hline
头疼 & tóuténg & v. & avoir mal à la tête \\
\hline
病人 & bìngrén & n. & malade, patient \\
\hline
医院 & yīyuàn & n. & hôpital \\
\hline
医生 & yīshēng & n. & médecin \\
\hline
药 & yào & n. & médicament \\
\hline
健康 & jiànkāng & n./adj. & santé ; en bonne santé \\
\hline
保护 & bǎohù & v. & protéger \\
\hline
马上 & mǎshàng & adv. & immédiatement \\
\hline
积水 & jīshuǐ & n. & eau stagnante \\
\hline
周围 & zhōuwéi & n. & alentours, aux alentours de \\
\hline
脏 & zāng & adj. & sale \\
\hline
雨季 & yǔjì & n. & saison des pluies \\
\hline
海报 & hǎibào & n. & affiche, poster \\
\hline
建议 & jiànyì & n./v. & conseil ; conseiller \\
\hline
越来越 & yuèláiyuè & adv. & de plus en plus \\
\hline
\end{tabularx}
\end{center}

\begin{propriete}[Rappel des connecteurs exigés]
\begin{itemize}
\item \cle{因为 A，所以 B} : cause puis conséquence. \zh{因为雨季蚊子很多，所以病人越来越多。} (Yīnwèi yǔjì wénzi hěn duō, suǒyǐ bìngrén yuèláiyuè duō.) « Parce qu'à la saison des pluies les moustiques sont nombreux, les malades sont de plus en plus nombreux. » En chinois, on peut garder les deux termes dans la même phrase, contrairement au français qui choisit « parce que » \emph{ou} « donc ».
\item \cle{如果 A，就 B} : hypothèse puis conséquence. \zh{如果发烧了，就要马上去医院。} (Rúguǒ fāshāo le, jiù yào mǎshàng qù yīyuàn.) « Si tu as de la fièvre, il faut aller immédiatement à l'hôpital. » Le 就 (jiù) dans la seconde proposition est quasi obligatoire.
\item \cle{首先…，然后…} : ordre des actions. \zh{首先挂好蚊帐，然后睡觉。} (Shǒuxiān guà hǎo wénzhàng, ránhòu shuìjiào.) « D'abord installer la moustiquaire, ensuite dormir. »
\item Conseil à l'impératif : verbe seul en tête de phrase, éventuellement précédé de 要 (il faut) ou de 别 / 不要 (ne… pas) : \zh{不要喝脏水。} (Búyào hē zāng shuǐ.) « Ne buvez pas d'eau sale. »
\end{itemize}
\end{propriete}

\begin{attention}[Pièges fréquents]
Ne confondez pas 蚊 (wén, moustique — clé 虫 « insecte ») et 纹 (wén, motif, ride — clé 纟« soie ») ; ni 帐 (zhàng, moustiquaire, tente — clé 巾 « tissu ») et 账 (zhàng, compte, comptabilité — clé 贝 « argent »). Dans 疟疾, le caractère 疾 porte la clé 疒 « maladie », comme 病 : on écrit d'abord la clé maladie (le point 丶, le trait horizontal 一, le trait descendant à gauche 丿, puis le point et le trait montant à l'intérieur gauche), ensuite la partie intérieure. Attention aussi à 越来越 + adjectif : jamais de 很 après (on dit 越来越多, non pas \emph{越来越很多}).
\end{attention}

\courssub{Proposition de production et corrigé commenté}

\begin{exoresolu}[Affiche modèle du groupe « Santé pour tous »]
Énoncé : rédiger l'affiche complète (titre, 6 à 8 phrases, 3 conseils, traduction française).
\tcblower
\textbf{Production modèle (texte chinois de l'affiche) :}

预防疟疾，人人有责\\
Yùfáng nüèji, rénrén yǒu zé.\\[4pt]
雨季到了，蚊子越来越多。因为蚊子会带来疟疾，所以我们要小心。\\
Yǔjì dào le, wénzi yuèláiyuè duō. Yīnwèi wénzi huì dàilái nüèji, suǒyǐ wǒmen yào xiǎoxīn.\\
如果发烧、头疼，就要马上去医院，不要自己吃药。\\
Rúguǒ fāshāo, tóuténg, jiù yào mǎshàng qù yīyuàn, búyào zìjǐ chī yào.\\
首先，我们要挂好蚊帐；然后，要清除房子周围的积水。\\
Shǒuxiān, wǒmen yào guà hǎo wénzhàng ; ránhòu, yào qīngchú fángzi zhōuwéi de jīshuǐ.\\
预防疟疾，保护健康！\\
Yùfáng nüèji, bǎohù jiànkāng !\\[4pt]
\textbf{Trois conseils :} 一、用蚊帐睡觉。二、发烧了，马上去看医生。三、不要喝脏水。\\
Yī, yòng wénzhàng shuìjiào. Èr, fāshāo le, mǎshàng qù kàn yīshēng. Sān, búyào hē zāng shuǐ.\\[6pt]
\textbf{Traduction française (version inversée) :}

« Prévenir le paludisme, l'affaire de tous. La saison des pluies arrive, les moustiques sont de plus en plus nombreux. Parce que les moustiques peuvent transmettre le paludisme, nous devons être prudents. Si tu as de la fièvre et mal à la tête, il faut aller immédiatement à l'hôpital, ne prends pas de médicaments tout seul. D'abord, nous devons bien installer la moustiquaire ; ensuite, éliminer les eaux stagnantes autour de la maison. Prévenons le paludisme, protégeons la santé ! Conseils : 1. Dormez sous moustiquaire. 2. En cas de fièvre, consultez immédiatement un médecin. 3. Ne buvez pas d'eau sale. »\\[6pt]
\textbf{Commentaire des choix de langue :}
\begin{itemize}
\item \textbf{Titre} : reprise du modèle 预防疟疾，人人有责, structure en deux segments de quatre syllabes, rythme typique des slogans chinois.
\item \textbf{Connecteurs} : 因为…所以… relie cause (les moustiques transmettent) et conséquence (prudence) ; 如果…就… formule l'hypothèse médicale ; 首先…然后… organise les gestes de prévention. Les trois exigences sont remplies.
\item \textbf{Impératifs} : 要… / 不要… pour les conseils ; 别 aurait été accepté (别喝脏水。Bié hē zāng shuǐ. « Ne buvez pas d'eau sale. »).
\item \textbf{Lexique} : 12 mots du thème employés (疟疾, 蚊子, 发烧, 头疼, 医院, 药, 蚊帐, 医生, 健康, 预防, 积水, 保护).
\item \textbf{Précision lexicale} : on dit 清除积水 « éliminer les eaux stagnantes » (les gîtes larvaires des moustiques) ; on ne peut pas dire \emph{打扫水}, car 打扫 « nettoyer, balayer » s'emploie avec un lieu (打扫房间 « nettoyer la chambre »), pas avec un liquide.
\item \textbf{Caractères} : 蚊 bien écrit avec la clé 虫 ; 帐 avec la clé 巾 ; 疾 et 病 avec la clé 疒.
\item \textbf{Version} : la traduction française reformule naturellement (« transmettre » pour 带来, littéralement « apporter ») : une bonne version n'est pas mot à mot.
\end{itemize}
\end{exoresolu}

\courssub{Bilan de l'activité}

Cette activité a permis de réinvestir l'ensemble des savoir-faire de la leçon dans une tâche authentique : comprendre une consigne en chinois (le message du médecin-chef), mobiliser le lexique du paludisme, écrire les caractères du thème en évitant les confusions entre caractères proches, structurer un texte avec les connecteurs 因为…所以…, 如果…就… et 首先…然后…, et pratiquer la traduction dans le sens chinois → français (version inversée). La présentation orale et l'évaluation par les pairs ont en outre développé l'expression orale et l'esprit critique. Les affiches exposées dans la salle constituent désormais un référentiel permanent de vocabulaire et de structures, utile pour les révisions du baccalauréat. Pour aller plus loin, les élèves volontaires pourront adapter leur affiche en court message radiophonique en chinois pour la radio du quartier.

\newpage

\coursec{Méthodes et exercices résolus}

\coursec{Leçon 14 — Expression écrite : le paludisme}

\courssub{Texte modèle et découverte du thème}

\begin{textebox}[Modèle de texte : « Le paludisme au Cameroun »]
\zh{在喀麦隆，疟疾是一种很常见的病。因为蚊子把疟疾传给人，所以很多人发烧、头疼。首先，我们应该用蚊帐睡觉；然后，如果生病了，就要马上去医院看病。这样，我们可以预防疟疾。} \\
\emph{Zài Kāmàilóng, nüèjí shì yì zhǒng hěn chángjiàn de bìng. Yīnwèi wénzi bǎ nüèjí chuán gěi rén, suǒyǐ hěn duō rén fāshāo, tóuténg. Shǒuxiān, wǒmen yīnggāi yòng wénzhàng shuìjiào; ránhòu, rúguǒ shēngbìng le, jiù yào mǎshàng qù yīyuàn kànbìng. Zhèyàng, wǒmen kěyǐ yùfáng nüèjí.} \\
« Au Cameroun, le paludisme est une maladie très courante. Parce que les moustiques transmettent le paludisme aux gens, beaucoup de personnes ont de la fièvre et mal à la tête. D'abord, nous devons dormir sous une moustiquaire ; ensuite, si l'on tombe malade, il faut aller tout de suite à l'hôpital se faire soigner. Ainsi, nous pouvons prévenir le paludisme. »
\end{textebox}

\begin{aretenir}[Analyse rapide du modèle]
\begin{itemize}
\item \textbf{Structure en 3 parties} : définition (phr. 1) → cause/symptômes (phr. 2) → prévention/traitement (phr. 3-4) → conclusion (phr. 5).
\item \textbf{Connecteurs} : \zh{因为……所以……} (cause-conséquence), \zh{首先……然后……} (étape), \zh{如果……就……} (condition-résultat), \zh{这样} (conclusion).
\item \textbf{Grammaire clé} : structure \zh{把} (\zh{蚊子把疟疾传给人}), compléments de lieu/temps avant le verbe (\zh{用蚊帐睡觉}, \zh{去医院看病}), aspect accompli \zh{了} (\zh{生病了}).
\end{itemize}
\end{aretenir}

\courssub{Lexique du paludisme}

\begin{definition}[Vocabulaire de base (niveau HSK 2-3)]
Le tableau ci-dessous regroupe les mots indispensables pour parler du paludisme. Mémorisez le caractère, le pinyin (tons sur voyelles), la nature et le sens.
\end{definition}

\begin{center}
\begin{tabularx}{\linewidth}{|X|X|X|X|}
\hline
\rowcolor{popblueL} \textbf{Caractères} & \textbf{Pinyin} & \textbf{Nature} & \textbf{Traduction} \\ \hline
疟疾 & nüèjí & n. & paludisme \\ \hline
蚊子 & wénzi & n. & moustique \\ \hline
传染 & chuánrǎn & v. & contaminer, transmettre (maladie) \\ \hline
发烧 & fāshāo & v. & avoir de la fièvre \\ \hline
头疼 & tóuténg & v. & avoir mal à la tête \\ \hline
症状 & zhèngzhuàng & n. & symptôme \\ \hline
蚊帐 & wénzhàng & n. & moustiquaire \\ \hline
预防 & yùfáng & v. & prévenir \\ \hline
治疗 & zhìliáo & v. & soigner, traiter \\ \hline
药 & yào & n. & médicament, remède \\ \hline
医院 & yīyuàn & n. & hôpital \\ \hline
医生 & yīshēng & n. & médecin \\ \hline
病人 & bìngrén & n. & malade, patient \\ \hline
严重 & yánzhòng & adj. & grave, sérieux \\ \hline
常见 & chángjiàn & adj. & courant, fréquent \\ \hline
马上 & mǎshàng & adv. & tout de suite, immédiatement \\ \hline
应该 & yīnggāi & mv. & devoir (conseil moral) \\ \hline
如果……就…… & rúguǒ…jiù… & conj. & si… alors… \\ \hline
因为……所以…… & yīnwèi…suǒyǐ… & conj. & parce que… donc… \\ \hline
首先……然后…… & shǒuxiān…ránhòu… & adv. & d'abord… ensuite… \\ \hline
\end{tabularx}
\end{center}

\begin{exemplebox}[Mots composés et collocations utiles]
\begin{itemize}
\item \zh{得疟疾} (dé nüèjí) : attraper le paludisme (\zh{得} = contracter une maladie, 2e ton).
\item \zh{睡蚊帐} (shuì wénzhàng) / \zh{用蚊帐} (yòng wénzhàng) : dormir sous une moustiquaire / utiliser une moustiquaire.
\item \zh{去医院看病} (qù yīyuàn kànbìng) : aller à l'hôpital se faire soigner (structure V1 + Lieu + V2).
\item \zh{吃药} (chī yào) : prendre un médicament (litt. « manger le médicament »).
\item \zh{打扫卫生} (dǎsǎo wèishēng) : faire le ménage, assainir les lieux (prévention).
\end{itemize}
\end{exemplebox}

\courssub{Écriture correcte des caractères}

\begin{methode}[Maîtriser les traits, radicaux et caractères proches du thème « santé »]
\begin{enumerate}
\item \textbf{Identifier le radical (clé) avant d'écrire} : il donne l'indice sémantique et dicte l'ordre des traits. Pour ce thème, quatre radicaux sont essentiels :
\begin{itemize}
\item \zh{疒} (bìngzìtóu, « clé de la maladie ») : \zh{病} (bìng, maladie), \zh{疟} (nüè, paludisme), \zh{疾} (jí, maladie grave), \zh{疼} (téng, douleur).
\item \zh{虫} (chóng, « clé de l'insecte ») : \zh{蚊} (wén, moustique).
\item \zh{巾} (jīn, « clé du tissu ») : \zh{帐} (zhàng, moustiquaire, tente).
\item \zh{火} (huǒ, « clé du feu ») : \zh{烧} (shāo, brûler ; \zh{发烧} = fièvre, le feu monte).
\end{itemize}
\item \textbf{Respecter l'ordre des traits universel} :
\begin{itemize}
\item De haut en bas (\zh{上} → \zh{下}), de gauche à droite (\zh{左} → \zh{droite}).
\item Horizontal avant vertical (\zh{横} avant \zh{竖}).
\item Extérieur avant intérieur ; on « ferme » la boîte en dernier (ex. \zh{囗}, \zh{门}).
\item Pour les caractères à clé enveloppante (\zh{疒}, \zh{广}, \zh{厂}) : écrire la clé \textbf{en premier} (sauf le trait de fermeture si elle se referme en bas), puis l'intérieur, enfin le trait de fermeture.
\end{itemize}
\item \textbf{Distinguer les caractères proches (faux amis graphiques)} :
\begin{itemize}
\item \zh{蚊} (wén, moustique, radical \zh{虫}) ≠ \zh{纹} (wén, motif/ride, radical \zh{纟}).
\item \zh{帐} (zhàng, moustiquaire, radical \zh{巾} tissu) ≠ \zh{账} (zhàng, compte/facture, radical \zh{贝} argent/coquillage).
\item \zh{疟} (nüè, paludisme, radical \zh{疒}) ≠ \zh{虐} (nüè, maltraiter, radical \zh{虎} tigre — composant phonétique différent).
\item \zh{病} (bìng, maladie) vs \zh{痛} (tòng, douleur, radical \zh{疒} + \zh{甬}) : \zh{头疼} (tóuténg, mal de tête) s'écrit avec \zh{疼} (téng, lecture altérée), pas \zh{痛}.
\end{itemize}
\item \textbf{S'entraîner par paires} : écris chaque caractère 5 fois en énonçant le nom des traits à voix haute (ex. « point, horizontal, trait tombant… »).
\end{enumerate}
\end{methode}

\begin{exoresolu}[Analyse et écriture guidée des caractères]
\textbf{Énoncé.} 1) Donne le radical, son sens et le nombre de traits du caractère : \zh{病}, \zh{蚊}, \zh{烧}, \zh{帐}. 2) Décris l'ordre des traits de \zh{病}. 3) Choisis le bon caractère : \zh{我们应该用蚊(帐 / 账)}.
\tcblower
\textbf{1) Radicaux, sens et nombre de traits (caractères simplifiés) :}
\begin{itemize}
\item \zh{病} (bìng) : radical \zh{疒} (maladie, 5 traits), total 10 traits. Sens : « maladie, tomber malade ».
\item \zh{蚊} (wén) : radical \zh{虫} (insecte, 6 traits), total 10 traits. Sens : « moustique ».
\item \zh{烧} (shāo) : radical \zh{火} (feu, 4 traits), total 10 traits. Sens : « brûler, chauffer » ; \zh{发烧} = « avoir de la fièvre ».
\item \zh{帐} (zhàng) : radical \zh{巾} (tissu, 3 traits), total 11 traits. Sens : « moustiquaire, tente, rideau ».
\end{itemize}

\textbf{2) Ordre des traits de \zh{病} (10 traits) :}
\begin{enumerate}
\item \zh{丶} (point, en haut à gauche de la clé)
\item \zh{一} (horizontal, traverse la clé)
\item \zh{丿} (trait tombant gauche, bord gauche de la clé)
\item \zh{丶} (point, à l'intérieur de la clé, en haut)
\item \zh{㇏} (trait tombant droit / \zh{提} levant, ferme la clé \zh{疒} en bas à droite) — \emph{règle : clé enveloppante écrite en premier}
\item \zh{一} (horizontal, début du composant intérieur \zh{丙})
\item \zh{丨} (vertical, traverse le milieu)
\item \zh{一} (horizontal, milieu)
\item \zh{一} (horizontal, bas du composant \zh{丙})
\item \zh{丿} (trait tombant gauche, dernier trait de \zh{丙}) — \emph{note : en écriture courante, le dernier trait de \zh{丙} peut être un \zh{㇏} (trait tombant droit) selon les polices, mais l'ordre standard est \zh{丿} puis \zh{㇏} pour \zh{丙} isolé ; ici dans \zh{病}, on écrit souvent le \zh{丿} final de \zh{丙} comme un trait court.}
\end{enumerate}
\emph{Resumé pédagogique :} « Clé \zh{疒} d'abord (5 traits), puis l'intérieur \zh{丙} de haut en bas. »

\textbf{3) Choix correct :} \zh{我们应该用蚊帐。} (\emph{Wǒmen yīnggāi yòng wénzhàng.})
\textbf{Justification} : \zh{帐} (radical \zh{巾} « tissu ») désigne une tente, une moustiquaire, un rideau — objet en tissu. \zh{账} (radical \zh{贝} « coquillage/argent ») désigne un compte, une facture, une dette. Le radical guide le sens : tissu vs argent.

\textbf{Conclusion} : identifier le radical, c'est accéder au sens et éviter les fautes de caractères « homographes ».
\end{exoresolu}

\courssub{Structure du texte et connecteurs logiques}

\begin{propriete}[Les 5 connecteurs indispensables pour un texte expositif de santé]
\begin{center}
\begin{tabularx}{\linewidth}{|X|X|X|}
\hline
\rowcolor{popblueL} \textbf{Fonction} & \textbf{Connecteur chinois} & \textbf{Structure / Place} \\ \hline
Énumération / Étapes & \zh{首先……然后……最后……} (shǒuxiān… ránhòu… zuìhòu…) & En tête de phrase, suivi de \zh{,} \\ \hline
Cause → Conséquence & \zh{因为……, 所以……} (yīnwèi…, suǒyǐ…) & \zh{因为} en tête prop. 1 ; \zh{所以} en tête prop. 2 (souvent après le sujet) \\ \hline
Condition → Résultat & \zh{如果……, 就……} (rúguǒ…, jiù…) & \zh{如果} en tête prop. 1 ; \zh{juste avant le verbe} de la prop. 2 \\ \hline
Opposition / Concession & \zh{但是 / 可是} (dànshì / kěshì) & En tête de la 2e proposition (après le sujet éventuel) \\ \hline
Conclusion / Conséquence finale & \zh{这样 / 所以} (zhèyàng / suǒyǐ) & En tête de la phrase finale \\ \hline
\end{tabularx}
\end{center}
\end{propriete}

\begin{attention}[Pièges classiques des connecteurs (spécifiques francophones)]
\begin{enumerate}
\item \textbf{Oublier la 2e moitié de la paire} : En français, « parce que » suffit seul. En chinois, \zh{因为} appelle \zh{所以} (ou \zh{就} dans certains cas). \zh{因为蚊子很多，所以很多人生病。} (pas \zh{因为蚊子很多，很多人生病。}).
\item \textbf{Mal placer \zh{就}} : \zh{就} se place \textbf{immédiatement avant le verbe} (ou l'auxiliaire) de la conséquence. \checkmark \zh{如果发烧了，就要去医院。} \quad \cross \zh{如果发烧了，要就去医院。}
\item \textbf{Confondre \zh{然后} et \zh{就}} : \zh{然后} = « ensuite » (chronologie, nouvelle phrase). \zh{就} = « alors » (conséquence immédiate, même phrase conditionnelle).
\item \textbf{Ponctuation} : Utilisez la virgule chinoise \zh{，} entre les propositions d'une même phrase ; le point-virgule \zh{；} pour séparer deux phrases proches dans une énumération (ex. \zh{首先……；然后……}).
\end{enumerate}
\end{attention}

\begin{exoresolu}[Composition : compléter un texte à trous avec les connecteurs]
\textbf{Énoncé.} Complète avec : \zh{因为}, \zh{所以}, \zh{如果}, \zh{就}, \zh{首先}, \zh{但是}.

\zh{①\_\_\_\_\_\_蚊子很多，②\_\_\_\_\_\_很多人得疟疾。③\_\_\_\_\_\_，我们要打扫房子周围；④\_\_\_\_\_\_发烧了，⑤\_\_\_\_\_\_要去看医生。疟疾很危险，⑥\_\_\_\_\_\_可以预防。}
\tcblower
\textbf{Raisonnement logique phrase par phrase :}
\begin{itemize}
\item \textbf{①-②} : Cause (beaucoup de moustiques) → Conséquence (beaucoup de malades) → Paire \zh{因为……所以……}.
\item \textbf{③} : Début d'une énumération de gestes de prévention → \zh{首先} (d'abord).
\item \textbf{④-⑤} : Condition (si fièvre) → Résultat immédiat (aller chez le médecin) → Paire \zh{如果……就……} ; \zh{就} se place avant \zh{要}.
\item \textbf{⑥} : Opposition entre « dangereux » et « prévenable » → \zh{但是}.
\end{itemize}

\textbf{Texte complété :}
\zh{因为蚊子很多，所以很多人得疟疾。首先，我们要打扫房子周围；如果发烧了，就要去看医生。疟疾很危险，但是可以预防。} \\
\emph{Yīnwèi wénzi hěn duō, suǒyǐ hěn duō rén dé nüèjí. Shǒuxiān, wǒmen yào dǎsǎo fángzi zhōuwéi; rúguǒ fāshāo le, jiù yào qù kàn yīshēng. Nüèjí hěn wēixiǎn, dànshì kěyǐ yùfáng.}

« Parce qu'il y a beaucoup de moustiques, beaucoup de gens attrapent le paludisme. D'abord, nous devons nettoyer les alentours de la maison ; si l'on a de la fièvre, il faut aller voir le médecin. Le paludisme est dangereux, mais on peut le prévenir. »
\end{exoresolu}

\courssub{Thème et version (traduction)}

\begin{methode}[Réussir le thème : du français vers le chinois (5 étapes)]
\begin{enumerate}
\item \textbf{Analyser le squelette français} : Qui (Sujet) ? Fait quoi (Verbe) ? Où (Lieu) ? Quand (Temps) ? Ne traduisez jamais mot à mot.
\item \textbf{Appliquer l'ordre chinois canonique} : \textbf{Sujet + Temps + Lieu + Verbe + Complément d'objet}. Ex. : « Hier, il va à l'hôpital » → \zh{他昨天去医院。} (Lui + hier + hôpital + va).
\item \textbf{Traiter les connecteurs logiques} : « Parce que… donc… » → \zh{因为……所以……} (garder les deux) ; « Si… alors… » → \zh{如果……就……} (\zh{juste avant le verbe}).
\item \textbf{Choisir les mots-outils et classificateurs} :
\begin{itemize}
\item \zh{了} : action accomplie / changement d'état (après le verbe).
\item \zh{的} : possession / qualification (avant le nom).
\item \zh{得} : évaluation de manière / degré (après le verbe).
\item Classificateurs : \zh{一种病}, \zh{一个医生}, \zh{一条蚊帐}, \zh{一些药} (pluriel indéfini).
\end{itemize}
\item \textbf{Relire en chinois mental} : La phrase sonne-t-elle naturelle pour un Sinophone ? (Pas de calque français).
\end{enumerate}
\end{methode}

\begin{exoresolu}[Thème : trois phrases types Bac]
\textbf{Énoncé.} Traduis en chinois :
1. « Ma sœur a de la fièvre, elle doit aller à l'hôpital. »
2. « Parce qu'il y a beaucoup de moustiques, nous dormons sous une moustiquaire. »
3. « Le médecin lui a donné des médicaments. »
\tcblower
\textbf{1) Squelette :} Sœur + fièvre ; Elle + hôpital + aller (obligation).
\textbf{Ordre :} Sujet + Verbe (fièvre) ; Sujet + Lieu + Verbe.
\zh{我妹妹发烧，她应该去医院。} \\
\emph{Wǒ mèimei fāshāo, tā yīnggāi qù yīyuàn.} \\
\textit{Note :} \zh{发烧} est verbe intransitif « avoir de la fièvre » (pas \zh{有发烧}). \zh{应该} exprime le conseil moral.

\textbf{2) Squelette :} Cause (moustiques nombreux) → Conséquence (nous + moustiquaire + dormir).
\textbf{Structure :} \zh{因为……所以……} + \zh{用蚊帐} (utiliser moustiquaire) + \zh{睡觉}.
\zh{因为蚊子很多，所以我们用蚊帐睡觉。} \\
\emph{Yīnwèi wénzi hěn duō, suǒyǐ wǒmen yòng wénzhàng shuìjiào.} \\
\textit{Note :} Classificateur \zh{条} pour \zh{蚊帐} si quantifié (\zh{一条蚊帐}), mais ici générique sans classif.

\textbf{3) Squelette :} Médecin + donner + lui + médicaments (action passée/accomplie).
\textbf{Structure :} \zh{给} (donner à) + \zh{了} (accompli) + personne + objet.
\zh{医生给了他一些药。} \\
\emph{Yīshēng gěi le tā yìxiē yào.} \\
\textit{Attention :} \zh{了} se place \textbf{juste après le verbe \zh{给}}, avant le bénéficiaire \zh{他}. \cross \zh{医生给他一些药了} (mal placé pour une action ponctuelle précise).
\end{exoresolu}

\begin{methode}[Réussir la version : du chinois vers le français (4 étapes)]
\begin{enumerate}
\item \textbf{Segmenter la phrase chinoise} en blocs syntaxiques : [Sujet] [Temps/Lieu] [Verbe + \zh{了}/\zh{着}/\zh{过}] [Objet] / [Connecteur] [Proposition].
\item \textbf{Traduire chaque bloc littéralement} (sous la phrase).
\item \textbf{Réorganiser selon la syntaxe française} : Sujet + Verbe + Compléments (Lieu/Temps souvent \textbf{après} le verbe). Rendre les aspects (\zh{le} = passé composé / accompli ; \zh{guo} = expérience ; \zh{zhe} = durée).
\item \textbf{Soigner le français final} : Phrase correcte, naturelle, sans mot-à-mot. Vérifier les temps (imparfait/passé composé), les prépositions (à, dans, sous), les articles.
\end{enumerate}
\end{methode}

\begin{exoresolu}[Version : texte court sur la prévention]
\textbf{Énoncé.} Traduis en français :
\zh{如果得了疟疾，就应该马上去医院。医生会给病人一些药。预防疟疾很重要：大家都要用蚊帐。}
\tcblower
\textbf{Découpage et traduction bloc par bloc :}
\begin{itemize}
\item \zh{如果得了疟疾} = « Si (on) a attrapé le paludisme » (\zh{得} = contracter, \zh{le} = accompli).
\item \zh{就应该马上去医院} = « alors il faut / on doit aller immédiatement à l'hôpital » (\zh{就} = alors ; \zh{应该} = devoir conseil ; \zh{马上} = tout de suite).
\item \zh{医生会给病人一些药} = « Le médecin donnera des médicaments au malade » (\zh{会} = futur / volonté ; \zh{给} + personne + objet).
\item \zh{预防疟疾很重要} = « Prévenir le paludisme est très important » (Sujet = verbe + objet ; \zh{很重要} = très important).
\item \zh{大家都要用蚊帐} = « Tout le monde doit utiliser une moustiquaire » (\zh{大家} = tout le monde ; \zh{都} = tous ; \zh{要} = devoir/obligation).
\end{itemize}

\textbf{Traduction rédigée (niveau Bac) :}
« Si l'on attrape le paludisme, il faut aller immédiatement à l'hôpital. Le médecin donnera des médicaments au malade. Prévenir le paludisme est très important : tout le monde doit utiliser une moustiquaire. »
\end{exoresolu}

\courssub{Production écrite guidée et évaluation}

\begin{methode}[Rédiger un texte court sur le paludisme en 5 étapes (Méthode Bac)]
\begin{enumerate}
\item \textbf{Choisir un plan en 3 parties} (1 paragraphe = 1 idée) :
① \textbf{Définition / Constat} : C'est quoi ? Où ? Fréquence ? (\zh{疟疾是……在喀麦隆很常见。})
② \textbf{Cause \& Symptômes} : Moustiques, transmission, fièvre, maux de tête. (\zh{因为蚊子传染，所以发烧头疼。})
③ \textbf{Prévention \& Traitement} : Moustiquaire, assainissement, hôpital, médicaments. (\zh{首先用蚊帐，然后去医院。})
\item \textbf{Lister le lexique imposé / ciblé} (5-7 mots) : \zh{疟疾, 蚊子, 发烧, 头疼, 蚊帐, 医院, 预防, 药}.
\item \textbf{Rédiger des phrases courtes} (Sujet + Verbe + Complément). Une idée = une phrase. \textbf{Interdit :} phrases à rallonge à la française.
\item \textbf{Insérer 2 à 3 connecteurs} variés : \zh{因为……所以……}, \zh{如果……就……}, \zh{首先……然后……}, \zh{但是}, \zh{这样}.
\item \textbf{Relire (check-list 30 secondes)} :
\begin{itemize}
\item Tons du pinyin vérifiés mentalement.
\item Ordre Sujet - Temps - Lieu - Verbe respecté ?
\item \zh{了} bien placé après le verbe (action passée) ? Absent sur conseil/habitude ?
\item Classificateurs présents si nombre + nom ?
\item Caractères proches vérifiés (\zh{帐/账}, \zh{蚊/纹}, \zh{疟/虐}) ?
\end{itemize}
\end{enumerate}
\end{methode}

\begin{exoresolu}[Rédaction guidée : « Le paludisme au Cameroun » (Sujet type Bac)]
\textbf{Consigne.} En 5 à 6 phrases, écris un texte intitulé « Le paludisme au Cameroun ». Contraintes : utiliser \zh{疟疾, 蚊子, 发烧, 蚊帐} + au moins 2 connecteurs (\zh{因为……所以……}, \zh{如果……就……}, \zh{首先/然后}).
\tcblower
\textbf{Étape 1 — Plan \& Lexique.}
Plan : 1. Constat (Cameroun, fréquent). 2. Cause/Symptômes (moustiques → fièvre/maux tête). 3. Prévention (moustiquaire, hôpital). 4. Conclusion.
Lexique coché : \zh{疟疾, 蚊子, 发烧, 头疼, 蚊帐, 医院, 预防}.

\textbf{Étape 2 — Construction phrase par phrase (vérification ordre des mots).}
1. \zh{在喀麦隆，疟疾是一种很常见的病。} (Lieu en tête, \zh{是一种} + classif.)
2. \zh{因为蚊子把疟疾传给人，所以很多人发烧、头疼。} (\zh{把} structure, \zh{因为……所以……})
3. \zh{首先，我们应该用蚊帐睡觉。} (\zh{首先}, \zh{应该} + V + O)
4. \zh{然后，如果生病了，就要马上去医院看病。} (\zh{然后}, \zh{如果……就……}, \zh{了} sur \zh{生病}, \zh{去医院看病})
5. \zh{这样，我们就能预防疟疾。} (\zh{这样} conclusion, \zh{就能} capacité/résultat).

\textbf{Étape 3 — Texte final.}
\zh{在喀麦隆，疟疾是一种很常见的病。因为蚊子把疟疾传给人，所以很多人发烧、头疼。首先，我们应该用蚊帐睡觉；然后，如果生病了，就要马上去医院看病。这样，我们就能预防疟疾。} \\
\emph{Zài Kāmàilóng, nüèjí shì yì zhǒng hěn chángjiàn de bìng. Yīnwèi wénzi bǎ nüèjí chuán gěi rén, suǒyǐ hěn duō rén fāshāo, tóuténg. Shǒuxiān, wǒmen yīnggāi yòng wénzhàng shuìjiào; ránhòu, rúguǒ shēngbìng le, jiù yào mǎshàng qù yīyuàn kànbìng. Zhèyàng, wǒmen jiù néng yùfáng nüèjí.}

« Au Cameroun, le paludisme est une maladie très courante. Parce que les moustiques transmettent le paludisme aux gens, beaucoup de gens ont de la fièvre et mal à la tête. D'abord, nous devons dormir sous une moustiquaire ; ensuite, si l'on tombe malade, il faut aller tout de suite à l'hôpital se faire soigner. Ainsi, nous pourrons prévenir le paludisme. »

\textbf{Auto-évaluation :} 5 phrases, 4 mots imposés + \zh{头疼}, 3 paires de connecteurs, ordre des mots correct, \zh{了} bien placé, \zh{帐} (tissu) correct. \textbf{Note attendue : 18-20/20.}
\end{exoresolu}

\begin{exercice}[Entraînement personnel — Expression écrite]
\textbf{Sujet :} « Comment prévenir le paludisme chez moi ? » (5-6 phrases).
\textbf{Contraintes :}
\begin{enumerate}
\item Utiliser \textbf{au moins 3 connecteurs} différents parmi : \zh{首先/然后/最后}, \zh{因为……所以……}, \zh{如果……就……}, \zh{但是}, \zh{这样}.
\item Intégrer \textbf{au moins 5 mots} du tableau de vocabulaire (dont \zh{蚊帐}, \zh{打扫卫生}, \zh{预防}).
\item Respecter l'ordre : Sujet + Temps/Lieu + Verbe.
\item Vérifier \zh{了} (action passée) et classificateurs.
\end{enumerate}
\textit{Écrivez votre brouillon sur une feuille, puis recopiez la version finale en caractères + pinyin + traduction.}
\end{exercice}

\begin{corrige}[Exercice — Expression écrite]
\textbf{Proposition de correction (une possibilité parmi d'autres correctes) :}
\zh{在我们家，预防疟疾很重要。因为蚊子喜欢脏水，所以首先我们要打扫房子周围的积水。然后，大家晚上都要睡蚊帐。如果有人发烧了，就要马上去医院吃药。最后，这样我们就能预防疟疾，身体健康。} \\
\emph{Zài wǒmen jiā, yùfáng nüèjí hěn zhòngyào. Yīnwèi wénzi xǐhuan zāng shuǐ, suǒyǐ shǒuxiān wǒmen yào dǎsǎo fángzi zhōuwéi de jīshuǐ. Ránhòu, dàjiā wǎnshang dōu yào shuì wénzhàng. Rúguǒ yǒurén fāshāo le, jiù yào mǎshàng qù yīyuàn chī yào. Zuìhòu, zhèyàng wǒmen jiù néng yùfáng nüèjí, shēntǐ jiànkāng.}

« Chez nous, prévenir le paludisme est très important. Parce que les moustiques aiment l'eau sale, d'abord nous devons nettoyer l'eau stagnante autour de la maison. Ensuite, tout le monde doit dormir sous une moustiquaire le soir. Si quelqu'un a de la fièvre, il faut aller tout de suite à l'hôpital prendre des médicaments. Enfin, ainsi nous pourrons prévenir le paludisme et rester en bonne santé. »

\textbf{Points de vigilance pour l'auto-correction :}
\begin{itemize}
\item \zh{积水} (jīshuǐ, eau stagnante) — vocabulaire bonus pertinent.
\item \zh{睡蚊帐} (shuì wénzhàng) — collocation correcte (pas \zh{用蚊帐睡觉} seulement, les deux sont justes).
\item \zh{吃药} (chī yào) — prendre des médicaments (courant à l'oral).
\item \zh{身体健康} (shēntǐ jiànkāng) — formule de conclusion standard.
\end{itemize}
\end{corrige}

\courssub{Élément socioculturel : Cameroun et Chine face au paludisme}

\begin{savaistu}[Le paludisme : un combat partagé, des stratégies différentes]
\begin{itemize}
\item \textbf{Cameroun :} Zone à transmission stable et intense (plasmodium \emph{falciparum}). Paludisme = 1ère cause de consultation (env. 30-40\% selon OMS). Stratégies nationales : Moustiquaires imprégnées à longue durée d'action (MILDA) distribuées gratuitement (campagnes « MILDA pour tous »), Traitement préventif intermittent chez la femme enceinte (TPIp), Pulvérisation intra-domiciliaire (PID) dans certaines zones, Prise en charge gratuite des < 5 ans. Défis : résistance aux insecticides, accès aux soins en zone rurale, automédication.
\item \textbf{Chine :} \textbf{Certifiée « exempte de paludisme » par l'OMS en juin 2021} (après 70 ans de lutte). Dernier cas autochtone en 2016 (province du Yunnan). Stratégie historique : « 1-3-7 » (déclaration sous 1 jour, investigation sous 3 jours, riposte sous 7 jours), utilisation massive de l'\textbf{artémisinine} (découverte par Tu Youyou, prix Nobel 2015, inspirée de la pharmacopée traditionnelle \zh{青蒿素} qīnghāosù), moustiquaires, gestion de l'environnement.
\item \textbf{Coopération :} Équipes médicales chinoises au Cameroun (hôpital de l'amitié sino-camerounais à Yaoundé, hôpital de district de Bafia, etc.), dons de médicaments (artémisinine), formation de personnel, centres de diagnostic de référence. Partage d'expérience « 1-3-7 » adapté au contexte camerounais.
\end{itemize}
\textbf{Piste de débat oral} : « Pourquoi la Chine a-t-elle réussi à éliminer le paludisme alors que le Cameroun reste endémique ? » (Facteurs : géographie, climat, système de santé, densité de population, financement, recherche locale).
\end{savaistu}

\courssub{Activité orale : Exposé et débat}

\begin{experience}[Débat structuré : « Stratégies de prévention au lycée »]
\textbf{Objectif :} Produire un discours suivi (2-3 min) en chinois, utiliser le vocabulaire et les connecteurs de la leçon.
\textbf{Déroulement (30 min) :}
\begin{enumerate}
\item \textbf{Préparation (10 min) :} Par groupes de 3. Chaque groupe prépare un mini-exposé (5 phrases) sur : « Comment notre lycée peut lutter contre les moustiques ? ». Mots imposés : \zh{校园} (xiàoyuán, campus), \zh{积水} (jīshuǐ), \zh{垃圾} (lājī, déchets), \zh{喷雾} (pēnwù, pulvériser), \zh{宣传} (xuānchuán, sensibiliser). Connecteurs obligatoires : \zh{首先}, \zh{因为……所以……}, \zh{如果……就……}.
\item \textbf{Présentation (10 min) :} Un élève par groupe lit/présente. Les autres notent : vocabulaire utilisé, connecteurs, prononciation (tons), fluidité.
\item \textbf{Débat / Vote (10 min) :} La classe vote pour le plan le plus réaliste. Le prof corrige 2-3 erreurs récurrentes au tableau (ex. place de \zh{了}, \zh{帐/账}, ordre temps/lieu).
\end{enumerate}
\textbf{Évaluation formative} : Grille 4 critères (Lexique / Grammaire-Connecteurs / Prononciation-Tons / Cohérence). Note sur 20 intégrée à la moyenne continue.
\end{experience}

\begin{attention}[Les 5 erreurs qui coûtent des points au Bac — Récapitulatif final]
\begin{enumerate}
\item \textbf{Ordre des mots calqué sur le français.} \cross \zh{他去医院昨天。} \checkmark \zh{他昨天去医院。} (Temps/Lieu \textbf{avant} le verbe). \cross \zh{我睡觉用蚊帐。} \checkmark \zh{我用蚊帐睡觉。} (Instrument \textbf{avant} l'action).
\item \textbf{Paire de connecteurs incomplète.} \cross \zh{因为蚊子很多，很多人生病。} \checkmark \zh{因为蚊子很多，\textbf{所以}很多人生病。} \cross \zh{如果发烧，要去医院。} \checkmark \zh{如果发烧\textbf{了}，\textbf{就}要去医院。}
\item \textbf{Confusion caractères proches (radical ignoré).} \cross \zh{蚊账} (moustique-compte), \zh{纹帐} (motif-tissu). \checkmark \zh{蚊帐} (insecte + tissu). \textbf{Réflexe :} Regarde le radical à gauche ou en haut !
\item \textbf{Mauvais placement de \zh{了}.} \cross \zh{医生给他一些药了。} (dans phrase narrative ponctuelle). \checkmark \zh{医生\textbf{给了}他一些药。} (\zh{了} collé au verbe). \cross \zh{我们应该用蚊帐了。} (conseil/habitude → pas de \zh{了}).
\item \textbf{Classificateur absent ou \zh{个} universel.} \cross \zh{一病}, \zh{一蚊帐}, \zh{三医生}. \checkmark \zh{一种病}, \zh{一条蚊帐}, \zh{三个医生} / \zh{三位医生}.
\end{enumerate}
\end{attention}

\newpage

\coursec{Exercices}

\courssub{Je vérifie mes connaissances}

\begin{exercice}[QCM : le paludisme]
Entoure la bonne réponse.
\begin{enumerate}
\item \zh{疟疾} (nüèji) signifie :
\begin{itemize}
\item A. le rhume \quad B. le paludisme \quad C. la fièvre typhoïde \quad D. la grippe
\end{itemize}
\item Le paludisme est transmis par :
\begin{itemize}
\item A. \zh{蚊子} (wénzi, le moustique) \quad B. \zh{水} (shuǐ, l'eau) \quad C. \zh{空气} (kōngqì, l'air) \quad D. \zh{食物} (shíwù, la nourriture)
\end{itemize}
\item \zh{发烧} (fāshāo) signifie :
\begin{itemize}
\item A. avoir mal à la tête \quad B. avoir de la fièvre \quad C. tousser \quad D. être fatigué
\end{itemize}
\item Pour se protéger du paludisme, on dort sous :
\begin{itemize}
\item A. \zh{被子} (bèizi, la couverture) \quad B. \zh{蚊帐} (wénzhàng, la moustiquaire) \quad C. \zh{床} (chuáng, le lit) \quad D. \zh{衣服} (yīfu, les vêtements)
\end{itemize}
\item \zh{预防比治疗更重要。} (Yùfáng bǐ zhìliáo gèng zhòngyào.) signifie :
\begin{itemize}
\item A. Le traitement est plus important que la prévention.
\item B. La prévention est plus importante que le traitement.
\item C. La prévention et le traitement sont identiques.
\item D. Le traitement est impossible.
\end{itemize}
\end{enumerate}
\end{exercice}

\begin{exercice}[Vrai ou faux ? Justifie ta réponse]
Réponds par VRAI ou FAUX, puis justifie en une phrase en français.
\begin{enumerate}
\item \zh{得了疟疾的人会发烧。} (Déle nüèji de rén huì fāshāo.) — Les malades du paludisme ont de la fièvre.
\item \zh{蚊帐} se prononce \emph{wénzhàng} et signifie « moustique ».
\item Dans \zh{我们应该保持环境干净。}, \zh{应该} (yīnggāi) exprime l'obligation ou le conseil.
\item Le caractère \zh{病} contient le radical de la maladie \zh{疒}.
\item \zh{雨季} (yǔjì) désigne la saison sèche.
\end{enumerate}
\end{exercice}

\begin{exercice}[Vocabulaire : complète le tableau]
Recopie et complète le tableau avec la nature grammaticale et la traduction française.
\begin{center}
\begin{tabularx}{\linewidth}{|X|X|X|X|}
\hline
\rowcolor{popblueL} Caractères & Pinyin & Nature & Traduction \\
\hline
疟疾 & nüèji & \ldots & \ldots \\
\hline
发烧 & fāshāo & \ldots & \ldots \\
\hline
蚊帐 & wénzhàng & \ldots & \ldots \\
\hline
预防 & yùfáng & \ldots & \ldots \\
\hline
治疗 & zhìliáo & \ldots & \ldots \\
\hline
医院 & yīyuàn & \ldots & \ldots \\
\hline
\end{tabularx}
\end{center}
\end{exercice}

\begin{exercice}[Associe les caractères au pinyin]
Relie chaque mot chinois à son pinyin.
\begin{center}
\begin{tabularx}{\linewidth}{|X|X|}
\hline
\rowcolor{popblueL} Caractères & Pinyin \\
\hline
1. 蚊子 & a. bìngdú \\
\hline
2. 病毒 & b. wénzi \\
\hline
3. 头疼 & c. gānjìng \\
\hline
4. 干净 & d. tóuténg \\
\hline
5. 危险 & e. wēixiǎn \\
\hline
\end{tabularx}
\end{center}
\end{exercice}

\courssub{Je m'entraîne}

\begin{exercice}[Traduis en chinois (thème)]
Traduis les phrases suivantes en chinois (caractères + pinyin).
\begin{enumerate}
\item Le moustique transmet le paludisme.
\item Pendant la saison des pluies, il y a beaucoup de moustiques.
\item Si tu as de la fièvre, il faut aller à l'hôpital.
\item Nous devons dormir sous la moustiquaire.
\item Il faut garder l'environnement propre.
\end{enumerate}
\end{exercice}

\begin{exercice}[Traduis en français (version)]
Traduis les phrases suivantes en français.
\begin{enumerate}
\item \zh{疟疾是一种很危险的病。} (Nüèji shì yī zhǒng hěn wēixiǎn de bìng.)
\item \zh{得了疟疾的人会发烧、头疼。} (Déle nüèji de rén huì fāshāo、tóuténg.)
\item \zh{我们应该睡在蚊帐里。} (Wǒmen yīnggāi shuì zài wénzhàng lǐ.)
\item \zh{预防疟疾比治疗疟疾更重要。} (Yùfáng nüèji bǐ zhìliáo nüèji gèng zhòngyào.)
\end{enumerate}
\end{exercice}

\begin{exercice}[Texte à trous : mots-outils et vocabulaire]
Complète le texte avec les mots de la liste : \zh{因为、所以、如果、就、比、应该、了、的}.
\begin{textebox}[Texte à trous]
\zh{\ldots 雨季蚊子很多，\ldots 很多人得疟疾。\ldots 你发烧，\ldots 要马上去医院。我们\ldots 睡在蚊帐里。预防\ldots 治疗更重要。得\ldots 疟疾\ldots 人要多喝水、多休息。}\\[0.3em]
\emph{... yǔjì wénzi hěn duō, ... hěn duō rén dé nüèji. ... nǐ fāshāo, ... yào mǎshàng qù yīyuàn. Wǒmen ... shuì zài wénzhàng lǐ. Yùfáng ... zhìliáo gèng zhòngyào. Dé ... nüèji ... rén yào duō hē shuǐ、duō xiūxi.}
\end{textebox}
\end{exercice}

\begin{exercice}[Le bon caractère : 病/疾, 治/活, 帐/账]
Complète chaque phrase avec le bon caractère parmi les deux proposés, puis traduis la phrase en français.
\begin{enumerate}
\item 疟\zh{…}是一种危险的病。(病 / 疾)
\item 得了疟疾要马上\zh{…}疗。(治 / 活)
\item 我们睡在蚊\zh{…}里。(帐 / 账)
\item 在喀麦隆，生\zh{…}很愉快。(活 / 治)
\item 生\zh{…}的人要去看医生。(病 / 疾)
\end{enumerate}
\end{exercice}

\begin{exercice}[Remets les phrases en ordre]
Les phrases du texte modèle ont été mélangées. Remets-les dans le bon ordre (note les lettres), puis justifie ton ordre en citant les connecteurs utilisés.
\begin{enumerate}
\item[A.] \zh{如果发烧、头疼，就要马上去医院。}
\item[B.] \zh{因为在雨季蚊子很多，所以很多人得疟疾。}
\item[C.] \zh{总之，预防比治疗更重要。}
\item[D.] \zh{疟疾是喀麦隆很常见的一种病。}
\item[E.] \zh{为了预防疟疾，我们应该睡在蚊帐里，保持环境干净。}
\end{enumerate}
\end{exercice}

\courssub{Je me prépare au Bac}

\begin{exercice}[Compréhension de texte (type Bac)]
Lis le texte suivant, puis réponds aux questions en français.

\begin{textebox}[Le paludisme à Yaoundé]
\zh{在喀麦隆，疟疾是一种很常见的病。雨季的时候，因为雨水多，蚊子也越来越多，所以很多人容易得疟疾。得了疟疾的人会发烧、头疼，有时候还很累。}\\[0.3em]
\emph{Zài Kāmàilóng, nüèji shì yī zhǒng hěn chángjiàn de bìng. Yǔjì de shíhou, yīnwèi yǔshuǐ duō, wénzi yě yuèláiyuè duō, suǒyǐ hěn duō rén róngyì dé nüèji. Déle nüèji de rén huì fāshāo、tóuténg, yǒu shíhou hái hěn lèi.}\\[0.3em]
\zh{为了预防疟疾，我们应该睡在蚊帐里，还应该保持环境干净。如果发烧了，就要马上去医院看医生。大家都说：预防比治疗更重要！}\\[0.3em]
\emph{Wèile yùfáng nüèji, wǒmen yīnggāi shuì zài wénzhàng lǐ, hái yīnggāi bǎochí huánjìng gānjìng. Rúguǒ fāshāo le, jiù yào mǎshàng qù yīyuàn kàn yīshēng. Dàjiā dōu shuō: yùfáng bǐ zhìliáo gèng zhòngyào!}\\[0.3em]
« Au Cameroun, le paludisme est une maladie très courante. Pendant la saison des pluies, comme il y a beaucoup d'eau de pluie, les moustiques sont de plus en plus nombreux, donc beaucoup de gens attrapent facilement le paludisme. Les malades ont de la fièvre, mal à la tête, et sont parfois très fatigués. Pour prévenir le paludisme, nous devons dormir sous la moustiquaire et garder l'environnement propre. Si l'on a de la fièvre, il faut aller tout de suite à l'hôpital voir le médecin. Tout le monde le dit : la prévention est plus importante que le traitement ! »
\end{textebox}

\textbf{Questions de compréhension}
\begin{enumerate}
\item Pourquoi y a-t-il plus de moustiques pendant la saison des pluies ?
\item Cite deux symptômes du paludisme mentionnés dans le texte.
\item Quels sont les deux moyens de prévention proposés ?
\item Que faut-il faire si l'on a de la fièvre ?
\end{enumerate}

\textbf{Questions de langue}
\begin{enumerate}
\item Releve dans le texte la structure de comparaison avec \zh{比} et recopie-la avec son pinyin.
\item Explique le rôle du connecteur \zh{为了} dans la troisième phrase.
\item Dans \zh{蚊子也越来越多}, quelle expression exprime la progression ? Traduis-la en français.
\end{enumerate}
\end{exercice}

\begin{exercice}[Thème et version (type Bac)]
\textbf{A. Thème.} Traduis en chinois (caractères + pinyin) :
\begin{enumerate}
\item « Parce que les moustiques sont nombreux pendant la saison des pluies, beaucoup de gens attrapent le paludisme. »
\item « Si tu as de la fièvre, il faut aller à l'hôpital. »
\end{enumerate}

\textbf{B. Version.} Traduis en français le texte suivant :
\begin{textebox}[Version]
\zh{疟疾是一种很危险的病。得了疟疾的人会发烧、头疼。我们应该睡在蚊帐里，保持环境干净。预防疟疾比治疗疟疾更重要。}\\[0.3em]
\emph{Nüèji shì yī zhǒng hěn wēixiǎn de bìng. Déle nüèji de rén huì fāshāo、tóuténg. Wǒmen yīnggāi shuì zài wénzhàng lǐ, bǎochí huánjìng gānjìng. Yùfáng nüèji bǐ zhìliáo nüèji gèng zhòngyào.}
\end{textebox}
\end{exercice}

\begin{exercice}[Expression écrite et dictée (type Bac)]
\textbf{A. Dictée de caractères.} Ton professeur dicte les huit mots suivants. Écris-les en caractères chinois, en respectant l'ordre des traits et les radicaux (la notation porte sur l'exactitude des traits et des radicaux) :
\zh{疟疾、蚊子、蚊帐、发烧、预防、治疗、医院、病毒}.

\textbf{B. Rédaction.} Sujet : \zh{预防疟疾} (« Prévenir le paludisme »).

Rédige en chinois un texte de \textbf{6 à 8 phrases} (environ 80 à 120 caractères) sur la prévention du paludisme au Cameroun. Tu devras :
\begin{itemize}
\item utiliser \textbf{au moins 3 connecteurs} parmi : \zh{因为…所以…、如果…就…、为了、还、总之} ;
\item employer \textbf{au moins 6 mots du glossaire} de la leçon (paludisme, moustique, moustiquaire, fièvre, prévention, traitement, hôpital, environnement, propre, dangereux…) ;
\item soigner l'écriture des caractères (traits, radicaux) et la ponctuation chinoise.
\end{itemize}
Ton texte sera noté sur 20 selon la grille : exactitude des caractères (4 pts), grammaire et structures (6 pts), vocabulaire du thème (4 pts), connecteurs et cohérence (4 pts), présentation et ponctuation (2 pts).
\end{exercice}



\newpage

\coursec{Corrigés des exercices}

\begin{corrige}[Exercice 1 — QCM : le paludisme]
\begin{enumerate}
\item \textbf{B.} \zh{疟疾} (nüèjí) = le paludisme. Attention au premier caractère \zh{疟} : radical de la maladie \zh{疒} (nè) + composant phonétique \zh{虐} (nüè).
\item \textbf{A.} Le paludisme est transmis par le moustique \zh{蚊子} (wénzi), et non par l'eau, l'air ou la nourriture.
\item \textbf{B.} \zh{发烧} (fāshāo) = avoir de la fièvre. « Avoir mal à la tête » se dit \zh{头疼} (tóuténg).
\item \textbf{B.} On dort sous la moustiquaire \zh{蚊帐} (wénzhàng). Ne pas confondre avec \zh{被子} (bèizi, la couverture).
\item \textbf{B.} Structure \cle{A + 比 + B + 更 + adjectif} : « A est plus ... que B ». Ici : la prévention est plus importante que le traitement.
\end{enumerate}
\end{corrige}

\begin{corrige}[Exercice 2 — Vrai ou faux ?]
\begin{enumerate}
\item \textbf{VRAI.} La fièvre est le symptôme principal du paludisme ; \zh{会} (huì) exprime ici la probabilité, la conséquence naturelle (ex. : \zh{得了疟疾会发烧}).
\item \textbf{FAUX.} La prononciation \emph{wénzhàng} est correcte, mais \zh{蚊帐} signifie « moustiquaire » ; « moustique » se dit \zh{蚊子} (wénzi).
\item \textbf{VRAI.} \zh{应该} (yīnggāi) est un verbe modal exprimant le devoir, le conseil : « il faut, on devrait ».
\item \textbf{VRAI.} \zh{病} = radical \zh{疒} (nè, la maladie) + composant phonétique \zh{丙} (bǐng). Ce radical se retrouve dans \zh{疾} (jí), \zh{疼} (téng), \zh{疟} (nüè).
\item \textbf{FAUX.} \zh{雨季} (yǔjì) = la saison des pluies (\zh{雨} = pluie). La saison sèche se dit \zh{旱季} (hànjì).
\end{enumerate}
\end{corrige}

\begin{corrige}[Exercice 3 — Vocabulaire : tableau complété]
\begin{center}
\begin{tabularx}{\linewidth}{|X|X|X|X|}
\hline
\rowcolor{popblueL} Caractères & Pinyin & Nature & Traduction \\
\hline
疟疾 & nüèjí & nom & le paludisme \\
\hline
发烧 & fāshāo & verbe & avoir de la fièvre \\
\hline
蚊帐 & wénzhàng & nom & la moustiquaire \\
\hline
预防 & yùfáng & verbe / nom & prévenir ; la prévention \\
\hline
治疗 & zhìliáo & verbe / nom & soigner, traiter ; le traitement \\
\hline
医院 & yīyuàn & nom & l'hôpital \\
\hline
\end{tabularx}
\end{center}
\end{corrige}

\begin{corrige}[Exercice 4 — Associe les caractères au pinyin]
1 -- b (\zh{蚊子} wénzi, le moustique) ; 2 -- a (\zh{病毒} bìngdú, le virus) ; 3 -- d (\zh{头疼} tóuténg, avoir mal à la tête) ; 4 -- c (\zh{干净} gānjìng, propre) ; 5 -- e (\zh{危险} wēixiǎn, dangereux).
\end{corrige}

\begin{corrige}[Exercice 5 — Traduis en chinois (thème)]
\begin{enumerate}
\item \zh{蚊子传播疟疾。} \emph{Wénzi chuánbō nüèjí.} — \zh{传播} (chuánbō) = transmettre, propager.
\item \zh{雨季的时候，蚊子很多。} \emph{Yǔjì de shíhou, wénzi hěn duō.} — « Il y a beaucoup de » se rend simplement par \zh{很多} après le sujet.
\item \zh{如果你发烧，就要去医院。} \emph{Rúguǒ nǐ fāshāo, jiù yào qù yīyuàn.} — Structure \cle{如果…就…} ; \zh{要} exprime ici la nécessité.
\item \zh{我们应该睡在蚊帐里。} \emph{Wǒmen yīnggāi shuì zài wénzhàng lǐ.} — \zh{在…里} = « dans, à l'intérieur de ».
\item \zh{我们应该保持环境干净。} \emph{Wǒmen yīnggāi bǎochí huánjìng gānjìng.} — \zh{保持} (bǎochí) = maintenir, garder.
\end{enumerate}
\end{corrige}

\begin{corrige}[Exercice 6 — Traduis en français (version)]
\begin{enumerate}
\item « Le paludisme est une maladie très dangereuse. » — \zh{一种} : classificateur \zh{种} (zhǒng) pour les sortes, les types.
\item « Les personnes qui ont attrapé le paludisme ont de la fièvre et mal à la tête. » — \zh{得了…的人} = « ceux qui ont contracté... » ; \zh{、} énumère les symptômes.
\item « Nous devons dormir sous la moustiquaire. »
\item « Prévenir le paludisme est plus important que le soigner. » (ou : « La prévention du paludisme est plus importante que son traitement. ») — structure comparative \cle{A + 比 + B + 更 + adjectif}.
\end{enumerate}
\end{corrige}

\begin{corrige}[Exercice 7 — Texte à trous]
\zh{因为雨季蚊子很多，所以很多人得疟疾。如果你发烧，就要马上去医院。我们应该睡在蚊帐里。预防比治疗更重要。得了疟疾的人要多喝水、多休息。}\\[0.3em]
\emph{Yīnwèi yǔjì wénzi hěn duō, suǒyǐ hěn duō rén dé nüèjí. Rúguǒ nǐ fāshāo, jiù yào mǎshàng qù yīyuàn. Wǒmen yīnggāi shuì zài wénzhàng lǐ. Yùfáng bǐ zhìliáo gèng zhòngyào. Déle nüèjí de rén yào duō hē shuǐ、duō xiūxi.}\\[0.3em]
« Parce qu'à la saison des pluies les moustiques sont nombreux, beaucoup de gens attrapent le paludisme. Si tu as de la fièvre, il faut aller tout de suite à l'hôpital. Nous devons dormir sous la moustiquaire. La prévention est plus importante que le traitement. Les malades du paludisme doivent boire beaucoup d'eau et bien se reposer. »
\end{corrige}

\begin{corrige}[Exercice 8 — Le bon caractère]
\begin{enumerate}
\item \zh{疟疾是一种危险的病。} \emph{Nüèjí shì yī zhǒng wēixiǎn de bìng.} « Le paludisme est une maladie dangereuse. » — \zh{疾} (jí) : radical \zh{疒} ; \zh{病} ne se combine pas avec \zh{疟}.
\item \zh{得了疟疾要马上治疗。} \emph{Déle nüèjí yào mǎshàng zhìliáo.} « Quand on a le paludisme, il faut se soigner tout de suite. » — \zh{治} (zhì, clé de l'eau \zh{氵}) = soigner ; \zh{活} (huó) = vivre.
\item \zh{我们睡在蚊帐里。} \emph{Wǒmen shuì zài wénzhàng lǐ.} « Nous dormons sous la moustiquaire. » — \zh{帐} (zhàng, clé du tissu \zh{巾}) = moustiquaire, tente ; \zh{账} (zhàng, clé de la coquille-monnaie \zh{贝}) = compte, comptabilité. Homophones parfaits, sens très différents !
\item \zh{在喀麦隆，生活很愉快。} \emph{Zài Kāmàilóng, shēnghuó hěn yúkuài.} « Au Cameroun, la vie est agréable. » — \zh{生活} (shēnghuó) = la vie (quotidienne).
\item \zh{生病的人要去看医生。} \emph{Shēngbìng de rén yào qù kàn yīshēng.} « Les malades doivent aller voir le médecin. » — on dit \zh{生病} (tomber malade), jamais \zh{生疾}.
\end{enumerate}
\end{corrige}

\begin{corrige}[Exercice 9 — Remets les phrases en ordre]
Ordre correct : \textbf{D -- B -- A -- E -- C}.
\begin{itemize}
\item \textbf{D} ouvre le texte : présentation générale du sujet (\zh{疟疾是…一种病}).
\item \textbf{B} donne la cause avec \zh{因为…所以…}.
\item \textbf{A} indique la conduite à tenir en cas de symptômes (\zh{如果…就…}).
\item \textbf{E} passe à la prévention avec \zh{为了} (but).
\item \textbf{C} conclut avec \zh{总之} (« en résumé »), connecteur typique de conclusion.
\end{itemize}
Texte reconstitué : \zh{疟疾是喀麦隆很常见的一种病。因为在雨季蚊子很多，所以很多人得疟疾。如果发烧、头疼，就要马上去医院。为了预防疟疾，我们应该睡在蚊帐里，保持环境干净。总之，预防比治疗更重要。}
\end{corrige}

\begin{corrige}[Exercice 10 — Compréhension de texte (type Bac)]
\textbf{Questions de compréhension}
\begin{enumerate}
\item Parce qu'il y a beaucoup d'eau de pluie (\zh{因为雨水多}), ce qui favorise la multiplication des moustiques.
\item La fièvre (\zh{发烧}) et le mal de tête (\zh{头疼}) ; le texte ajoute parfois la fatigue (\zh{很累}).
\item Dormir sous la moustiquaire (\zh{睡在蚊帐里}) et garder l'environnement propre (\zh{保持环境干净}).
\item Il faut aller tout de suite à l'hôpital voir le médecin (\zh{马上去医院看医生}).
\end{enumerate}
\textbf{Questions de langue}
\begin{enumerate}
\item \zh{预防比治疗更重要。} \emph{Yùfáng bǐ zhìliáo gèng zhòngyào.} — structure \cle{A + 比 + B + 更 + adjectif}.
\item \zh{为了} (wèile) introduit le but, l'objectif : « pour, afin de » + verbe. Ici : « pour prévenir le paludisme ».
\item C'est l'expression \zh{越来越} (yuèláiyuè) + adjectif : « de plus en plus ». \zh{蚊子也越来越多} = « les moustiques aussi sont de plus en plus nombreux ».
\end{enumerate}
\end{corrige}

\begin{corrige}[Exercice 11 — Thème et version (type Bac)]
\textbf{A. Thème.}
\begin{enumerate}
\item \zh{因为雨季蚊子很多，所以很多人得疟疾。} \emph{Yīnwèi yǔjì wénzi hěn duō, suǒyǐ hěn duō rén dé nüèjí.} — En chinois, on peut utiliser \zh{因为…所以…} ensemble dans la même phrase (contrairement au français qui évite « parce que... donc »).
\item \zh{如果你发烧，就要去医院。} \emph{Rúguǒ nǐ fāshāo, jiù yào qù yīyuàn.}
\end{enumerate}
\textbf{B. Version.} « Le paludisme est une maladie très dangereuse. Les personnes qui ont attrapé le paludisme ont de la fièvre et mal à la tête. Nous devons dormir sous la moustiquaire et garder l'environnement propre. Prévenir le paludisme est plus important que le soigner. »
\end{corrige}

\begin{corrige}[Exercice 12 — Expression écrite et dictée (type Bac)]
\textbf{A. Dictée.} Vérifier pour chaque mot :
\begin{itemize}
\item \zh{疟} (nüè) : radical \zh{疒} (maladie) + \zh{虐} (phonétique).
\item \zh{疾} (jí) : radical \zh{疒} + composant \zh{矢} (flèche/serment).
\item \zh{蚊} (wén) : clé de l'insecte \zh{虫} + \zh{文} (phonétique).
\item \zh{帐} (zhàng) : clé du tissu \zh{巾} (pas \zh{账} qui a la clé \zh{贝} coquille-monnaie).
\item \zh{烧} (shāo) : clé du feu \zh{火} + \zh{少}.
\item \zh{预} (yù) : radical \zh{页} (tête/page) + \zh{予} ; \zh{防} (fáng) : clé de la colline \zh{阝} (oreille droite) + \zh{方}.
\item \zh{治} (zhì) : clé de l'eau \zh{氵} (trois traits d'eau) + \zh{台} (pas \zh{活} qui a \zh{氵} + \zh{舌}).
\item \zh{医} (yī) : radical \zh{匚} (boîte) + \zh{矢} ; \zh{院} (yuàn) : clé \zh{阝} + \zh{完}.
\item \zh{毒} (dú) : radical \zh{毋} (mère/non) + \zh{虫} (insecte/venin).
\end{itemize}
Ordre des traits : de haut en bas, de gauche à droite, l'horizontal avant le vertical, l'extérieur avant l'intérieur pour les caractères à enclos (ex. \zh{医}, \zh{院}).

\textbf{B. Rédaction : modèle de texte attendu.}
\begin{textebox}[Modèle de rédaction]
\zh{疟疾是喀麦隆很常见的一种病，也是一种很危险的病。因为雨季蚊子很多，所以很多人容易得疟疾。为了预防疟疾，我们应该每天睡在蚊帐里，还应该保持环境干净。如果发烧、头疼，就要马上去医院看医生。总之，预防比治疗更重要！}\\[0.3em]
\emph{Nüèjí shì Kāmàilóng hěn chángjiàn de yī zhǒng bìng, yě shì yī zhǒng hěn wēixiǎn de bìng. Yīnwèi yǔjì wénzi hěn duō, suǒyǐ hěn duō rén róngyì dé nüèjí. Wèile yùfáng nüèjí, wǒmen yīnggāi měitiān shuì zài wénzhàng lǐ, hái yīnggāi bǎochí huánjìng gānjìng. Rúguǒ fāshāo、tóuténg, jiù yào mǎshàng qù yīyuàn kàn yīshēng. Zǒngzhī, yùfáng bǐ zhìliáo gèng zhòngyào!}\\[0.3em]
« Le paludisme est une maladie très courante au Cameroun, et aussi une maladie très dangereuse. Parce qu'à la saison des pluies les moustiques sont nombreux, beaucoup de gens attrapent facilement le paludisme. Pour prévenir le paludisme, nous devons dormir chaque jour sous la moustiquaire et aussi garder l'environnement propre. Si l'on a de la fièvre et mal à la tête, il faut aller tout de suite à l'hôpital voir le médecin. En résumé, la prévention est plus importante que le traitement ! »
\end{textebox}
Ce modèle utilise 4 connecteurs (\zh{因为…所以…、为了、还、如果…就…、总之}) et 9 mots du glossaire. Accepter toute production correcte respectant la grille : caractères exacts, structures correctes, vocabulaire du thème, cohérence logique (introduction -- causes -- moyens de prévention -- conclusion).
\end{corrige}

\newpage

\coursec{Fiche bilan}

\begin{popfigure}
\begin{tikzpicture}[mindmap, grow cyclic, every node/.style={concept, font=\small},
root concept/.append style={concept color=popblue, font=\bfseries, minimum size=3cm},
level 1/.append style={level distance=3.4cm, sibling angle=60},
level 1 concept/.append style={font=\footnotesize, minimum size=2.2cm}]
\node[root concept, align=center]{Leçon 14\\Expression écrite\\le paludisme}
child[concept color=poporange]{node[align=center]{Lexique\\疟疾 发烧\\蚊子 医院}}
child[concept color=popgreen]{node[align=center]{Caractères\\radicaux 疒 虫\\ordre des traits}}
child[concept color=poppurple]{node[align=center]{Structure\\du texte\\3 parties}}
child[concept color=popgold]{node[align=center]{Connecteurs\\因为 所以\\首先 然后}}
child[concept color=poppink]{node[align=center]{Traduction\\thème et version\\phrases simples}}
child[concept color=popblueL]{node[align=center]{Production\\écrite guidée\\modèle à suivre}};
\end{tikzpicture}
\legende{Carte mentale de la leçon : les six compétences à maîtriser pour écrire un texte simple sur le paludisme.}
\end{popfigure}

\begin{bilan}
\textbf{1. Lexique clé du paludisme}
\begin{center}
\begin{tabularx}{\linewidth}{|X|X|X|X|}\hline
\rowcolor{popblueL} Caractères & Pinyin & Nature & Traduction \\ \hline
疟疾 & nüèji & n. & le paludisme \\ \hline
发烧 & fāshāo & v. & avoir de la fièvre \\ \hline
头疼 & tóuténg & v. & avoir mal à la tête \\ \hline
蚊子 & wénzi & n. & le moustique \\ \hline
蚊帐 & wénzhàng & n. & la moustiquaire \\ \hline
生病 & shēngbìng & v. & tomber malade \\ \hline
医院 & yīyuàn & n. & l'hôpital \\ \hline
医生 & yīshēng & n. & le médecin \\ \hline
药 & yào & n. & le médicament \\ \hline
预防 & yùfáng & v. & prévenir \\ \hline
健康 & jiànkāng & n./adj. & la santé ; en bonne santé \\ \hline
睡觉 & shuìjiào & v. & dormir \\ \hline
\end{tabularx}
\end{center}

\textbf{2. Règles de grammaire à connaître par cœur}
\begin{center}
\begin{tabularx}{\linewidth}{|X|X|X|}\hline
\rowcolor{popblueL} Structure & Exemple & Traduction \\ \hline
Sujet + 因为 + cause, 所以 + résultat & \zh{因为蚊子多，所以很多人生病。} Yīnwèi wénzi duō, suǒyǐ hěn duō rén shēngbìng. & Parce qu'il y a beaucoup de moustiques, beaucoup de gens tombent malades. \\ \hline
Sujet + 应该 + verbe (conseil) & \zh{我们应该用蚊帐。} Wǒmen yīnggāi yòng wénzhàng. & Nous devrions utiliser une moustiquaire. \\ \hline
Sujet + 要 + verbe (il faut) & \zh{生病要去看医生。} Shēngbìng yào qù kàn yīshēng. & Quand on est malade, il faut aller voir le médecin. \\ \hline
Connecteurs : 首先、然后、最后 & \zh{首先预防，然后治疗，最后保持健康。} Shǒuxiān yùfáng, ránhòu zhìliáo, zuìhòu bǎochí jiànkāng. & D'abord prévenir, ensuite soigner, enfin rester en bonne santé. \\ \hline
\end{tabularx}
\end{center}

\textbf{3. Écriture des caractères}
\begin{itemize}
\item \zh{疟} : radical de la maladie \zh{疒} (à gauche et en haut), puis la partie phonétique à l'intérieur.
\item \zh{疾} : radical \zh{疒} + \zh{矢} ; on écrit d'abord l'extérieur, puis l'intérieur.
\item \zh{蚊} : radical de l'insecte \zh{虫} à gauche, \zh{文} à droite ; gauche avant droite.
\item \zh{帐} : radical \zh{巾} (tissu) à gauche ; attention à ne pas confondre avec \zh{账} (compte).
\item Règles d'or : haut avant bas, gauche avant droite, extérieur avant intérieur, horizontal avant vertical.
\end{itemize}

\textbf{4. Méthodes express}
\begin{itemize}
\item \textbf{Rédiger} : 3 parties — introduction (le problème), développement (causes, symptômes, prévention), conclusion (conseil). Utiliser 首先、然后、最后.
\item \textbf{Thème (français → chinois)} : repérer sujet + verbe + complément ; traduire mot à mot les structures vues ; vérifier les tons et l'ordre des mots.
\item \textbf{Version (chinois → français)} : identifier d'abord le sujet, puis le verbe ; traduire les connecteurs (因为 = parce que, 所以 = donc).
\item \textbf{Vérifier} : relire en comptant les caractères, contrôler les radicaux 疒 et 虫, et la place de 应该 / 要 avant le verbe.
\end{itemize}
\end{bilan}

\begin{aretenir}[Checklist avant le Bac]
\begin{itemize}
\item[$\square$] Je connais par cœur le lexique du paludisme : \zh{疟疾、发烧、蚊子、蚊帐、医院、预防、健康}.
\item[$\square$] J'écris correctement \zh{疟} et \zh{疾} avec le radical de la maladie \zh{疒}.
\item[$\square$] Je distingue \zh{蚊} (moustique) et \zh{帐} (moustiquaire) sans confondre les radicaux.
\item[$\square$] Je sais construire une phrase de cause avec 因为 … 所以 …
\item[$\square$] Je place correctement 应该 et 要 avant le verbe principal.
\item[$\square$] J'utilise les connecteurs 首先、然后、最后 pour organiser mon texte.
\item[$\square$] Mon texte a trois parties : introduction, développement, conclusion.
\item[$\square$] Je sais traduire une phrase simple du français vers le chinois (thème).
\item[$\square$] Je sais traduire une phrase simple du chinois vers le français (version).
\item[$\square$] Je mets le pinyin avec les bons tons sur les voyelles.
\item[$\square$] Je relis ma production pour vérifier l'ordre des mots : sujet + verbe + complément.
\item[$\square$] Je respecte la ponctuation chinoise pleine chasse ：，。？
\end{itemize}
\end{aretenir}

\findecours

\end{document}
